% concatenated from main.tex
\documentclass[12pt,a4paper]{amsart}
\usepackage{amsfonts}
\usepackage{mathrsfs}
\usepackage{amsthm,amsxtra}
\usepackage{amssymb}
\usepackage{amsmath}
\usepackage{amscd}
\usepackage{hyperref}
\hypersetup{
colorlinks=true,
linkcolor=blue,
citecolor=gray,
%filecolor=magenta,      
%urlcolor=cyan,
}
\usepackage[utf8]{inputenc}
\usepackage{t1enc}
\usepackage[mathscr]{eucal}
\usepackage{indentfirst}
\usepackage{graphicx}
\graphicspath{ {images/} }
\usepackage{subfig}
\usepackage{enumerate}
\usepackage{graphics}
\usepackage{pict2e}
\usepackage{epic}
\numberwithin{equation}{section}
\usepackage[margin=2.9cm]{geometry}
\usepackage{epstopdf} 
\usepackage{longtable}
\usepackage{enumerate}
\usepackage[inline]{enumitem}
\usepackage{tikz-cd}
\usepackage{tikz}
\usepackage{color,soul}


%%%%%%%%%%%%%%%%%%%%%%%%%%%%%%%%%%%%%%%%%%
% Todonotes Setting %%%%%%%%%%%%%%%%%%%%%%%%%%%%%%%
%%%%%%%%%%%%%%%%%%%%%%%%%%%%%%%%%%%%%%%%%%

% \usepackage{todonotes}

% \newcommand{\argue}[1]{\todo[color=red!40, tickmarkheight=0.1cm]{#1}}
% \newcommand{\plan}[1]{\todo[color=blue!40, tickmarkheight=0.1cm]{#1}}
% \newcommand{\etcnote}[1]{\todo[color=yellow!40, tickmarkheight=0.1cm]{#1}}

%%%%%%%%%%%%%%%%%%%%%%%%%%%%%%%%%%%%%%%%%%



%%%%%%%%%%%%%%%%%%%%%%%%%%%%%%%%%%%%%%%%%%
% Declaring operators %%%%%%%%%%%%%%%%%%%%%%%%%%%%%%%
%%%%%%%%%%%%%%%%%%%%%%%%%%%%%%%%%%%%%%%%%%

\DeclareMathOperator{\supp}{supp}
\DeclareMathOperator{\Lk}{Lk}
\DeclareMathOperator{\interior}{interior}
\DeclareMathOperator{\St}{St}
\DeclareMathOperator{\Mod}{Mod}
\DeclareMathOperator{\MCG}{MCG}
\DeclareMathOperator{\ML}{\mathcal{ML}}
\DeclareMathOperator{\PML}{\mathbb{P}\mathcal{ML}}
\DeclareMathOperator{\Out}{Out}
\DeclareMathOperator{\Aut}{Aut}
\DeclareMathOperator{\Inn}{Inn}
\DeclareMathOperator{\Stab}{Stab}
\DeclareMathOperator{\Cay}{Cay}
\DeclareMathOperator{\Fix}{Fix}

\newcommand{\FF}{\mathbb{F}}
\newcommand{\MM}{\mathbb{M}}
\newcommand{\ZZ}{\mathbb{Z}}

\newcommand{\cev}[1]{\reflectbox{\ensuremath{\vec{\reflectbox{\ensuremath{#1}}}}}}


%%%%%%%%%%%%%%%%%%%%%%%%%%%%%%%%%%%%%%%%%%



\def\numset#1{{\\mathbb #1}}



\theoremstyle{plain}
\newtheorem{theorem}{Theorem}[section]
\newtheorem{proposition}[theorem]{Proposition}
\newtheorem{lemma}[theorem]{Lemma}
\newtheorem{corollary}[theorem]{Corollary}

\theoremstyle{definition}
\newtheorem{definition}[theorem]{Definition}
\newtheorem{notation}[theorem]{Notation}
\newtheorem{conjecture}[theorem]{Conjecture}
\newtheorem{question}[theorem]{Question}

\theoremstyle{remark}
\newtheorem{remark}[theorem]{Remark}
\newtheorem*{remark*}{Remark}
\newtheorem{example}[theorem]{Example}
\newtheorem{exercise}[theorem]{Exercise}



\begin{document}

\title[Dynamics of Dehn twists in $\mathrm{Out}(\FF)$]{Dynamics of Dehn Twists in the Outer Automorphism Group of a Free Group}
	
\author{Donggyun Seo}
\address{Institute of Mathematical Sciences \\ Chungnam National University \\ Daejeon, Korea}
\email{seodonggyun@cnu.ac.kr \\ seodonggyun7@gmail.com}
\urladdr{https://www.seodonggyun.com}

\keywords{outer automorphism group of free group, doubled handlebody, right-angled Artin group, Dehn twist, intersection number, bigon--bihedron criterion, ping-pong lemma} 
\subjclass[2020]{primary: 20F28, 20F65, 57M07 secondary: 20E08, 20F36, 57K30, 57M60}


\begin{abstract}
    We study Dehn twists in the outer automorphism group of a finitely generated non-abelian free group.  
    Our main result states that, under certain compatibility conditions, sufficiently large powers of finitely many Dehn twists generate a right-angled Artin group.
    The proof proceeds by analyzing the geometry of spheres, tori, and simple closed curves in a doubled handlebody.
    Along the way, we establish the bigon--bihedron criterion and an equivalent condition for commuting Dehn twists.
    Furthermore, we construct a compact topological space on which multi-twists act parabolically, fixing simplices of weighted multicores, and fully irreducible elements act with north–south dynamics.
\end{abstract}

\maketitle




\section{Introduction}


Let $\mathbb{F} = \mathbb{F}_g$ be a free group of rank $g$.
An \emph{$\mathbb{F}$-tree} is a simplicial tree on which $\mathbb{F}$ acts simplicially without edge inversions.
An $\mathbb{F}$-tree is called a \emph{$\mathbb{Z}$-splitting} if every edge stabilizer is infinite cyclic.
It is said to be \emph{one-edge} if the induced action of $\mathbb{F}$ on the set of edges is transitive.
In this case, the conjugacy class of an edge stabilizer of a one-edge $\mathbb{Z}$-splitting $T$ is uniquely determined, and we call it the \emph{core} of $T$.

For a conjugacy class $\gamma$ of an infinite cyclic subgroup of $\mathbb{F}$, the \emph{intersection number} of $\gamma$ and $T$, denoted by $i(\gamma, T)$, is defined as the minimal translation length of a nontrivial element of a subgroup representing $\gamma$ on $T$:
\[
    i(\gamma, T)
    = \min\{\, d_T(x, wx) \mid 1 \neq w \in c \in \gamma,\; x \in T \,\},
\]
where $d_T$ denotes the combinatorial metric on $T$.
For example, if $\gamma$ is the core of $T$, then $i(\gamma, T) = 0$.

% Let $\mathbb{F} = \mathbb{F}_g$ be a free group of rank $g$.
% A simplicial tree where $\FF$ acts simplicially without edge inversion is called an \emph{$\FF$-tree}.
% An $\FF$-tree is called a \emph{$\mathbb{Z}$-splitting} if every edge stabilizer is an infinite cyclic group.
% An $\FF$-tree is said to be \emph{one-edge} if the induced action of $\FF$ on the set of edges is transitive.
% Then the conjugacy class of an edge stabilizer of a one-edge $\mathbb{Z}$-splitting $T$ is uniquely determined, which will be called the \emph{core} of $T$.

% For a conjugacy class $\gamma$ of an infinite cyclic subgroup of $\FF$, the \emph{intersection number} of $\gamma$ and $T$, denoted by $i(\gamma, T)$, is the minimal translation length of a nontrivial element of a subgroup representing $\gamma$ on $T$, that is, \[ i(\gamma, T) = \min\{ d_T(x,wx) \mid 1 \neq w \in c \in \gamma ~\text{and}~ x \in T \} \] where $d_T$ is the combinatorial metric of $T$.
% For example, if $\gamma$ is the core of $T$, then we have $i(\gamma, T) = 0$.

% Let $\mathbb{F} = \mathbb{F}_g$ be a free group of rank $g$.
The automorphism group $\Aut(\mathbb{F})$ contains the inner automorphism group
$\Inn(\mathbb{F})$ as a normal subgroup, and the quotient
\[
    \Out(\mathbb{F}) := \Aut(\mathbb{F}) / \Inn(\mathbb{F})
\]
is called the \emph{outer automorphism group} of $\mathbb{F}$.
Given a one-edge $\mathbb{Z}$-splitting $T$ of $\mathbb{F}$, one can associate a \emph{Dehn twist} (or \emph{generalized transvection}) $\delta_{\vec{T}} \in \Out(\mathbb{F})$, defined as an outer automorphism induced by twisting along $T$.

Our main theorem concerns the existence of right-angled Artin subgroups generated by powers of such Dehn twists.
For a finite simplicial graph $\Gamma = (V, E)$, the \emph{right-angled Artin group (RAAG)} on $\Gamma$ is
\[
    A(\Gamma) := \langle\, v \in V \mid [u,v] = 1 \text{ whenever } \{u,v\} \in E \,\rangle.
\]
Given a finite collection $\{T_1,\dots,T_k\}$ of one-edge $\mathbb{Z}$-splittings of $\mathbb{F}$, one defines the \emph{coincidence graph} $\Gamma$ with vertices $T_1,\dots,T_k$, where $T_i$ and $T_j$ are adjacent precisely when they admit a common refinement.
For integers $n_1,\dots,n_k$, this determines a homomorphism
\[
    A(\Gamma) \longrightarrow \Out(\mathbb{F}), \qquad T_i \longmapsto \delta_{\vec{T}_i}^{\,n_i}.
\]

Our main result asserts that, under suitable hypotheses, this map is injective and its image is a right-angled Artin subgroup of $\Out(\mathbb{F})$.
Here the collection is \emph{compatible} if the cores are pairwise distinct and indivisible, and any two splittings that fail to admit a common refinement have cores penetrating each other.

\begin{theorem}[Theorem \ref{thm:main}] \label{thm:1}
    Let $V = \{ T_1, \dots, T_k \}$ be a compatible collection of one-edge $\mathbb{Z}$-splittings with cores $\gamma_1, \dots, \gamma_k$, and let $\Gamma$ be the coincidence graph of $V$.
    For each $i$, let $\delta_i = \delta_{\vec{T}_i}$ be a Dehn twist along $T_i$.
    Then for all
    \[
        \lvert n_1 \rvert, \dots, \lvert n_k \rvert
        \;\ge\;
        \max_{i \ne j} i(\gamma_i, T_j)
        \;+\;
        14 \max_{i \ne j} i(T_i, T_j)^2
        \;+\; 8,
    \]
    the homomorphism
    \[
        \Phi \colon A(\Gamma) \longrightarrow \Out(\FF), \qquad T_i \longmapsto \delta_i^{n_i}
    \]
    is injective.
\end{theorem}

The problem of constructing free subgroups of $\Out(\mathbb{F})$ has been extensively studied.  
Bestvina, Feighn, and Handel \cite{MR1765705} proved that $\Out(\mathbb{F})$ satisfies the Tits alternative: every subgroup is either virtually solvable or contains a free subgroup of rank two.  
Kapovich and Lustig \cite{MR2652906} showed that sufficiently large powers of two hyperbolic iwip elements generate either a virtually cyclic subgroup or a nonabelian free group.  
Uyanik \cite{MR3954286} obtained a related result in the setting of hyperbolic extensions.  
Clay and Pettet \cite{MR2720227} proved that if two one-edge $\mathbb{Z}$-splittings are filling, then sufficiently large powers of the associated Dehn twists generate a free group of rank two consisting of atoroidal fully irreducible elements except for the conjugacy classes of the twists themselves.  
G\"ultepe \cite{MR3677931} provided another criterion under which powers of two Dehn twists generate a free subgroup of rank two.

A complementary development is due to Taylor \cite{MR3343354}, who proved that every right-angled Artin group embeds quasi-isometrically into some $\Out(\mathbb{F})$.  
He constructed a suitable collection of fully irreducible outer automorphisms supported on free factors, whose sufficiently large powers generate the desired RAAG.  
Since these automorphisms are fully irreducible, Dehn twists do not appear in his construction.  
Our work may therefore be viewed as a counterpart to Taylor's theorem: we construct right-angled Artin subgroups of $\Out(\mathbb{F})$ generated instead by powers of Dehn twists.

Theorem~\ref{thm:1} gives a broad and explicit method for constructing RAAG subgroups inside $\Out(\mathbb{F})$.
However, the compatibility condition does not encompass all right-angled Artin phenomena arising from Dehn twists.
For instance, $\Out(\mathbb{F}_g)$ contains a direct product of $2g-4$ free groups generated by Dehn twists (Horbez and Wade \cite{MR4069231}, and Bridson and Wade \cite{MR4792836}) and a free-abelian subgroup of rank $2g-3$ generated by Dehn twists (Feighn and Handel \cite{MR2496054}), yet the associated one-edge $\mathbb{Z}$--splittings in these examples are not compatible.
Thus, Theorem~\ref{thm:1} captures a large and well-structured class of RAAGs generated by powers of Dehn twists, but it does not apply to every such instance found in $\Out(\mathbb{F}_g)$.

We approach the proof of Theorem~\ref{thm:1} by working in the mapping class group of the doubled handlebody
\[
    \MM = \MM_g = \#_g(\mathbb{S}^1 \times \mathbb{S}^2).
\]
The interaction between $\Out(\FF)$ and $\MCG(\MM)$ was first studied by Whitehead \cite{MR1575455, MR1503309} in his solution to the automorphic equivalence problem for $\FF$.
In particular, he proved that the natural map
\[
    \MCG(\MM) \longrightarrow \Out(\FF)
\]
is surjective.
Laudenbach \cite{MR356056} later established the short exact sequence
\[
    1 \longrightarrow (\mathbb{Z}/2\mathbb{Z})^{g} \longrightarrow \MCG(\MM) \longrightarrow \Out(\FF) \longrightarrow 1.
\]
Brendle, Broaddus, and Putman \cite{MR4557874} subsequently proved that this sequence splits; in particular, there exists a section
\[
    \Out(\FF) \longrightarrow \MCG(\MM). 
\]

In this setting, a Dehn twist may be interpreted as the isotopy class of a homeomorphism obtained by twisting $\MM$ along an essential torus.
G{\"u}ltepe \cite{MR3677931} likewise studied Dehn twists realized as twists along essential tori.

% We approach the proof of Theorem~\ref{thm:1} by working in the mapping class group of the doubled handlebody \(\MM = \MM_g = \#_g(\mathbb{S}^1 \times \mathbb{S}^2)\).
% The interaction between $\Out(\FF)$ and $\MCG(\MM)$ has been initiatively studied by Whitehead \cite{MR1575455, MR1503309} to solve the automorphic equivalent problem of $\FF$.
% He proved that the induced map $\MCG(\MM) \to \Out(\FF)$ is surjective.
% Laudenbach \cite{MR356056} established the following short exact sequence:
% \[
%     1 \longrightarrow (\mathbb{Z}/2\mathbb{Z})^{g} \longrightarrow \MCG(\MM) \longrightarrow \Out(\mathbb{F}) \longrightarrow 1.
% \]
% Brendle, Broaddus, and Putman \cite{MR4557874} later proved that this sequence splits; in particular, there exists a section \(\Out(\mathbb{F}) \to \MCG(\MM)\).
% In this setting, a Dehn twist may be interpreted as an isotopy class of a homeomorphism twisting \(\MM\) along an essential torus.
% G{\"u}ltepe \cite{MR3677931} likewise studied Dehn twists realized as twists along essential tori.

We begin by analyzing pairs of commuting Dehn twists.
Given homotopy classes of essential tori \(T\) and \(T'\), their \emph{geometric intersection number} \(i(T, T')\) is defined as the minimum number of intersection circles among all representatives \(t \in T\) and \(t' \in T'\).
If \(i(T, T') = 0\), that is, if \(T\) and \(T'\) admit disjoint representatives, then the corresponding Dehn twists commute.
Our next result establishes the converse: two Dehn twists commute if and only if their supports may be chosen disjointly.
In particular, no accidental commutation occurs.


% We approach to the proof of Theorem \ref{thm:1} by focusing on the mapping class group of the doubled handlebody \(\MM = \MM_g = \#_g(\mathbb{S}^1 \times \mathbb{S}^2)\).
% Laudenbach \cite{MR356056} established the following short exact sequence:
% \[
%     1 \longrightarrow (\mathbb{Z} / 2\mathbb{Z})^g \longrightarrow \MCG(\MM) \longrightarrow \Out(\mathbb{F}) \longrightarrow 1.
% \]
% Brendle, Broaddus and Putman \cite{MR4557874} proved that this sequence splits; in particular, there exists a section \(\Out(\mathbb{F}) \to \MCG(\MM)\).  
% In this setting, a Dehn twist is regarded as an isotopy class of a homeomorphism that twists \(\MM\) along an essential torus.
% G{\"u}ltepe \cite{MR3677931} also treated Dehn twists as isotopy classes of homeomorphisms twisting along essential tori.

% We begin by focusing on pairs of commuting Dehn twists.  
% For homotopy classes of essential tori \(T\) and \(T'\), the \emph{(geometric) intersection number} \(i(T, T')\) is defined as the minimum number of circles in \(t \cap t'\), taken over all representatives \(t \in T\) and \(t' \in T'\).  
% If \(i(T, T') = 0\), that is, if \(T\) and \(T'\) admit disjoint representatives, then the corresponding Dehn twists commute.
% Our next result establishes the converse: two Dehn twists commute if and only if their supports can be chosen disjointly---in particular, accidental commutation does not occur.

\begin{theorem}[Theorem \ref{thm:commuting}]
    Let $T$ and $T'$ be homotopy classes of essential tori.
    Then the following statement holds:
    \begin{center}
        $[\delta_{\vec{T}}, \delta_{\vec{T}'}] = 1$ if and only if $i(T, T') = 0$.
    \end{center}
\end{theorem}


From an algebraic perspective, two Dehn twists commute precisely when their associated one-edge $\mathbb{Z}$-splittings admit a common refinement.
On the geometric side, recall that two simple closed curves on a surface intersect minimally if and only if they form no bigon---the classical \emph{bigon criterion}, which underlies much of surface topology.
In an analogous spirit, our work is grounded in the following criterion governing intersections of essential tori in doubled handlebodies.

% In the algebraic view, one can say that two Dehn twists commute if and only if their corresponding one-edge $\mathbb{Z}$-splittings have a common refinement.
% Note that two simple closed curves in a surface minimally intersect if and only if they have no bigon, which is generally called the bigon criterion.
% The bigon criterion for simple closed curves of surfaces serves as a basis for the surface theory.
% Our results are also based on the following criterion.

\begin{proposition}[Bigon--bihedron criterion, Proposition \ref{prop:bigon-bihedron}]
    Two essential tori intersect minimally if and only if they form neither a bigon nor a bihedron.
\end{proposition}

One can formulate analogous criteria in the cases of two essential spheres or of an essential sphere together with an essential torus.  
In addition, one obtains a \emph{bigon criterion} in this setting: an essential simple closed curve and an essential sphere (or an essential torus) intersect minimally if and only if they do not form a bigon.  
We note that Hatcher \cite{MR1314940} introduced a normal form for sphere systems, and this normal form guarantees that any two sphere systems placed in it are in minimal position.

Meanwhile, Theorem~\ref{thm:1} is proved by a ping--pong argument on the set of homotopy classes of essential simple closed curves.  
We employ a variation of the ping--pong method due to Koberda \cite{MR3000498}.  
Careful control of intersection numbers within this framework yields the proof of Theorem~\ref{thm:1}.

% In addition, we construct a compact topological space \(\overline{(p \circ i_*)(\mathcal{C})}\) on which \(\Out(\mathbb{F})\) acts continuously.
% This space is precisely the closure of the set of points represented by homotopy classes of essential simple closed curves that represent $1$-free factors of $\FF$.
% The definition of \(\overline{(p \circ i_*)(\mathcal{C})}\) is closely analogous to that of the projective space \(\mathcal{PML}(S)\) of measured laminations on a surface \(S\).
% In \(\mathcal{PML}(S)\), the action of a Dehn twist is parabolic, and a Dehn twist in \(\Out(\mathbb{F})\) exhibits the same dynamical behavior on \(\overline{(p \circ i_*)(\mathcal{C})}\), as described below.

% \begin{theorem}[Theorem \ref{thm:converge}]
%     Let $T$ be a homotopy class of an essential torus with core $\gamma$.
%     Let $\alpha$ be a homotopy class of an essential simple closed curve that represents a free factor.
%     Then, we have $i(\alpha, T) > 0$ if and only if $\delta_{\vec{T}}^n(\alpha)$ converges to $\gamma$ as $n \to \infty$.
% \end{theorem}

% The dynamics of a Dehn twist described in Theorem~\ref{thm:converge} is parabolic, fixing the core of its torus and pushing primitive curves toward it.
% This is the behavior of a Dehn twist on the projective space of measured laminations of a surface.
% On a surface, the complementary dynamical type is realized by pseudo-Anosov mapping classes, which act on $\mathcal{PML}(S)$ with north–south dynamics.
% The analog of a pseudo-Anosov in $\Out(\FF)$ is a fully irreducible (iwip) outer automorphism, which, by Martin \cite{MR2693216} and Uyanik \cite{MR3273532, MR3370027}, acts with north–south dynamics on the unique minimal set of projective geodesic currents.
% Our final result transports this behavior to the space $\overline{(p \circ i_*)(\mathcal{C})}$, showing that fully irreducible elements act on it with precisely this complementary, loxodromic behavior.
% Together with Theorem~\ref{thm:converge}, this exhibits on a single compact $\Out(\FF)$-space the same parabolic–loxodromic dichotomy that governs the action of a mapping class group on the space of projective measured laminations of a surface.

% \begin{theorem}[Theorem~\ref{thm:loxodromic}]\label{thm:loxodromic-intro}
%     Let $\phi \in \Out(\FF)$ be a fully irreducible outer automorphism.
%     Then $\phi$ acts on the compact space $\overline{(p \circ i_*)(\mathcal{C})}$ with north–south dynamics.
% \end{theorem}

Our next results describe the dynamics of Dehn twists on the compact space $\overline{(p \circ i_*)(\mathcal{C})}$.
Recall that on a surface $S$, a multitwist along a multicurve $c_1 \cup \dots \cup c_r$ acts on the projective space $\mathcal{PML}(S)$ of measured laminations parabolically: it fixes the simplex of projective measured laminations supported on $c_1 \cup \dots \cup c_r$ pointwise, and it attracts every lamination crossing the multicurve into that simplex.
We prove that multi-twists in $\Out(\FF)$ exhibit precisely this behavior on $\overline{(p \circ i_*)(\mathcal{C})}$.

By Theorem~\ref{thm:commuting}, the Dehn twists along pairwise disjoint homotopy classes $T_1, \dots, T_r$ of essential tori commute, so every product $\Delta = \delta_{\vec{T}_1}^{m_1} \dots \delta_{\vec{T}_r}^{m_r}$, which we call a \emph{multi-twist}, is well-defined independently of the order of the factors.

\begin{theorem}[Corollary~\ref{cor:multitwist-parabolic}] \label{thm:3}
    Let $T_1, \dots, T_r$ be pairwise disjoint homotopy classes of essential tori with cores $\gamma_1, \dots, \gamma_r$, and let $\Delta = \delta_{\vec{T}_1}^{m_1} \dots \delta_{\vec{T}_r}^{m_r}$ be a multi-twist.
    For every homotopy class $\alpha$ of a primitive curve, exactly one of the following holds:
    \begin{enumerate}
        \item $i(\alpha, T_v) = 0$ for every $v$ with $m_v \neq 0$, and $\Delta^n(\alpha) = \alpha$ for all $n \in \ZZ$; or
        \item as $\lvert n \rvert \to \infty$, the sequence $\Delta^n(\alpha)$ converges in $\overline{(p \circ i_*)(\mathcal{C})}$ to the projective class of the weighted sum
              \[
                  \sum_{v=1}^{r} \lvert m_v \rvert \, i(\alpha, T_v)\, i_*(\gamma_v).
              \]
    \end{enumerate}
\end{theorem}

For $r = 1$, Theorem~\ref{thm:3} describes the parabolic dynamics of a single Dehn twist: the twist fixes the core of its torus and pushes every primitive curve crossing the torus toward it.
For $r > 1$, however, the limit in Theorem~\ref{thm:3} is in general not a curve at all: it is a genuine projective combination of the cores, whose weights record the intersection pattern of $\alpha$ with the system $T_1, \dots, T_r$.
This raises the question of which combinations actually arise, and the answer turns out to be: all of them.

\begin{theorem}[Corollary~\ref{cor:core_simplex}] \label{thm:4}
    Let $T_1, \dots, T_r$ be pairwise disjoint homotopy classes of essential tori with cores $\gamma_1, \dots, \gamma_r$, and suppose some homotopy class $\alpha$ of a primitive curve satisfies $i(\alpha, T_v) > 0$ for all $v$.
    Then for every choice of nonnegative reals $s_1, \dots, s_r$, not all zero, the projective class
    \[
        \Bigl[\, \sum_{v=1}^{r} s_v \, i_*(\gamma_v) \,\Bigr]
    \]
    lies in $\overline{(p \circ i_*)(\mathcal{C})}$.
    In other words, $\overline{(p \circ i_*)(\mathcal{C})}$ contains the entire projectivized cone spanned by the cores $\gamma_1, \dots, \gamma_r$, and this cone is fixed pointwise by every multi-twist along $T_1, \dots, T_r$.
\end{theorem}

Theorem~\ref{thm:4} shows that $\overline{(p \circ i_*)(\mathcal{C})}$ is considerably richer than the set of curve classes it is defined from: alongside the primitive curves and the cores of essential tori, it contains continuous families of ``weighted multicores,'' organized into simplices attached to the disjoint torus systems of $\MM$.
This is the exact counterpart of the simplices of weighted multicurves in $\mathcal{PML}(S)$, and it is proved by the same mechanism: the weighted sum in Theorem~\ref{thm:4} is exhibited as a limit of primitive curves $\Delta_{\vec{N}}(\alpha)$ under multi-twists whose exponent vectors $\vec{N}$ tend to infinity in a direction prescribed by $s_1, \dots, s_r$ (Theorem~\ref{thm:multi-converge}).

On a surface, the dynamical type complementary to the parabolic one is realized by pseudo-Anosov mapping classes, which act on $\mathcal{PML}(S)$ with north--south dynamics.
The analog of a pseudo-Anosov in $\Out(\FF)$ is a fully irreducible (iwip) outer automorphism, which, by Martin \cite{MR2693216} and Uyanik \cite{MR3273532, MR3370027}, acts with north--south dynamics on the unique minimal set of projective geodesic currents.
Our final result transports this behavior to the space $\overline{(p \circ i_*)(\mathcal{C})}$, showing that fully irreducible elements act on it with precisely this complementary, loxodromic behavior.

\begin{theorem}[Theorem~\ref{thm:loxodromic}]\label{thm:loxodromic-intro}
    Let $\phi \in \Out(\FF)$ be a fully irreducible outer automorphism.
    Then $\phi$ acts on the compact space $\overline{(p \circ i_*)(\mathcal{C})}$ with north--south dynamics.
\end{theorem}

Together, Theorems~\ref{thm:3}, \ref{thm:4} and~\ref{thm:loxodromic-intro} exhibit on a single compact $\Out(\FF)$-space the same structure that governs the action of a mapping class group on the space of projective measured laminations: multi-twists act parabolically, fixing simplices of weighted multicores and attracting crossing curves into them, while fully irreducible elements act loxodromically with two fixed points.

% There is a natural relationship between \(\overline{(p \circ i_*)(\mathcal{C})}\) and the projective space \(\mathbb{P}\mathrm{Curr}(\mathbb{F})\) of geodesic currents on \(\mathbb{F}\).  
% The Kapovich--Lustig intersection form  
% \[
%     \langle \cdot,\cdot \rangle : \overline{cv} \times \mathrm{Curr}(\mathbb{F}) \to \mathbb{R}_+
% \]
% introduced in \cite{MR2197815, MR2496058} defines a continuous map  
% \[
%     i_{**} : \mathrm{Curr}(\mathbb{F}) \to \mathbb{R}_+^{\overline{cv}}
% \]
% defined by
% \[
%     i_{**}(\nu)(T) = \langle T, \nu \rangle
% \]
% for all \(T \in \overline{cv}\) and \(\nu \in \mathrm{Curr}(\mathbb{F})\).  
% Let \(\iota_{\mathcal{C}} : \mathcal{C} \hookrightarrow \mathrm{Curr}(\mathbb{F})\) be the natural injection sending each primitive curve \(\gamma \in \mathcal{C}\) to its counting current \(\eta_\gamma\).  
% With the projection \(\pi : \mathbb{R}_+^{\overline{cv}} \to \mathbb{R}_+^{\mathcal{S}}\), we obtain  
% \[
%     i_* = \pi \circ i_{**} \circ \iota_{\mathcal{C}}.
% \]
% Meanwhile, the action of \(\Out(\mathbb{F})\) on \(\mathbb{P}\mathrm{Curr}(\mathbb{F})\) is not minimal.  
% By Martin \cite{MR2693216}, there exists a unique proper minimal closed subset \(\mathbb{P}M(\mathbb{F}) \subset \mathbb{P}\mathrm{Curr}(\mathbb{F})\) containing the projective classes of counting currents of primitive (i.e., primitive) conjugacy classes of \(\mathbb{F}\).  
% This yields a continuous surjective map  
% \[
%     \Psi : \mathbb{P}M(\mathbb{F}) \to \overline{(p \circ i_*)(\mathcal{C})}
% \]
% defined by
% \[
%     [\nu] \mapsto (p \circ \pi \circ i_{**})(\nu),
% \]
% for each \([\nu] \in \mathbb{P}M(\mathbb{F})\).  
% Kapovich, Levitt, Schupp and Shpilrain \cite{MR2272138} showed that \(i_{**}\) is not injective, and the restriction \(i_{**}|_{\mathbb{P}M(\mathbb{F})}\) remains non-injective.
% Consequently, the induced map \(\Psi\) is also not injective.



\subsection{Acknowledgments}

The author thanks Sang-hyun Kim for many helpful and inspiring conversations during the preparation of this work.
During the revision process, discussions with Michael Kopreski, Thomas Hill, George Shaji, and Brian Udall
highlighted an error in a previous version of Lemma \ref{lem:dehn_disj}, leading to its improvement.
The author is also grateful to Thomas Koberda for valuable comments that strengthened the manuscript.
Thanks to Ignat Soroko, the author became aware of Whitehead's work.

This paper is dedicated to my son, Jian, who arrived as this manuscript was taking shape.









\section{Dehn twist}





In this section, we review the basic background on outer automorphism groups of free groups and doubled handlebodies required for our discussion.
Let $\FF = \FF_g$ be a free group of rank $g \geq 2$.
We write $\Aut(\FF)$ and $\Inn(\FF)$ for the automorphism group and the inner automorphism group of $\FF$, respectively.
Note $\Inn(\FF)$ is a normal subgroup of $\Aut(\FF)$.
The \emph{outer automorphism group} of $\FF$, denoted by $\Out(\FF)$, is the quotient of $\Aut(\FF)$ by $\Inn(\FF)$.





\subsection{Dehn twist automorphisms of free groups}





An $\FF$-tree $T$ is called a \emph{one-edge $\mathbb{Z}$-splitting} if the action of $\FF$ on the set of edges is transitive and a stabilizer of an edge is an infinite cyclic subgroup of $\FF$.
Let $\vec{E}T$ be the set of orientations of edges of $T$.
Then there exists an orientation reversing $*: \vec{E}T \to \vec{E}T$, that is, $*(\vec{e})$ is the reverse orientation of $\vec{e}$ for each $\vec{e} \in \vec{E}T$.
Let $o: \vec{E}T \to \FF$ be a map defined by
\begin{enumerate}
    \item $\langle o(\vec{e}) \rangle = \operatorname{Stab}(e)$,
    \item $o(*(\vec{e}))= o(\vec{e})^{-1}$ and
    \item $o(w\vec{e}) = wo(\vec{e})w^{-1}$
\end{enumerate}
for all $\vec{e} \in \vec{E}T$ and $w \in \FF$.
We call $o$ an \emph{orientation} of $T$ and write $\vec{T} := (T, o)$ for an oriented one-edge $\mathbb{Z}$-splitting.

Fix a base-vertex $v_0 \in T$.
For each vertex $v \in T$, there exists a unique geodesic path from $v_0$ to $v$, denoted by $\vec{e}_1, \dots, \vec{e}_k$.
Then we define \[\delta_{\vec{T}, v_0}(v) := o(\vec{e}_1) \dots o(\vec{e}_k)v,\] which is called a \emph{Dehn twist automorphism based at $v_0$}.
We then show the following.

\begin{lemma}
    Let $\delta_{\vec{T}, v_0}: VT \to VT$ be defined as above.
    Then $\delta_{\vec{T}, v_0}$ can be extended to a simplicial automorphism of $T$.
\end{lemma}

\begin{proof}
    Suppose $v$ and $v'$ are adjacent.
    Without loss of generality, let us assume that $v'$ is closer to $v_0$ than $v$.
    Let $\vec{e}_1, \dots, \vec{e}_k$ be the geodesic path from $v_0$ to $v$.
    Since $o(\vec{e}_k)$ stabilizes the edge joining $v'$ and $v$, we have
    \begin{align*}
        d_T(\delta_{\vec{T}, v_0}(v), \delta_{\vec{T}, v_0}(v')) &= d_T(o(\vec{e}_1)\dots o(\vec{e}_{k-1}) o(\vec{e}_{k})v, o(\vec{e}_1)\dots o(\vec{e}_{k-1})v') \\
        &= d_T(o(\vec{e}_k)v,v') = 1.
    \end{align*}
    That is, $\delta_{\vec{T}, v_0}(v)$ and $\delta_{\vec{T}, v_0}(v')$ are adjacent.
    So $\delta_{\vec{T}, v_0}$ can be extended to a simplicial automorphism of $T$.
\end{proof}

\begin{lemma}
    For each vertex $v \in T$, if $\vec{e}_1, \dots, \vec{e}_k$ is the geodesic path from $v_0$ to $v$, let us define \[\delta_{\vec{T}, v_0}^{-1}(v) := o(\vec{e}_1)^{-1} \dots o(\vec{e}_k)^{-1}(v).\]
    Then the map $\delta_{\vec{T}, v_0}^{-1}: VT \to VT$ is the inverse of $\delta_{\vec{T}, v_0}$.
\end{lemma}

\begin{proof}
    Since the sequence \[\vec{e}_1, o(\vec{e}_1)^{-1}\vec{e}_2, \dots, o(\vec{e}_1)^{-1} \dots o(\vec{e}_{k-1})^{-1}\vec{e}_k\] is the geodesic path from $v_0$ to $o(\vec{e}_1)^{-1} \dots o(\vec{e}_k)^{-1}v$, we obtain
    \begin{align*}
        \delta_{\vec{T}, v_0}\delta_{\vec{T}, v_0}^{-1}(v) &= \delta_{\vec{T}, v_0}(o(\vec{e}_1)^{-1} \dots o(\vec{e}_k)^{-1}v) \\
        & = o(\vec{e}_1) o(o(\vec{e}_1)^{-1}\vec{e}_2) \dots o(o(\vec{e}_1)^{-1} \dots o(\vec{e}_{k-1})^{-1}\vec{e}_k)( o(\vec{e}_1)^{-1} \dots o(\vec{e}_k)^{-1}v ) \\
        & = v.
    \end{align*}
    Note $\vec{e}_1, o(\vec{e}_1)\vec{e}_2, \dots, o(\vec{e}_1) \dots o(\vec{e}_{k-1})\vec{e}_k$ is the geodesic path from $v_0$ to $\delta_{\vec{T}, v_0}(v)$.
    So we have
    \begin{align*}
        \delta_{\vec{T}, v_0}^{-1} \delta_{\vec{T}, v_0}(v) &= \delta_{\vec{T}, v_0}^{-1}(o(\vec{e}_1) \dots o(\vec{e}_k)v)\\
        &= o(\vec{e}_1)^{-1}o(o(\vec{e}_1)\vec{e}_2)^{-1} \dots o(o(\vec{e}_1) \dots o(\vec{e}_{k-1})\vec{e}_k)^{-1}(o(\vec{e}_1) \dots o(\vec{e}_k)v) \\
        &= v
    \end{align*}
    Therefore, the statement holds.
\end{proof}

\begin{lemma}
    For each $w \in \FF$, it holds that $w\delta_{\vec{T}, v_0} w^{-1} = \delta_{\vec{T},wv_0}$.
\end{lemma}

\begin{proof}
    For a vertex $v \in VT$, let $\vec{e}_1, \dots, \vec{e}_k$ be the geodesic from $wv_0$ to $v$.
    Since $w^{-1}\vec{e}_1, \dots, w^{-1}\vec{e}_k$ is a geodesic from $v_0$ to $w^{-1}v$, we have
    \[
        w\delta_{\vec{T}, v_0}w^{-1}v = wo(w^{-1}\vec{e}_1) \dots o(w^{-1}\vec{e}_k)w^{-1}v = o(\vec{e}_1) \dots o(\vec{e}_k)v = \delta_{\vec{T}, wv_0}(v)
    \]
    So the statement holds.
\end{proof}

\begin{lemma}
    For a vertex $v \in T$, if $\vec{e}_1, \dots, \vec{e}_k$ is any path from $v_0$ to $v$, then we have $\delta_{\vec{T}, v_0}(v) = o(\vec{e}_1) \dots o(\vec{e}_k) v$.
\end{lemma}

\begin{proof}
    If $\vec{e}_1, \dots, \vec{e}_k$ is not geodesic, then there exists $i$ such that $\vec{e}_{i+1} = *(\vec{e}_i)$.
    So we have $o(\vec{e}_1) \dots o(\vec{e}_k) = o(\vec{e}_1) \dots o(\vec{e}_{i-1}) o(\vec{e}_{i+2}) \dots o(\vec{e}_k)$, which corresponds to the path $\vec{e}_1, \dots, \vec{e}_{i-1}, \vec{e}_{i+2}, \dots, \vec{e}_k$.
    So we can inductively reduce the length of the path until it is a geodesic.
    Therefore, we get $o(\vec{e}_1) \dots o(\vec{e}_k)v = \delta_{\vec{T}, v_0}(v)$.
\end{proof}

\begin{proposition}
    For each $w \in \FF$, we have $\delta_{\vec{T}, v_0} w \delta_{\vec{T}, v_0}^{-1} \in \FF$.
\end{proposition}

\begin{proof}
    Let $\vec{e}_1, \dots, \vec{e}_k$ be the geodesic path from $v_0$ to $wv_0$.
    For a vertex $v \in T$, if $\vec{e}_1', \dots, \vec{e}_m'$ is a geodesic path from $v_0$ to $v$, then we have
    \begin{align*}
        \delta_{\vec{T}, v_0} w \delta_{\vec{T}, v_0}^{-1} (v) &= \delta_{\vec{T}, v_0} (w\delta_{\vec{T}, v_0}^{-1}w^{-1})(wv) \\
        &= \delta_{\vec{T},v_0} \delta_{\vec{T},wv_0}^{-1}(wv) \\
        &=\delta_{\vec{T}, v_0}(o(w\vec{e}_1')^{-1} \dots o(w \vec{e}_m')^{-1}(wv)) \\
        &= o(\vec{e}_1) \dots o(\vec{e}_k) \cdot \\ &\quad\,\, o(w\vec{e}_1') o(o(w\vec{e}_1')^{-1}w\vec{e}_2') \dots o(o(w\vec{e}_1')^{-1} \dots o(w\vec{e}_{m-1}')^{-1}w\vec{e}_m) \cdot \\ &\quad\,\, o(w\vec{e}_1')^{-1} \dots o(w \vec{e}_m')^{-1}(wv) \\
        &= o(\vec{e}_1) \dots o(\vec{e}_k)wv.
    \end{align*}
    Therefore, we conclude that $\delta_{\vec{T}, v_0} w \delta_{\vec{T}, v_0}^{-1} = o(\vec{e}_1) \dots o(\vec{e}_k)w \in \FF$.
\end{proof}

So one considers the conjugation by $\delta_{\vec{T}, v_0}$ on $\FF$ as an automorphism of $\FF$.
Then the outer automorphism represented by the conjugation by $\delta_{\vec{T}, v_0}$ is called a \emph{Dehn twist}, denoted by $\delta_{\vec{T}}$.





\subsection{Dehn twists of doubled handlebodies}





Let $\MM = \MM_g$ denote the connected sum of $g$ copies of $\mathbb{S}^1 \times \mathbb{S}^2$, namely  
\[
    \MM = \MM_g = \#_g \, \mathbb{S}^1 \times \mathbb{S}^2 = \underbrace{(\mathbb{S}^1 \times \mathbb{S}^2) \# \dots \# (\mathbb{S}^1 \times \mathbb{S}^2)}_{g~\text{times}}.
\]
The fundamental group $\pi_1(\MM, x)$, for a chosen basepoint $x \in \MM$, is a free group of rank $g$. 
We henceforth identify $\FF$ with $\pi_1(\MM, x)$.

Let $\MCG(\MM)$ denote the mapping class group of $\MM$, that is, the group of isotopy classes of orientation-preserving homeomorphisms.
Laudenbach \cite{MR356056} shows the following short exact sequence holds:
\[
    1 \longrightarrow (\ZZ / 2 \ZZ)^g \longrightarrow \MCG(\MM) \longrightarrow \Out(\FF) \longrightarrow 1
\]
Brendle, Broaddus and Putman \cite{MR4557874} discover that this sequence is splitting, that is, there exists a section $\Out(\FF) \hookrightarrow \MCG(\MM)$.
It is therefore natural to refer to the image of a Dehn twist in $\MCG(\MM)$ simply as a Dehn twist.
What kind of a homeomorphism is associated with a Dehn twist?
To answer to this question, we need to dive into the geometry of $\MM$.

% For a torus $t$ embedded in $\MM$, we say $t$ is \emph{essential} if 
% \begin{enumerate}
%     \item it does not bound a solid torus and
%     \item \label{def:essential_torus} there exists a simple closed curve in $t$ which is essential in $\MM$.
% \end{enumerate}
% By the definition, an induced map from a fundamental group of an essential torus $t$ to $\FF$ has an infinite cyclic image.
% Thus, the simple closed curve in \eqref{def:essential_torus} is uniquely determined up to homotopy.
% We refer to the homotopy class of this curve as the \emph{core} of $t$.
For a torus $t$ embedded in $\MM$, we say $t$ is \emph{essential} if
\begin{enumerate}
    \item it does not bound a solid torus and
    \item the induced homomorphism $\pi_1(t) \to \pi_1(\MM)$ is nontrivial.
\end{enumerate}
Since $\FF$ contains no subgroup isomorphic to $\ZZ \times \ZZ$, the image of $\pi_1(t) \to \pi_1(\MM)$ is an infinite cyclic subgroup of $\FF$, well defined up to conjugation.
The \emph{core} of $t$ is the conjugacy class of a generator of this image; it is well defined up to inversion.
The kernel of $\pi_1(t) \to \pi_1(\MM)$ is a rank-one direct summand of $\pi_1(t) \cong \ZZ \times \ZZ$, generated by a primitive class of $H_1(t)$; an embedded essential curve on $t$ representing this class is called a \emph{meridian} of $t$, and it is unique up to isotopy on $t$.
If $m$ is a meridian and $\ell$ is an embedded curve on $t$ with $\{[\ell],[m]\}$ a basis of $H_1(t)$, then $\ell$ represents the core; any two such curves are homotopic in $\MM$, since their classes in $\pi_1(t)$ differ by a power of $[m]$, which maps to the identity.
We call such a curve a \emph{core curve} of $t$.

Let $\widetilde\MM$ be a universal cover of $\MM$ with a covering map $\pi: \widetilde\MM \to \MM$.
For an essential torus $t \subset \MM$, the \emph{dual tree} of $t$ is a tree $T$ embedded in $\widetilde\MM$ satisfying
\begin{enumerate}
    \item each component of $\widetilde\MM \setminus \pi^{-1}(t)$ contains exactly one vertex of $T$ and
    \item each component of $\pi^{-1}(t)$ intersects exactly one edge of $T$.
\end{enumerate}
For each $v \in VT$ and $w \in \FF$, if $C$ is the component of $\widetilde\MM$ that contains $v$, we define $wv$ as the vertex contained in $wC$.
This gives a simplicial action of $\FF$ on $T$, which makes $T$ a one-edge $\ZZ$-splitting.

For an orientation on $T$, the Dehn twist $\delta_{\vec{T}}$ along $\vec{T}$ is associated with some mapping class of $\MCG(\MM)$ through the section $\Out(\FF) \hookrightarrow \MCG(\MM)$.
This is represented by the following homeomorphism.

Let \(t\) be an essential torus, and let \(N(t)\) denote a closed regular neighborhood of \(t\).  
Choose an orientation on the core of \(t\), and denote by \(\vec{t}\) the torus \(t\) equipped with this orientation.  

Note $t$ is homeomorphic to $\mathbb{S}^1 \times \mathbb{S}^1$.
Assume $\mathbb{S}^1$ is oriented.
There exists a homeomorphism
\[
    h : N(t) \to [0,1] \times \mathbb{S}^1 \times \mathbb{S}^1
\]
such that the restriction
\[
    h^{-1}\big|_{\{1/2\} \times \{0\} \times \mathbb{S}^1}
\]
is orientation-preserving between \(\{1/2\} \times \{0\} \times \mathbb{S}^1\) and the core of \(t\).
Define a map \[\delta : [0,1] \times \mathbb{S}^1 \times \mathbb{S}^1 \to [0,1] \times \mathbb{S}^1 \times \mathbb{S}^1\] by
\[
    \delta(t, s_0, s_1) = (t, s_0, s_1+t)
\]
for all $(t, s_0, s_1) \in [0,1] \times \mathbb{S}^1 \times \mathbb{S}^1$.
Then the \emph{Dehn twist homeomorphism along $\vec{t}$}, denoted $\delta_{\vec{t}}$, is defined by
\[
    \delta_{\vec{t}}(x) \;=\;
    \begin{cases}
    h^{-1}\,\delta\,h(x), & x \in N(t), \\
    x, & \text{otherwise.}
    \end{cases}
\]

\begin{lemma}[{\cite[Lemma 2.6]{MR3677931}}] \label{lem:gul}
    Let $t$ be an essential torus of $\MM$.
    Let $T$ be the dual tree of $t$.
    For each orientation of $t$, there exists an orientation of $T$ such that $\delta_{\vec{t}}$ represents $\delta_{\vec{T}}$.
\end{lemma}

Note if two essential tori are homotopic, then they give the same Dehn twist as outer automorphisms.
By Brendle, Broaddus and Putman \cite{MR4557874}, this outer automorphism can be considered as a mapping class.
Note the mapping class group consists of isotopy classes of orientation-preserving homeomorphisms.
So the following holds.

\begin{proposition} \label{prop:coincide}
    Let $\vec{t}$ and $\vec{t}'$ be oriented essential tori.
    Then, $\delta_{\vec{t}}$ is isotopic to $\delta_{\vec{t}'}$ if and only if $\vec{t}$ is homotopic to $\vec{t}'$.
\end{proposition}

We postpone the proof of Proposition \ref{prop:coincide} until after Lemma \ref{lem:dehn_disj}.
Throughout this section, we will abuse notation by writing $T$ both for a homotopy class of an essential torus $t$ and for the associated dual tree of $t$.
Similarly, we use $\vec T$ to denote either oriented version.
With this convention, the Dehn twist along $\vec T$ will be denoted by $\delta_{\vec T}$.






\section{Intersection number}







Let $A$ and $B$ be homotopy classes of embedded submanifolds (e.g. simple closed curves, spheres and tori) of $\MM$.
The \emph{intersection number} of $A$ and $B$, denoted by $i(A, B)$, is defined by \[i(A, B) := \min \{ \lvert \pi_0(a \cap b) \rvert \mid a \in A, b \in B \} \]
where $\pi_0(a\cap b)$ denotes the collection of components of $a \cap b$, that is, the zeroth homotopy group of $a \cap b$.

\begin{example}
    If $A$ and $B$ are homotopy classes of essential simple closed curves, then $i(A, B)$ is always zero.
    If $A$ is a homotopy class of an essential simple closed curve and $B$ is a homotopy class of an essential torus or an essential sphere, then $i(A, B)$ can be written by $\min \{ \lvert a \cap b \rvert \mid a \in A, b \in B \}$.
    If $A$ and $B$ are homotopy classes of essential spheres and essential tori, then for each $a \in A$ and $b \in B$, the intersection of $a$ and $b$ is a disjoint union of circles, so $i(A, B)$ indicates the minimum number of circles among such intersections.
\end{example}


\subsection{Bigon, bihedron and minimal position}

Let $c$ be an essential simple closed curve, and let $s$ denote an essential sphere or an essential torus.
We say $c$ and $s$ \emph{form a bigon} if there are arcs $a \subset c$ and $a' \subset s$ such that $a \cup a'$ is a closed curve that bounds a disk.
If $c$ and $s$ form a bigon, one can reduce the number of bigons by homotoping $c$.
This induces the cardinality of $c \cap s$ is strictly larger than the intersection number of $[c]$ and $[s]$.
What we want to mention is that the converse is also true.

\begin{proposition}[Bigon criterion]
    Let $c$ be an essential simple closed curve, and let $s$ be an essential sphere or an essential torus.
    Then, $\lvert c \cap s \rvert = i([c], [s])$ if and only if $c$ and $s$ do not form a bigon.
\end{proposition}

The following proof is adapted from Farb and Margalit \cite{MR2850125}.

\begin{proof}
    Suppose $\lvert c \cap s \rvert > i([c], [s])$.
    Let $c'$ be a simple closed curve that is homotopic to $c$ such that $\lvert c' \cap s \rvert = i([c], [s])$.
    Then there exists a homotopy $h$ that moves $c$ to $c'$.
    Note the domain of $h$ is an annulus $A$.
    Let $\partial_0$ be the boundary component of $A$ whose image is $c$.
    
    Note $h^{-1}(s)$ is a union of arcs whose endpoints are contained in the boundary of $A$.
    Because $c \cap s$ is strictly larger than $c' \cap s$, there exists an arc $a$ of $h^{-1}(s)$ that connects two points of $\partial_0$.
    Then one can find a subarc $a_0 \subset \partial_0$ such that $a \cup a_0$ bounds a disk in $A$.
    That is, the union of $h(a)$ and $h(a_0)$ bounds a disk.
    Hence, $c$ and $s$ form a bigon.
\end{proof}

We may focus on an intersection number of two essential tori.
Before describing how such an intersection can be simplified, we record what the intersection circles themselves can look like.

The point is that they cannot be arbitrary.
An intersection circle lies on both surfaces at once, so its class in $\FF$ lies simultaneously in the images of $\pi_1(s)$ and of $\pi_1(t)$, and these images are severely constrained: a sphere contributes the trivial group, and an essential torus contributes an infinite cyclic group generated by its core.
Since the intersection of two distinct maximal cyclic subgroups of a free group is trivial, one expects such a circle to be null-homotopic in $\MM$, and hence to be as simple as possible on each surface: either bounding a disk, or, on a torus, running parallel to the meridian.

This expectation is correct, but the cyclic-subgroup argument alone does not establish it, since two essential tori may well have commensurable cores.
What rescues the statement is a normal form.
Because $\FF$ contains no free abelian group of rank two, every essential torus is compressible; compressing it along a disk turns it into an essential sphere, and reversing the compression exhibits the torus as that sphere with a single tube attached along an embedded arc.
In this picture the meridians of the torus are precisely the cross-sections of the tube, and the core is the loop closing the arc through the sphere.
Two tori presented in this way interact along their spheres, which contribute only null-homotopic circles, and along their tubes; and since the spines are one-dimensional while $\MM$ is three-dimensional, the spines may be made disjoint, so the tubes cannot meet each other.
A circle carrying a nontrivial class would have to run the length of a tube, which the disjointness forbids.

The following lemma makes this precise.
It requires no hypothesis on the cores, and it will be used throughout to classify the circles of an intersection and, in Lemma \ref{lem:meridian_count}, to count them.

\begin{lemma} \label{lem:circle_types}
    Let $s$ and $t$ be essential surfaces of $\MM$, each a sphere or a torus.
    Then there are transverse representatives, again denoted $s$ and $t$, such that every circle of $s \cap t$ is null-homotopic in $\MM$.
    Consequently each circle of $s \cap t$ either bounds a disk in the surface containing it or, when that surface is a torus, is isotopic on it to the meridian.
\end{lemma}

\begin{proof}
    We make objects transverse wherever the corresponding intersections are intended to be transverse; in particular the final representatives $s$ and $t$ will be transverse to each other.

    If at least one of the two surfaces is a sphere, say $s$, then every circle of $s \cap t$ lies on $s$ and $\pi_1(s) = 1$, so it is null-homotopic in $\MM$ and there is nothing to prove.
    Assume therefore that $s$ and $t$ are essential tori.

    Let $u$ be an essential torus with core $\gamma_u$.
    Since $u$ is essential, the image of $\pi_1(u) \to \FF$ is nontrivial; since it is abelian and every abelian subgroup of a free group is cyclic, the image is infinite cyclic.
    Hence the kernel of $\pi_1(u) \to \FF$ is an infinite cyclic direct summand of $\pi_1(u) \cong \ZZ^2$; let $m \subset u$ be an embedded curve representing a generator of this kernel, that is, a meridian of $u$.

    As $\pi_1(u) \to \pi_1(\MM)$ is not injective and $u$ is two-sided, the loop theorem provides an embedded disk $D \subset \MM$ with $D \cap u = \partial D$ and with $\partial D$ an essential simple closed curve of $u$ lying in the kernel of $\pi_1(u) \to \FF$.
    An essential simple closed curve on a torus represents a primitive class of $\pi_1(u)$, so $\partial D$ represents $m^{\pm 1}$ and is isotopic to $m$ on $u$.
    Adjusting $D$ by this isotopy, we may assume
    \[
        D \cap u = \partial D = m .
    \]

    Compressing $u$ along $D$ yields an embedded sphere $\Sigma$.
    We claim that $\Sigma$ is essential.
    Suppose that $\Sigma$ bounds a ball $B$.
    In the standard compression picture there are two possibilities.
    If $D$ lies on the side of $\Sigma$ opposite $B$, then reversing the compression attaches a $1$--handle to $B$, producing a solid torus bounded by $u$, contrary to essentiality of $u$.
    If instead $D$ lies in $B$, then reversing the compression shows that $u \subset B$; hence every loop on $u$ is null-homotopic in $\MM$, again contrary to essentiality.
    Thus $\Sigma$ does not bound a ball.

    Reversing the compression exhibits $u$ as a sphere with a tube attached: there are disjoint open disks $\mathring{D}^{-}, \mathring{D}^{+} \subset \Sigma$ and an embedded arc $a$ with endpoints in $\partial D^{-} \cup \partial D^{+}$ and interior disjoint from $\Sigma$, such that
    \[
        u = \bigl(\Sigma \setminus (\mathring{D}^{-} \cup \mathring{D}^{+})\bigr) \cup A_u ,
        \qquad A_u \cong S^1 \times [0,1] ,
    \]
    where $A_u$ is the boundary annulus of a thin regular neighborhood of $a$, attached along $\partial D^{-}$ and $\partial D^{+}$.
    We call $A_u$ the \emph{tube} of $u$ and $a$ its \emph{spine}.
    The cross-section circles $S^1 \times \{r\}$ of $A_u$ are meridians of $u$, and $\gamma_u$ is represented by the loop obtained by closing $a$ through $\Sigma$.

    Write accordingly
    \[
        s = \bigl(\Sigma_s \setminus (\mathring{D}_s^{-} \cup \mathring{D}_s^{+})\bigr) \cup s_a ,
        \qquad
        t = \bigl(\Sigma_t \setminus (\mathring{D}_t^{-} \cup \mathring{D}_t^{+})\bigr) \cup t_a ,
    \]
    with spines $a_s, a_t$ and tubes $s_a, t_a \cong S^1 \times [0,1]$, and put $K_s := \Sigma_s \cup a_s$ and $K_t := \Sigma_t \cup a_t$.

    Since $\dim \MM = 3$ and the spines are $1$--dimensional, a small homotopy makes $a_s \cap a_t = \emptyset$.
    A further small homotopy makes $a_s$ transverse to $\Sigma_t$, $a_t$ transverse to $\Sigma_s$, and $\Sigma_s$ transverse to $\Sigma_t$, so that $a_s \cap \Sigma_t$ and $a_t \cap \Sigma_s$ are finite sets of interior points.
    The endpoints of $a_s$ may be slid freely on $\Sigma_s$, so we place them in $\Sigma_s \setminus K_t$, which is nonempty because $\Sigma_s \cap \Sigma_t$ is a union of circles on the sphere $\Sigma_s$ and $\Sigma_s \cap a_t$ is finite; symmetrically for the endpoints of $a_t$.

    Let $J_s \subseteq a_s$ be the closed subarc underlying the tube $s_a$, obtained by deleting small neighborhoods of the two endpoints of $a_s$; since those endpoints are disjoint from the compact set $K_t$, we may choose the deleted neighborhoods disjoint from $K_t$, so that
    \[
        a_s \cap K_t \subset \operatorname{int}(J_s) .
    \]
    Fix an identification $J_s \cong [0,1]$ compatible with the product structure $s_a \cong S^1 \times [0,1]$, and write
    \[
        \pi_s \colon s_a \cong S^1 \times [0,1] \longrightarrow [0,1] \cong J_s
    \]
    for the projection to the second factor, using the same symbol for the product projection of a regular neighborhood of $a_s$ onto $a_s$.

    For each $x \in a_s \cap K_t = a_s \cap \Sigma_t$ choose a closed interval $I_x \subset \operatorname{int}(J_s)$ with $x \in \operatorname{int}(I_x)$, the $I_x$ pairwise disjoint, and put
    \[
        I_s := \bigcup_x I_x \;\subsetneq\; J_s ,
    \]
    so that neither endpoint of $J_s$ lies in $I_s$.
    Then $a_s \setminus \operatorname{int}(I_s)$ is compact and disjoint from $K_t$, hence at positive distance from it for any fixed Riemannian metric on $\MM$.
    Choosing the regular neighborhoods of $K_s$ and $K_t$ inside sufficiently small metric neighborhoods, we obtain
    \[
        \pi_s\bigl(N(a_s) \cap N(K_t)\bigr) \subseteq I_s ,
    \]
    and since $t \subset N(K_t)$,
    \begin{equation} \label{eq:tube_projection_control}
        \pi_s(s_a \cap t) \subseteq I_s \subsetneq J_s .
    \end{equation}
    Shrinking further if necessary, we arrange the symmetric statement for $t_a \cap s$.

    After these isotopies, reconstruct the two sphere-with-tube surfaces using sufficiently thin regular neighborhoods, and continue to denote them by $s$ and $t$; they are isotopic to the original representatives through the preceding modifications.
    We choose the neighborhoods generically so that $s$ and $t$ are transverse.

    Let $c$ be a component of $s \cap t$ and suppose $[c] \neq 1$ in $\FF$.
    Since the image of $\pi_1(s) \to \FF$ is generated by $\gamma_s$, writing the slope of $c$ in a core--meridian basis of $s$ as $(p,q)$ gives $[c] = \gamma_s^{\,p}$; hence $p \neq 0$.

    After a small isotopy of the two boundary circles of $s_a$ we may assume they are transverse to $c$, so that $c \cap s_a$ is a disjoint union of properly embedded arcs and circles.
    Consider
    \[
        [\,c \cap s_a\,] \in H_1\bigl(s_a, \partial s_a\bigr) \cong \ZZ ,
    \]
    the isomorphism being given by algebraic intersection number with a meridian $S^1 \times \{r\}$.
    A circle in $s_a$, and an arc with both endpoints on the same boundary component of $s_a$, each have algebraic intersection number zero with a meridian, so only arcs joining the two boundary components contribute; the total is the algebraic intersection number of $c$ with a meridian of $s$, namely $p$.
    As $p \neq 0$, some component $\beta \subseteq c \cap s_a$ joins the two boundary components of $s_a$.
    Since the endpoints of $\beta$ lie on the two boundary components of $s_a$, its projection contains both endpoints of $J_s$; as $\pi_s(\beta)$ is connected, $\pi_s(\beta) = J_s$.
    On the other hand $\beta \subseteq c \subseteq t$, so $\beta \subseteq s_a \cap t$, and \eqref{eq:tube_projection_control} gives $\pi_s(\beta) \subseteq I_s \subsetneq J_s$.
    This contradiction shows $[c] = 1$.

    Let $c$ be a component of $s \cap t$ and let $u$ be either of the two surfaces, so that $c \subset u$.
    If $u$ is a sphere, then $c$ bounds a disk on $u$.
    Suppose $u$ is a torus.
    If $c$ is inessential on $u$, then $c$ bounds a disk on $u$.
    If $c$ is essential on $u$, then $c$ represents a primitive class of $\pi_1(u)$, of slope $(p,q)$ with $\gcd(p,q) = 1$ in a core--meridian basis, and its image in $\FF$ is $\gamma_u^{\,p}$.
    By the above this image is trivial, and $\gamma_u$ has infinite order, so $p = 0$; primitivity then forces $q = \pm 1$, so $c$ is isotopic on $u$ to the meridian.
\end{proof}

With Lemma \ref{lem:circle_types} in hand, every intersection circle of two essential tori bounds a disk on each surface or is a meridian of one of them, and we may ask how such circles are removed.
There are two ways to reduce the number of circles of two essential tori via a homotopy.

\begin{example}
    Let $t$ and $t'$ be essential tori.
    \begin{enumerate}
        \item Suppose there exist disks $D \subset t$ and $D' \subset t'$ such that $D \cup D'$ forms a sphere bounding a $3$-ball. In this case, $t$ and $t'$ are said to form a \emph{bihedron}. By homotoping $t$ across this $3$-ball, one can eliminate the intersection circle $D \cap D'$, thereby reducing the number of intersection circles between $t$ and $t'$.

        \item Suppose there are arcs $a \subset t$ and $a' \subset t'$ connecting two distinct intersection circles such that $a \cup a'$ bounds a disk. Then $t$ and $t'$ are said to form a \emph{bigon}. By homotoping $t$ across this disk, one can reduce the number of intersection circles.
    \end{enumerate}
\end{example}

The above example applies both to pairs of essential spheres and to pairs consisting of an essential sphere and an essential torus.
This naturally raises the question of whether forming a bigon or a bihedron is the only obstruction to realizing the intersection number.
The following result gives a positive answer, which we refer to as the \emph{bigon--bihedron criterion}.

\begin{proposition}[Bigon--bihedron criterion] \label{prop:bigon-bihedron}
    Let \(s\) and \(t\) be essential spheres, essential tori, or one of each.
    Then, the number of intersection circles of $s$ and $t$ is exactly $i([s], [t])$ if and only if $s$ and $t$ form neither a bigon nor a bihedron.
\end{proposition}

\begin{proof}
    If $s$ and $t$ form either a bigon or a bihedron, then the number of intersection circles can be reduced, as in the preceding example.
    
    Conversely, suppose \(s\) and \(t\) are not in minimal position.
    Then there exists a homotopy 
    \[
        H : \hat{s} \times [0,1] \to \MM
    \]
    that deforms \(s\) into \(s'\) such that \(s'\) and \(t\) are in minimal position.
    We may assume \(H\) is smooth and each intermediate surface \(H_t(\hat{s})\) intersects \(t\) in a disjoint union of circles except for finitely many times.
    
    The preimage \(H^{-1}(t)\) is a (possibly disconnected) orientable surface with boundary.  
    If \(H^{-1}(t)\) contains a disk whose boundary lies in \(s\), then \(s \cup t\) contains a bihedron.  
    Otherwise, since the number of circles in \(s' \cap t\) is strictly smaller than in \(s \cap t\), there exists a connected component \(S \subset H^{-1}(t)\) such that \(S \cap (\hat{s} \times \{0\})\) is a disjoint union of at least two circles.
    
    Choose two such circles, denoted by $C_1$ and $C_2$.
    Then there exist arcs $a \subset \hat{s} \times \{0\}$ and $b \subset S$, each connecting $C_1$ to $C_2$, such that the union $a \cup b$ bounds a disk in $\hat{s} \times [0,1]$.
    Since the union $H(a) \cup H(b)$ also bounds a disk, it follows that $s$ and $t$ form a bigon.
    
    Therefore, if two essential tori are not in minimal position, they necessarily form either a bigon or a bihedron.
\end{proof}

By Laudenbach \cite{MR356056}, two homotopic sphere systems are isotopic.
Hatcher \cite{MR1314940} shows that two sphere systems can be in ``normal form.''
This is equivalent to minimal position.
That is, two sphere systems are in normal form if and only if they are in minimal position.
Refer to Horbez's article \cite{MR3020214}.

The situation is not different if we consider an essential torus and a sphere system.
G{\"u}ltepe \cite{MR3043125} defines the normal form of a torus with respect to a sphere system.
This is equivalent to minimal position.

\begin{theorem}[Minimal position] \label{thm:minimal-position}
    Let $A_1, \dots, A_k$ be homotopy classes of essential simple closed curves, essential spheres and essential tori.
    Then there are representatives $$a_1 \in A_1, \dots, a_k \in A_k$$ such that $$\lvert \pi_0 (a_i \cap a_j) \rvert = i(A_i, A_j)$$ for all distinct $i$ and $j$.
\end{theorem}

\begin{proof}
    We proceed by induction on the number \(k\) of homotopy classes.  
    The statement is already seen to be true for \(k \le 2\).  
    Assume now that \(k > 2\) and that the statement holds for all smaller values of \(k\).  
    
    By the inductive hypothesis, there exist representatives \(a_1, \dots, a_{k-1}\) of \(A_1, \dots, A_{k-1}\), respectively, such that  
    \[
    \lvert \pi_0(a_i \cap a_j) \rvert = i(A_i, A_j)
    \]
    for all distinct $i, j \in \{1, \dots, k-1\}$.
    Choose a representative \(a_k \in A_k\).  
    It suffices to show that if
    \[
    \sum_{i=1}^{k-1} \lvert \pi_0(a_i \cap a_k) \rvert
    >
    \sum_{i=1}^{k-1} i(A_i, A_k),
    \]
    then there exists another representative \(a_k' \in A_k\) satisfying
    \[
    \sum_{i=1}^{k-1} \lvert \pi_0(a_i \cap a_k') \rvert
    <
    \sum_{i=1}^{k-1} \lvert \pi_0(a_i \cap a_k) \rvert.
    \]
    This will produce a representative of \(A_k\) in minimal position with respect to all \(A_i\) for \(i < k\).
    
    \medskip
    \noindent
    \textbf{Case 1.} Some \(a_i\) is an essential simple closed curve.  

    \medskip
    Without loss of generality, assume \(a_k\) is such a curve.  
    Let \(\mathcal{A}\) denote the set of all subarcs \(b \subset a_k\) for which \(b \cup c\) bounds a bigon with some subarc \(c \subset a_i\).  
    The set \(\mathcal{A}\) is partially ordered by inclusion.  
    Choose a minimal element \(b_0 \in \mathcal{A}\) and let \(c_0 \subset a_i\) be such that \(b_0 \cup c_0\) bounds a disk \(D_0\).  
    For each \(j \ne i\), minimality of \(b_0\) implies that every component of \(a_j \cap D_0\) is an arc joining \(b_0\) and \(c_0\).  
    (If a component is a circle or a point, one may shrink \(D_0\) slightly to remove it.)  
    Let \(a_k'\) be the simple closed curve obtained by homotoping \(b_0\) to \(c_0\).  
    Then \(\lvert a_j \cap a_k' \rvert \le \lvert a_j \cap a_k \rvert\) for all \(j \ne i\), as required.
    
    \medskip
    \noindent
    \textbf{Case 2.} None of \(a_1, \dots, a_k\) is an essential simple closed curve.

    \medskip
    The argument is analogous.  
    Suppose \(a_k\) and \(a_i\) form a bihedron for some \(i\); that is, there exist disks \(d_i \subset a_i\) and \(d_k \subset a_k\) such that \(d_i \cup d_k\) is a sphere bounding a \(3\)--ball.  
    Assume \(d_k\) is minimal in the sense that no proper subset of \(d_k\) forms a bigon or a bihedron with any \(a_j\).  
    Let \(B_0\) denote the \(3\)--ball bounded by \(d_i \cup d_k\).  
    By the inductive hypothesis, and by minimality of \(d_k\), each intersection \(a_j \cap B_0\) is a disjoint union of disks and cylinders connecting \(d_i\) and \(d_k\).  
    Let \(a_k'\) be obtained from \(a_k\) by homotoping \(d_k\) to \(d_i\).  
    Then \(\lvert \pi_0(a_j \cap a_k') \rvert \le \lvert \pi_0(a_j \cap a_k) \rvert\) for all \(j \ne i\).  
    The case in which \(a_k\) and \(a_i\) form a bigon is treated similarly and will be omitted.
    
    \medskip
    By iterating this procedure, we obtain a representative \(a_k'' \in A_k\) satisfying  
    \(\lvert \pi_0(a_i \cap a_k'') \rvert = i(A_i, A_k)\) for all \(i \ne k\).  
    Hence, all pairs among \(a_1, \dots, a_{k-1}, a_k''\) are in minimal position, completing the proof.
\end{proof}


\subsection{Intersection number and Dehn twist}

In this section, we analyze when two Dehn twists commute.  
If their representatives have disjoint supports, then the corresponding Dehn twists clearly commute.
We are interested in the converse question: if two Dehn twists commute, must they admit representatives with disjoint supports?

\begin{lemma} \label{lem:meridian_count}
    Suppose $s$ and $t$ are in minimal position, and let $\gamma'$ be the core of $t$.
    Then the number of circles of $s \cap t$ of meridian type on $t$ equals $i(\gamma', [s])$.
\end{lemma}

\begin{proof}
    Let $c_1, \dots, c_n$ be the circles of $s \cap t$ of meridian type on $t$, and let $d_1, \dots, d_r$ be the remaining circles of $s \cap t$, each of which bounds a disk on $t$ by Lemma \ref{lem:circle_types}.
    We construct a core curve $\ell$ of $t$ meeting each $c_i$ in exactly one point and disjoint from every $d_k$.

    The circles $c_1, \dots, c_n$ are disjoint and all of meridian slope on $t$, hence pairwise isotopic on $t$, so cutting $t$ along them yields $n$ annuli $A_1, \dots, A_n$, cyclically arranged.
    Each $d_k$ lies in the interior of some $A_p$ and is null-homotopic in $t$; since $\pi_1(A_p) \to \pi_1(t)$ is injective, $d_k$ bounds a disk contained in $A_p$.
    Removing from $A_p$ the open disks bounded by the outermost such circles leaves a connected surface containing both boundary circles of $A_p$, so we may choose a properly embedded arc $\alpha_p \subset A_p$ joining the two boundary circles and disjoint from every $d_k$.
    Concatenating $\alpha_1, \dots, \alpha_n$ cyclically produces an embedded closed curve $\ell \subset t$ meeting each $c_i$ transversely in exactly one point and disjoint from every $d_k$.
    If $n = 0$ we instead take for $\ell$ any embedded curve on $t$ of non-meridian slope disjoint from the $d_k$, chosen by the same argument applied to $t$ cut along a single meridian.

    Since $\ell$ meets a meridian of $t$ in a single point, $\{[\ell], [m]\}$ is a basis of $H_1(t)$, so $\ell$ is a core curve of $t$ and represents $\gamma'$.
    As $\ell \subset t$, we have $\ell \cap s = \ell \cap (s \cap t)$, and by construction this set consists of exactly one point on each $c_i$; hence
    \[
        \lvert \ell \cap s \rvert = n .
    \]
    Moreover $\ell$ is transverse to $s$: at a point $p \in \ell \cap c_i$ we have $T_p s \cap T_p t = T_p c_i$ by transversality of $s$ and $t$, while $T_p \ell \not\subset T_p c_i$, so $T_p \ell \not\subset T_p s$ and therefore $T_p \ell + T_p s = T_p \MM$.

    We claim that $\ell$ and $s$ form no bigon.
    Suppose they do, with sides $a \subset \ell$ and $b \subset s$, so that $a \cup b$ bounds a disk in $\MM$.
    The endpoints of $a$ are two distinct points of $\ell \cap s$, lying on circles $c_i$ and $c_j$ say.
    Since $\ell$ meets each $c_i$ in exactly one point, distinct points of $\ell \cap s$ lie on distinct circles, so $i \neq j$.
    But $a \subset t$ and $b \subset s$, so $a$ and $b$ are arcs of $t$ and of $s$ joining two distinct circles of $s \cap t$ whose union bounds a disk; that is, $s$ and $t$ form a bigon, contradicting minimal position.

    By the bigon criterion, $\lvert \ell \cap s \rvert = i([\ell], [s]) = i(\gamma', [s])$, and combining with the displayed equality gives $n = i(\gamma', [s])$.
\end{proof}

\begin{corollary} \label{cor:type_invariance}
    The number of circles of $s \cap t$ of each of the types listed in Lemma \ref{lem:circle_types} is independent of the choice of representatives in minimal position.
\end{corollary}

\begin{proof}
    By Lemma \ref{lem:meridian_count} the number of circles of meridian type on $t$ equals $i(\gamma', [s])$, which depends only on the homotopy classes; when $s$ is a torus, the same applies with the roles exchanged.
    The total number of circles is $i([s],[t])$, again depending only on the homotopy classes, and the counts of the remaining types are determined by these together with the total.
\end{proof}

\begin{lemma} \label{lem:dehn_disj}
    Let $T$ be a homotopy class of an essential torus with core $\gamma$, and let $S$ be a homotopy class of an essential sphere or of an essential torus.
    If $i(S, \delta_{\vec T}^{\,k}(S)) = 0$ for all $k > 0$, then $i(S, T) = 0$.
\end{lemma}

\begin{proof}
    We prove the contrapositive: if $i(S, T) > 0$, then $i(S, \delta_{\vec{T}}^k(S)) > 0$ for sufficiently large $k$.
    We treat the case where $S$ is a sphere and the case where $S$ is a torus separately.
    We investigate how essential spheres and essential tori evolve under iteration of a Dehn twist by examining the intersection between an essential surface and its image under a power of the twist.
    Let $T$ be a homotopy class of an essential torus with core $\gamma$.
    Let $k$ be a positive integer.
    
    Fix a homotopy class $S$ of an essential sphere with $i(S,T)>0$, and choose representatives $s\in S$ and $t\in T$ in minimal position.
    By Lemma~\ref{lem:circle_types}, for each component $c$ of $s\cap t$, exactly one of the following holds:
    \begin{enumerate}
        \item\label{enum:sphere-torus:disk} $c$ bounds a disk in $t$; or
        \item\label{enum:sphere-torus:meridian} $c$ represents a meridian of $t$.
    \end{enumerate}
    The number of intersection circles of type \eqref{enum:sphere-torus:meridian} is equal to $i(\gamma,S)$ by Lemma~\ref{lem:meridian_count}.
    Consequently, the number of intersection circles of type \eqref{enum:sphere-torus:disk} is $i(S,T)-i(\gamma,S)$.
    The local modification of a neighborhood of $c$ under the Dehn twist $\delta_{\vec T}^k$ depends essentially on whether $c$ is of type \eqref{enum:sphere-torus:disk} or \eqref{enum:sphere-torus:meridian}.
    
    \begin{figure}
        \centering
        \includegraphics[width=0.8\linewidth]{images/intersection_sphere1.jpeg}
        \caption{$c$ bounds a disk in $t$.}
        \label{fig:intersection_sphere1}
    \end{figure}
    
    If \eqref{enum:sphere-torus:disk} holds, that is, $c$ bounds a disk in $t$, then a neighborhood of $c$ in $s$ is stretched along the core $\gamma$ of $t$ under a positive power of the Dehn twist.
    As a result, this neighborhood intersects $s$ in
    \[
        k\, i(\gamma,S)+2
    \]
    circles.
    See Figure \ref{fig:intersection_sphere1}.
    
    \begin{figure}
        \centering
        \includegraphics[width=0.8\linewidth]{images/intersection_sphere2.jpeg}
        \caption{$c$ represents a meridian of $t$.}
        \label{fig:intersection_sphere2}
    \end{figure}
    
    If \eqref{enum:sphere-torus:meridian} holds, that is, $c$ represents a meridian of $t$, then a neighborhood of $c$ in $s$ is stretched along the entire torus $t$ by a positive power of the Dehn twist.
    In this case, the neighborhood intersects $s$ in
    \[
        k\, i(S,T)-1
    \]
    circles.
    See Figure \ref{fig:intersection_sphere2}.
    
    \begin{figure}
        \centering
        \includegraphics[width=\linewidth]{images/example_intersection.jpeg}
        \caption{The surfaces $s$ and $t$ intersect in three circles. In this configuration, there exists an essential torus homotopic to $\delta_{\vec{t}}(s)$ intersects $s$ in eight circles.}
        \label{fig:example_intersection}
    \end{figure}
    
    From the above local pictures, one obtains a corresponding global configuration.
    Figure~\ref{fig:example_intersection} illustrates how these local modifications assemble into a global intersection pattern.
    Combining the contributions from each type of intersection circle, we obtain the following estimate: for every $k>0$,
    \begin{align*}
        i(S,\delta_{\vec T}^k(S)) &\geq (i(S,T)-i(\gamma,S)) \, (k\,i(\gamma,S)+2) + i(\gamma, S) (k \, i(S,T) - 1) \\
        &= k \, i(\gamma,S) \, (2 \, i(S,T) - i(\gamma, S)) + 2\,i(S,T)-3\,i(\gamma,S).
    \end{align*}
    
    To justify this inequality, it suffices to verify that the resulting global configuration contains neither bigons nor bihedra.
    First, none of the local modifications described above produces a bihedron; consequently, no bihedron appears in the global picture.
    Next, suppose for contradiction that there exists a bigon between $s$ and $s_k$ where $s_k$ is a representative of $\delta_{\vec{T}}^k(S)$ locally described as the above.
    Then the arc of $s_k$ forming one side of the bigon can be decomposed into a concatenation of arcs, each of which lies entirely in either $s$ or $t$.
    It follows that there exists a bigon formed by $s$ and $t$.
    This contradicts the assumption that $s$ and $t$ are in minimal position.
    Therefore, no bigons occur between $s$ and $s_k$, and the inequality follows.
    
    If $i(S,T)>0$ and $k\ge 2$, then
    \[
        i(S,\delta_{\vec T}^k(S))>0.
    \]
    This is the desired conclusion.
    
    Suppose two homotopy classes $T$ and $T'$ of essential tori have a positive intersection number, that is, $i(T, T') > 0$.
    Choose representatives $t \in T$ and $t' \in T'$ in minimal position.
    We say an essential simple closed curve of $t$ is a \emph{meridian} if it bounds a disk in $\MM$.
    Then each intersection circle $c$ of $t$ and $t'$ satisfies one of the following:
    \begin{enumerate}
        \item \label{enum:intersection:1} $c$ bounds a disk in each of $t$ and $t'$;
        \item \label{enum:intersection:2} $c$ bounds a disk in $t$ and represents a meridian of $t'$;
        \item \label{enum:intersection:3} $c$ represents a meridian of $t$ and bounds a disk in $t'$; or
        \item \label{enum:intersection:4} $c$ represents a meridian of both $t$ and $t'$.
    \end{enumerate}
    The type of each intersection circle in the above list is invariant under homotopy.
    More precisely, for any other choice of representatives of $T$ and $T'$ in minimal position, each of the above types of intersection circles occurring in $t\cap t'$ also occurs for the new representatives.
    
    Let us examine how the local picture near an intersection circle changes after applying a positive power of the Dehn twist $\delta_{\vec{T}}^k$ to $T'$.
    Fix representatives $t\in T$ and $t'\in T'$ in minimal position, and let $c$ be a component of $t\cap t'$.
    
    \smallskip
    \noindent
    \textbf{Type \eqref{enum:intersection:1}.}
    Suppose that $c$ bounds a disk in each of $t$ and $t'$.
    Then there exist disks $D\subset t$ and $D'\subset t'$ with $\partial D=\partial D'=c$ such that $D\cup D'$ is an essential sphere; see the left-hand side of Figure~\ref{fig:intersection:1}.
    A representative of $\delta_{\vec{T}}^k(T')$ can be chosen so that, in a neighborhood of $c$, it has the local form shown on the right-hand side of Figure~\ref{fig:intersection:1}.
    In particular, this local modification creates at least two transverse intersection curves with $t'$.
    
    \begin{figure}
        \centering
        \includegraphics[width=0.8\linewidth]{images/intersection1.JPG}
        \caption{An intersection circle bounds a disk in each of $t$ and $t'$.}
        \label{fig:intersection:1}
    \end{figure}
    
    \smallskip
    \noindent
    \textbf{Type \eqref{enum:intersection:2}.}
    Suppose that $c$ bounds a disk in $t$ and represents a meridian of $t'$.
    If $c$ appears as on the left-hand side of Figure~\ref{fig:intersection:2}, then a representative of
    $\delta_{\vec{T}}^k(T')$ has the local form shown on the right-hand side of Figure~\ref{fig:intersection:2}.
    In particular, in this neighborhood the surfaces $\delta_{\vec{T}}^k(t')$ and $t'$ intersect in at least two circles.
    
    \begin{figure}
        \centering
        \includegraphics[width=0.8\linewidth]{images/intersection2.JPG}
        \caption{An intersection circle bounds a disk in $t$ and represents a meridian of $t'$.}
        \label{fig:intersection:2}
    \end{figure}
    
    \smallskip
    \noindent
    \textbf{Type \eqref{enum:intersection:3}.}
    Suppose that $c$ represents a meridian of $t$ and bounds a disk in $t'$.
    Then $t$ and $t'$ meet locally as on the left-hand side of Figure~\ref{fig:intersection:3},
    and a representative of $\delta_{\vec{T}}^k(T')$ has the local form shown on the right-hand side.
    In this case the number of intersection circles of $t'$ and $\delta_{\vec{T}}^k(t')$
    contained in a sufficiently small neighborhood of $c$ grows linearly with $k$.
    More precisely, this number is bounded below by $k\, i(T,T')-1$.
    
    \begin{figure}
        \centering
        \includegraphics[width=0.8\linewidth]{images/intersection3.JPG}
        \caption{An intersection circle represents a meridian of $t$ and bounds a disk in $t'$.}
        \label{fig:intersection:3}
    \end{figure}
    
    \smallskip
    \noindent
    \textbf{Type \eqref{enum:intersection:4}.}
    Suppose that $c$ represents a meridian of both $t$ and $t'$.
    Then $t$ and $t'$ meet locally as on the left-hand side of Figure~\ref{fig:intersection:4},
    and a representative of $\delta_{\vec{T}}^k(T')$ has the local form shown on the right-hand side.
    In this case the number of intersection circles of $t'$ and $\delta_{\vec{T}}^k(t')$
    contained in a sufficiently small neighborhood of $c$ equals $k \, i(T,T')$.
    
    \begin{figure}
        \centering
        \includegraphics[width=0.8\linewidth]{images/intersection4.JPG}
        \caption{An intersection circle represents a meridian of both $t$ and $t'$.}
        \label{fig:intersection:4}
    \end{figure}
    
    \smallskip
    
    We now estimate $i(\delta_{\vec{T}}^k(T'),T')$.
    For each $k>0$, construct a representative $t_k''$ of $\delta_{\vec{T}}^k(T')$ by applying the local modifications described in \eqref{enum:intersection:1}--\eqref{enum:intersection:4} in a neighborhood of every intersection circle of $t\cap t'$.
    In general, the pair $(t_k'',t')$ need not be in minimal position: the local construction may create bigons between $t_k''$ and $t'$.
    However, the number of such bigons is uniformly bounded (independent of $k$), since bigons can only occur in finitely many prescribed neighborhoods determined by the original configuration of $t\cap t'$.
    Therefore, after removing all bigons by homotopy (i.e.\ by repeatedly applying the bigon criterion), the total number of intersection circles decreases by at most a uniform constant.
    
    On the other hand, the local models in Figures~\ref{fig:intersection:3} and \ref{fig:intersection:4} produce at least $k \, i(T, T')^2 - i(T, T')$ intersection circles between $t_k''$ and $t'$ whenever there is at least one intersection circle of type \eqref{enum:intersection:3} or \eqref{enum:intersection:4}.
    Consequently, for $k$ sufficiently large, the intersection number remains positive after bigon removal, that is,
    \[
        i(\delta_{\vec{T}}^k(T'),T')>0
    \]
    for all sufficiently large $k$.
    This is the desired statement.
\end{proof}

One of the implications from Lemma \ref{lem:dehn_disj} is Proposition \ref{prop:coincide}.

\begin{proof}[Proof of Proposition \ref{prop:coincide}]
    Lemma~\ref{lem:gul} implies that a homotopy class of an (oriented) essential torus determines an isotopy class of a Dehn twist.
    Thus, to prove Proposition~\ref{prop:coincide}, it suffices to show that if two Dehn twist homeomorphisms are isotopic, then their corresponding essential tori are homotopic.
    
    Suppose that $t$ and $t'$ are essential tori that are not homotopic, however, in minimal position.
    We claim that there exists an essential sphere $s$ such that either
    \[
        i([s],[t])=0<i([s],[t']) \quad \text{or} \quad i([s],[t])>0=i([s],[t']).
    \]
    
    Let $c$ be a meridian of $t$.
    If $D$ is an embedded disk with $t\cap D=c=\partial D$, then one obtains an essential sphere $s_D$ as a boundary component of a regular neighborhood of $t\cup D$.
    If there exists such a disk $D$ with $i([s_D],[t'])>0$, then $s_D$ satisfies the claim.
    
    Otherwise, that is, if $i([s_D],[t']) = 0$ for all $D$, then there exists an essential sphere $s$ that intersects $t$ exactly at $c$.
    By the bigon--bihedron criterion, $s$ and $t$ are in minimal position.
    Since $t'$ is disjoint from $s$, this sphere $s$ satisfies the claim.
    In conclusion, the claim holds.
    
    By Lemma~\ref{lem:dehn_disj}, there exists $k>0$ such that either
    \[
        i([s],[\delta_{\vec t}^k(s)])=0<i([s],[\delta_{\vec t'}^k(s)]) \quad \text{or} \quad i([s],[\delta_{\vec t}^k(s)])>0=i([s],[\delta_{\vec t'}^k(s)]).
    \]
    Therefore, $\delta_{\vec t}$ is not isotopic to $\delta_{\vec t'}$.
\end{proof}

If two homeomorphisms have disjoint supports, then they commute.
Consequently, their mapping classes also commute.
In other words, mapping classes that admit representatives with disjoint supports must commute.
However, in general, commutativity does not imply the existence of representatives with disjoint supports.
On surfaces, it is known that Dehn twists commute if and only if they admit representatives with disjoint supports (see \cite[Fact~3.6]{MR2850125}).
The following theorem shows that the same phenomenon holds for doubled handlebodies.

\begin{theorem} \label{thm:commuting}
    Let $T$ and $T'$ be homotopy classes of essential tori.
    Then, we have $[\delta_{\vec{T}}, \delta_{\vec{T}'}] = 1$ for any orientations of $T$ and $T'$ if and only if $i(T, T') = 0$.
\end{theorem}

\begin{proof}
    Suppose first that \(i(T,T') = 0\).  
    Then there exist representatives \(t \in T\) and \(t' \in T'\) which are disjoint.  
    Since the supports of \(\delta_{\vec{t}}\) and \(\delta_{\vec{t}'}\) are disjoint, these homeomorphisms commute.  
    It follows that \(\delta_{\vec{T}}\) and \(\delta_{\vec{T}'}\) commute.  
    
    Conversely, assume that \(\delta_{\vec{T}}\) and \(\delta_{\vec{T}'}\) commute.  
    Then for each $k > 0$, we have
    \[
        \delta_{\delta_{\vec{T}'}^k(\vec{T})} = \delta_{\vec{T}'}^k \delta_{\vec{T}} \delta_{\vec{T}'}^{-k} = \delta_{\vec{T}}.
    \]
    So one gets \(\delta_{\vec{T}'}^k(T) = T\) for every $k > 0$.
    By Lemma \ref{lem:dehn_disj}, this implies that \(i(T,T') = 0\).  
    Therefore, \(\delta_{\vec{T}}\) and \(\delta_{\vec{T}'}\) commute if and only if \(i(T,T') = 0\).
\end{proof}

\begin{corollary} \label{cor:disjoint_twists_free_abelian}
    Let $T_1, \dots, T_r$ be pairwise distinct, pairwise disjoint homotopy classes of essential tori.
    Then $\delta_{\vec T_1}, \dots, \delta_{\vec T_r}$ generate a free abelian subgroup of $\Out(\FF)$ of rank $r$.
\end{corollary}

\begin{proof}
    The twists commute pairwise by Theorem \ref{thm:commuting}.
    Suppose $\delta_{\vec T_1}^{N_1} \dots \delta_{\vec T_r}^{N_r} = 1$ with $N_{v_0} \ne 0$ for some $v_0$.
    Raising to the $q$-th power and applying Lemma \ref{lem:disjoint_twist_growth} to any $\alpha$ and $S$,
    \[
        i(\alpha, S)
        = i\bigl(\delta_{\vec T_1}^{q N_1} \dots \delta_{\vec T_r}^{q N_r}(\alpha), S\bigr)
        \;\ge\; q \sum_{v} \lvert N_v \rvert \, i(\alpha, T_v) \, i(\gamma_v, S) - i(\alpha, S)
    \]
    for every $q \ge 1$, whence $\sum_{v} \lvert N_v \rvert \, i(\alpha, T_v) \, i(\gamma_v, S) = 0$.
    All terms are nonnegative, so in particular $i(\alpha, T_{v_0}) \, i(\gamma_{v_0}, S) = 0$ for every $\alpha$ and $S$.

    We contradict this by choosing $\alpha$ and $S$ suitably.
    If every element of $\FF$ were elliptic on the Bass--Serre tree of $T_{v_0}$, then $\FF$, being finitely generated, would fix a vertex and the splitting would be trivial; so some $w \in \FF$ is hyperbolic on $T_{v_0}$.
    Represent $w$ by an embedded circle in $\MM$, which is possible since $\dim \MM = 3$, and let $\alpha$ be its homotopy class; then $i(\alpha, T_{v_0}) > 0$.
    Next, fix a basis of $\FF$ and write $\gamma_{v_0}$ cyclically reduced; some basis element $a$ occurs in it.
    Let $S$ be the nonseparating essential sphere dual to $a$, so that $i(\gamma_{v_0}, S)$ is the number of occurrences of $a^{\pm 1}$, which is positive.
    This contradiction shows $N_v = 0$ for all $v$.
\end{proof}




\section{On the universal cover}

Let $p: \widetilde{\MM} \to \MM$ be a universal cover.
Any complete geodesic Riemannian metric on $\MM$ induces Gromov hyperbolicity on $\widetilde{\MM}$.
So we obtain a Gromov boundary of $\widetilde{\MM}$, denoted by $\partial\widetilde{\MM}$.
Two bi-infinite lines on $\widetilde{\MM}$ are said to be \emph{equivalent} if the Gromov--Hausdorff distance between them is finite.
For a bi-infinite line $\tilde{c}$ and an essential sphere $\tilde{s}$ in $\widetilde{\MM}$, we say $\tilde{c}$ \emph{penetrates} $\tilde{s}$ if $\tilde{c}$ and $\tilde{s}$ intersect without bigon.
For two integers $a$ and $b$, let $\operatorname{lcm}(a, b)$ denote the least common multiple of $a$ and $b$.

\begin{lemma} \label{lem:parallel}
    Let $c$ and $c'$ be essential simple closed curves, and let $s$ be an essential sphere, which are in minimal position.
    Let $\tilde{c}$ and $\tilde{c}'$ be lifts of $c$ and $c'$, respectively.
    If $\tilde{c}$ is not equivalent to $\tilde{c}'$, then the number of lifts of $s$ penetrated by both $\tilde{c}$ and $\tilde{c}'$ is at most $\operatorname{lcm}(i([c],[s]), i([c'],[s]))$.
\end{lemma}

\begin{proof}
    If $\operatorname{lcm}(i([c],[s]), i([c'],[s])) = 0$, then $s$ is disjoint from either $c$ or $c'$, that is, there is no lifts of $s$ penetrated by both $\tilde{c}$ and $\tilde{c}'$.
    So the statement holds in this case.
    Let us assume that $\operatorname{lcm}(i([c],[s]), i([c'],[s])) > 0$.
    
    Suppose, for contradiction, that there are at least $\operatorname{lcm}(i([c],[s]), i([c'], [s])) + 1$ lifts of $s$, each of which penetrated by both $\tilde{c}$ and $\tilde{c}'$.
    Then there are two distinct lifts $\tilde{s}_1$ and $\tilde{s}_2$ of $s$ such that \begin{center}$p(\tilde{c} \cap \tilde{s}_1) = p(\tilde{c} \cap \tilde{s}_2)$ and $p(\tilde{c}' \cap \tilde{s}_1) = p(\tilde{c}' \cap \tilde{s}_2)$.\end{center}
    If $w \in \FF$ is an element that maps $\tilde{s}_1$ to $\tilde{s}_2$, then $w$ preserves $\tilde{c}$ and $\tilde{c}'$.
    So $\tilde{c}$ and $\tilde{c}'$ are equivalent, which contradicts the hypothesis.
    Therefore, the statement holds.
\end{proof}

A bi-infinite cylinder in $\widetilde{\MM}$ is said to be \emph{essential} if it does not bound a solid cylinder.
Note that every lift of an essential torus is an essential cylinder.
We also say a bi-infinite line $\tilde{c}$ \emph{penetrates} an essential cylinder $\tilde{t}$ in $\widetilde{\MM}$ if they intersect without bigon.

\begin{proposition} \label{prop:parallel_torus}
    Let $c$ and $c'$ be essential simple closed curves, and let $t$ be an essential torus, which are in minimal position.
    Let $\tilde{c}$ be a lift of $c$, and let $\tilde{c}'$ be a lift of $c'$.
    If $\tilde{c}$ is not equivalent to $\tilde{c}'$, then the number of lifts of $t$ penetrated by both $\tilde{c}$ and $\tilde{c}'$ is at most $2 \operatorname{lcm}(i([c],[t]), i([c'],[t]))$.
\end{proposition}

Meanwhile, $\FF$ acts properly discontinuously on $\widetilde{\MM}$.
An element of $\FF$ is said to be \emph{indivisible} if it is not a nontrivial power of any other element.
We need the following property of $\FF$ for a proof of Proposition \ref{prop:parallel_torus}.

\begin{theorem}[Lyndon and Sch{\"u}tzenberger \cite{MR162838}] \label{thm:equation}
    Let $x, y, z \in \FF$ be nontrivial indivisible elements.
    Suppose the equation $x^a y^b z^c = 1$ holds for some $a, b, c \geq 2$.
    Then there exists an element $w \in \FF$ such that $x, y, z \in \{ w, w^{-1}\}$.
\end{theorem}

We say an element $g \in \FF$ \emph{minimally preserves} a bi-infinite line $\tilde{c}$ if $h$ is any element of $\FF$ that preserves $\tilde{c}$, then one gets $h = g^n$ for some $n \in \mathbb{Z}$.

\begin{proof}[Proof of Proposition \ref{prop:parallel_torus}]
    Suppose there are at least $2 \operatorname{lcm}(i([c],[t]), i([c'],[t])) + 1$ lifts of $t$ which are penetrated by both $\tilde{c}$ and $\tilde{c}'$.
    Then there are three lifts $\tilde{t}_1$, $\tilde{t}_2$ and $\tilde{t}_3$ of $t$ such that \begin{center}$p(\tilde{c} \cap \tilde{t}_1) = p(\tilde{c} \cap \tilde{t}_2) = p(\tilde{c} \cap \tilde{t}_3)$ and $p(\tilde{c}' \cap \tilde{t}_1) = p(\tilde{c}' \cap \tilde{t}_2) = p(\tilde{c}' \cap \tilde{t}_3)$.\end{center}

    Let $g$ and $g'$ be elements of $\FF$ minimally preserving $\tilde{c}$ and $\tilde{c}'$, respectively.
    Let us take signs for $g$ and $g'$ to satisfy that \begin{center}$g^{m+n}(\tilde{t}_1) = g^{n}(\tilde{t}_2) = \tilde{t}_3$ and $(g')^{m'+n'}(\tilde{t}_1) = (g')^{n'}(\tilde{t}_2) = \tilde{t}_3$\end{center} for some positive integers $m, n, m', n'$.
    Then $(g')^{m'}g^{-m}$ and $(g')^{m'+n'}g^{-(m+n)}$ preserve $\tilde{t}_2$ and $\tilde{t}_3$, respectively.
    If $(g')^{m'}g^{-m}$ is trivial, then $\tilde{c}$ and $\tilde{c}'$ are equivalent, so this contradicts the hypothesis.
    Otherwise, there exists a nontrivial element $h_2 \in \FF$ that preserves $\tilde{t}_2$ such that $h_2 = (g')^{m'}g^{-m}$.

    Similarly, since $g^{-(m+n)}(g')^{m'+n'}$ is nontrivial, there exists a nontrivial element $h_3 \in \FF$ that preserves $\tilde{t}_3$ such that $h_3 = (g')^{m'+n'}g^{-(m+n)}$.
    If $h_3$ is conjugate to $h_2$, then this implies that $\tilde{c}$ and $\tilde{c}'$ are equivalent.
    Otherwise, $h_3$ can be expressed as a nontrivial power of some indivisible element $h$.
    That is, we get an equation $h^n = (g')^{m'+n'}g^{-(m+n)}$ for some $n \geq 2$.
    By Theorem \ref{thm:equation}, $h, g, g'$ are powers of a common element.
    That is, $\tilde{c}$ and $\tilde{c}'$ are equivalent.
    This is a contradiction.
    Therefore, there are at most $2 \operatorname{lcm}(i([c],[t]), i([c'],[t]))$ lifts of $t$ that are penetrated by both $\tilde{c}$ and $\tilde{c}'$.
\end{proof}

\begin{proposition} \label{prop:between}
    Let $c$ be an essential simple closed curve, and let $t$ and $t'$ be essential tori in minimal position.
    Let $\tilde{t}_{-1}, \tilde{t}_0, \tilde{t}_1$ be a consecutive sequence of lifts of $t$ penetrated by some lift $\tilde{c}$ of $c$.
    For each $i$, let $\tilde{c}_i$ be the core of $\tilde{t}_i$.
    Then the number of lifts of $t'$ penetrated by $\tilde{c}_0$ that separate $\tilde{c}_{-1}$ from $\tilde{c}_1$ is at most $i([c], [t'])$.
\end{proposition}

\begin{proof}
    Let $g$ be an element of $\FF$ minimally preserving $\tilde{c}$.
    Let $\tilde{t}'$ be a lift of $t'$ that is penetrated by $\tilde{c}_0$ and separates $\tilde{c}_{-1}$ from $\tilde{c}_1$.
    Because $\tilde{t}_{-1}, \tilde{t}_0, \tilde{t}_1$ are consecutive, $g(\tilde{t}')$ cannot separate $\tilde{c}_{-1}$ and $\tilde{c}_1$.
    So all the lifts of $t'$ that is penetrated by $\tilde{c}_0$ and separates $\tilde{c}_{-1}$ from $\tilde{c}_1$ are distinguished as $g$-orbits.
    So they are at most $i([c], [t'])$.
\end{proof}







\section{Compatible collection and fellow traveler}





In this section, we introduce compatible collections and the notion of fellow travelers.  

\subsection{Compatible collection}

A homotopy class of an essential simple closed curve is said to be \emph{indivisible} if it is not represented by a nontrivial power of some nontrivial element of $\FF$.
A homotopy class of an essential simple closed curve is said to \emph{penetrate} a homotopy class of an essential torus if their intersection number is positive.
Given two homotopy classes \(T\) and \(T'\) of essential tori, if the core of \(T\) penetrates \(T'\), we simply say that \(T\) \emph{penetrates} \(T'\).
We say that \(T\) and \(T'\) \emph{co-penetrate} if they penetrate each other.
A finite collection \(\{T_1, \dots, T_k\}\) of homotopy classes of essential tori is called \emph{compatible} if their cores are pairwise distinct and indivisible, and every pair that intersects also co-penetrates.

\begin{lemma} \label{lem:penetration_noncomm}
    Let $T$ and $T'$ be homotopy classes of essential tori with cores $\gamma$ and $\gamma'$.
    If $\gamma$ and $\gamma'$ are commensurable, then $i(\gamma, T') = i(\gamma', T) = 0$.
    Equivalently, if $T$ penetrates $T'$, then $\gamma$ and $\gamma'$ are non-commensurable.
\end{lemma}

\begin{proof}
    Let $g$ and $g'$ be elements of $\FF$ representing $\gamma$ and $\gamma'$.
    By hypothesis there are nonzero integers $m, n$ and an element $w \in \FF$ with $g^m = w (g')^n w^{-1}$.
    Since $\gamma'$ is the core of the one-edge $\ZZ$-splitting $T'$, the element $(g')^n$ stabilizes an edge of $T'$, hence so does $g^m$, and therefore $\lVert g^m \rVert_{T'} = 0$.
    An element of $\FF$ acting on a simplicial tree satisfies $\lVert g^m \rVert_{T'} = \lvert m \rvert \, \lVert g \rVert_{T'}$, so $i(\gamma, T') = \lVert g \rVert_{T'} = 0$.
    The other equality is symmetric.
\end{proof}

\begin{lemma} \label{lem:compatible_noncomm}
    The cores of a compatible collection are pairwise non-commensurable.
\end{lemma}

\begin{proof}
    Let $\gamma_a$ and $\gamma_b$ be commensurable cores, with representatives $g_a$ and $g_b$, so that $g_a^{m} = w g_b^{n} w^{-1}$ for some $w \in \FF$ and nonzero integers $m, n$.
    Both $g_a$ and $w g_b w^{-1}$ centralize $g_a^{m}$.
    Centralizers of nontrivial elements of $\FF$ are maximal cyclic, so there is $u \in \FF$ with $g_a = u^{p}$ and $w g_b w^{-1} = u^{q}$ for some integers $p, q$.
    Since $\gamma_a$ and $\gamma_b$ are indivisible, $p, q \in \{1, -1\}$, whence $\langle g_a \rangle = \langle w g_b w^{-1} \rangle$ and $\gamma_a = \gamma_b$ as unoriented conjugacy classes.
    This contradicts the assumption that the cores are pairwise distinct.
\end{proof}

\subsection{Penetration and lifts}

Let \(V = \{T_{1}, \dots, T_{k}\}\) be a compatible collection of homotopy classes of essential tori, and let \(\Gamma\) be the \emph{coincidence graph} of \(V\).  
Thus the vertices of \(\Gamma\) are \(T_{1}, \dots, T_{k}\), and \(T_{i}\) is joined to \(T_{j}\) by an edge if and only if \(i(T_{i}, T_{j}) = 0\).  
We adopt standard graph-theoretic notation: for instance, \(\operatorname{lk}(T_{i})\) denotes the link of \(T_{i}\) and \(\operatorname{st}(T_{i})\) its star.

For a homotopy class \(\gamma\) of an essential simple closed curve, a \emph{lift} $\tilde\gamma$ of \(\gamma\) means an equivalence class of lifts of representatives of \(\gamma\).  
Likewise, for a homotopy class \(T\) of an essential torus, a \emph{lift} $\tilde{T}$ of \(T\) refers to an equivalence class of lifts of representatives of \(T\).
We say that \(\tilde{\gamma}\) \emph{penetrates} \(\tilde{T}\) if any representative of \(\tilde{\gamma}\) intersects any representative of \(\tilde{T}\).

Throughout this subsection we fix representatives of the tori of $V$ and of the curves under consideration in simultaneous minimal position, which is possible by Theorem \ref{thm:minimal-position}.
In particular, if $i(T_a, T_b) = 0$, then the chosen representatives $t_a$ and $t_b$ are disjoint.
Recall that each lift of an essential torus corresponds to an edge of the associated dual tree; what we use below is the resulting fact that a lift $\tilde T$ separates $\widetilde\MM$ into two components, which we call its \emph{sides}.

\begin{lemma} \label{lem:geodesic_in_dual_tree}
    Let $T$ be a homotopy class of an essential torus with representative $t$ and dual tree $X_T$, and let $\tilde c$ be a lift of an essential simple closed curve $c$ such that $c$ and $t$ are in minimal position.
    Then the following hold.
    \begin{enumerate}
        \item \label{item:geodesic:disjoint} If $\tilde c$ does not penetrate a lift $\tilde T$ of $t$, then $\tilde c \cap \tilde T = \emptyset$; in particular $\tilde c$ lies in one side of $\tilde T$.
        \item \label{item:geodesic:once} If $\tilde c$ penetrates $\tilde T$, then $\tilde c$ meets $\tilde T$ in exactly one point.
        \item \label{item:geodesic:geodesic} The lifts of $t$ penetrated by $\tilde c$ are precisely the edges of a geodesic of $X_T$, and the order in which $\tilde c$ meets them agrees with the order along that geodesic.
    \end{enumerate}
\end{lemma}

\begin{proof}
    The components of $\widetilde\MM \setminus \rho^{-1}(t)$ are the vertices of $X_T$ and the components of $\rho^{-1}(t)$ are its edges.
    Following $\tilde c$ and recording the components it passes through determines an edge path in $X_T$, in which each intersection point of $\tilde c$ with $\rho^{-1}(t)$ contributes one traversed edge.
    If this path were not reduced, then $\tilde c$ would traverse some edge $\tilde T$ and immediately afterwards traverse it back; choosing an outermost such backtracking, the corresponding subarc of $\tilde c$ together with an arc of $\tilde T$ bounds a disk, that is, $\tilde c$ and $\tilde T$ form a bigon.
    This contradicts minimal position.
    Hence the path is reduced, so it is a geodesic and traverses each edge at most once, which gives \eqref{item:geodesic:once} and \eqref{item:geodesic:geodesic}.
    Finally, an intersection point of $\tilde c$ with a lift $\tilde T$ would force the path to traverse the edge $\tilde T$, so a lift that is not penetrated is disjoint from $\tilde c$, which gives \eqref{item:geodesic:disjoint}.
\end{proof}

\begin{corollary} \label{cor:common_block}
    Let $\tilde c$ and $\tilde c'$ be lifts of essential simple closed curves in minimal position with $t \in T$.
    Then the set of lifts of $t$ penetrated by both $\tilde c$ and $\tilde c'$ is the edge set of a (possibly empty) geodesic segment of $X_T$.
    In particular these lifts are consecutive along $\tilde c$ and along $\tilde c'$, and the two induced orders agree up to reversal.
\end{corollary}

\begin{proof}
    By Lemma \ref{lem:geodesic_in_dual_tree}\eqref{item:geodesic:geodesic} the two sets of penetrated lifts are the edge sets of two geodesics of the tree $X_T$, and the intersection of two geodesics in a tree is a geodesic segment.
\end{proof}

\begin{proposition} \label{prop:two_cylinders}
    Let $T_a$, $T_b$ and $T_c$ be pairwise distinct members of a compatible collection $V$, and fix representatives in minimal position.
    Let $\tilde T_a$ and $\tilde T_b$ be lifts of $T_a$ and $T_b$.
    Then the number of lifts of $T_c$ having non-empty intersection with both $\tilde T_a$ and $\tilde T_b$ is at most
    \[
        7 \, i(T_a, T_c) \, i(T_b, T_c) .
    \]
\end{proposition}

\begin{proof}
    Let $\mathcal{E}$ be the set of lifts of $T_c$ having non-empty intersection with both $\tilde T_a$ and $\tilde T_b$, and let $\langle g_a \rangle = \operatorname{Stab}(\tilde T_a)$ and $\langle g_b \rangle = \operatorname{Stab}(\tilde T_b)$, so that $g_a$ and $g_b$ represent the cores of $T_a$ and $T_b$ and are indivisible.
    If $i(T_a, T_c) = 0$ or $i(T_b, T_c) = 0$, then $\mathcal{E}$ is empty and there is nothing to prove, so assume both are positive.

    The covering $\tilde T_a \to t_a$ is regular with deck group $\langle g_a \rangle$, so the components of $\rho^{-1}(t_c) \cap \tilde T_a$ fall into exactly $i(T_a, T_c)$ orbits under $\langle g_a \rangle$, one for each component of $t_a \cap t_c$.
    Distinct lifts of $T_c$ are disjoint, so a lift meeting $\tilde T_a$ is determined by any one of the components it contains; hence the lifts of $T_c$ meeting $\tilde T_a$ fall into at most $i(T_a, T_c)$ orbits under $\langle g_a \rangle$.
    Symmetrically, those meeting $\tilde T_b$ fall into at most $i(T_b, T_c)$ orbits under $\langle g_b \rangle$.

    Suppose, for contradiction, that $\# \mathcal{E} > 7 \, i(T_a, T_c) \, i(T_b, T_c)$.
    Sorting the elements of $\mathcal{E}$ by their $\langle g_a \rangle$--orbit and their $\langle g_b \rangle$--orbit, some pair of orbits contains at least eight elements $E_0, E_1, \ldots, E_7$.
    Write $\langle h \rangle = \operatorname{Stab}(E_0)$, so that $h$ is a conjugate of the core of $T_c$ and is therefore indivisible.
    For each $s$ choose integers $a_s$ and $b_s$ with
    \[
        E_s = g_a^{a_s} E_0 = g_b^{b_s} E_0 ,
    \]
    so that $g_b^{-b_s} g_a^{a_s}$ preserves $E_0$ and hence equals $h^{c_s}$ for a unique integer $c_s$.
    Normalizing, $a_0 = b_0 = c_0 = 0$.

    The three sequences consist of pairwise distinct integers.
    If $a_s = a_{s'}$ for $s \ne s'$, then $E_s = g_a^{a_s} E_0 = g_a^{a_{s'}} E_0 = E_{s'}$, contradicting that $E_0, \ldots, E_7$ are distinct; the case $b_s = b_{s'}$ is identical.
    If $c_s = c_{s'}$ for $s \ne s'$, then $g_b^{-b_s} g_a^{a_s} = g_b^{-b_{s'}} g_a^{a_{s'}}$, so $g_b^{\,b_{s'} - b_s} = g_a^{\,a_{s'} - a_s}$ with both exponents nonzero by the previous case; hence the cores of $T_a$ and $T_b$ are commensurable, contradicting Lemma \ref{lem:compatible_noncomm}.

    Consequently, among the seven indices $s = 1, \ldots, 7$, at most two satisfy $\lvert a_s \rvert \le 1$, since $a_s \ne 0$ and the $a_s$ are distinct; likewise at most two satisfy $\lvert b_s \rvert \le 1$ and at most two satisfy $\lvert c_s \rvert \le 1$.
    As $7 > 2 + 2 + 2$, there is an index $s$ with
    \[
        \lvert a_s \rvert \ge 2, \qquad \lvert b_s \rvert \ge 2, \qquad \lvert c_s \rvert \ge 2 .
    \]
    For this index, $g_b^{b_s} h^{c_s} g_a^{-a_s} = 1$.
    Replacing $g_b$, $h$ or $g_a$ by its inverse so that all three exponents become positive, this is an equation $x^{\alpha} y^{\beta} z^{\varepsilon} = 1$ with $\alpha, \beta, \varepsilon \ge 2$ in the indivisible elements $x, y, z \in \{g_b^{\pm 1}, h^{\pm 1}, g_a^{\pm 1}\}$.
    By Theorem \ref{thm:equation} there is $w \in \FF$ with $x, y, z \in \{w, w^{-1}\}$, so the cores of $T_a$, $T_b$ and $T_c$ are pairwise commensurable, contradicting Lemma \ref{lem:compatible_noncomm}.
\end{proof}

For a compatible collection $V$ with coincidence graph $\Gamma$, put
\[
    D(\Gamma) := 7 \, \max_{a \neq b} i(T_a, T_b)^2 .
\]

\begin{lemma} \label{lem:end_block}
    Let $T_i \in V$, let $T_j \in \Gamma \setminus \operatorname{st}(T_i)$ and let $T_\ell \in \operatorname{lk}(T_i)$.
    Let $\gamma$ be a homotopy class of an essential simple closed curve, let $\tilde\gamma$ be a lift of $\gamma$ and let $\tilde\gamma_i$ be a lift of $\gamma_i$, and let $\tilde T_i$ be the lift of $T_i$ with core $\tilde\gamma_i$.
    Let
    \[
        E_1, \ldots, E_J
    \]
    be the lifts of $T_j$ penetrated by both $\tilde\gamma$ and $\tilde\gamma_i$, indexed in the order in which $\tilde\gamma$ penetrates them, as in Corollary \ref{cor:common_block}.
    Let $\tilde T_\ell$ be a lift of $T_\ell$ penetrated by $\tilde\gamma$.
    Then the set
    \[
        I_\ell' = \{\, r \in \{1, \ldots, J\} \mid \tilde T_\ell \cap E_r \ne \emptyset \,\}
    \]
    is an interval of $\{1, \ldots, J\}$ which, if nonempty, contains $1$ or $J$; moreover $\# I_\ell' \le D(\Gamma)$.
\end{lemma}

\begin{proof}
    Orient $\tilde\gamma$ so that it penetrates $E_1, \ldots, E_J$ in this order.
    For each $r$, let $P_r$ be the side of $E_r$ containing the initial ray of $\tilde\gamma$ and let $Q_r$ be the other side.
    Distinct lifts of $T_j$ are disjoint, so each $E_{r'}$ with $r' \ne r$ is connected and disjoint from $E_r$, and it meets $\tilde\gamma$ in a single point by Lemma \ref{lem:geodesic_in_dual_tree}\eqref{item:geodesic:once}; hence $E_{r'} \subset P_r$ for $r' < r$ and $E_{r'} \subset Q_r$ for $r' > r$.

    Since $T_\ell \in \operatorname{lk}(T_i)$, the chosen representatives of $T_i$ and $T_\ell$ are disjoint, hence $\tilde T_i \cap \tilde T_\ell = \emptyset$; in particular $\tilde\gamma_i \cap \tilde T_\ell = \emptyset$.
    Let $A$ be the side of $\tilde T_\ell$ containing $\tilde\gamma_i$ and let $B$ be the other side.
    By Lemma \ref{lem:geodesic_in_dual_tree}\eqref{item:geodesic:once}, the curve $\tilde\gamma$ meets $\tilde T_\ell$ in a single point $x$, and the two rays of $\tilde\gamma \setminus \{x\}$ lie in $A$ and in $B$ respectively.
    Reversing the indexing if necessary, we may assume that the ray lying in $A$ is the initial one; we show that $I_\ell'$ contains $J$ whenever it is nonempty.

    We claim $I_\ell'$ is an interval.
    Let $a < c < a'$ with $a, a' \in I_\ell'$, and suppose $c \notin I_\ell'$.
    Then $\tilde T_\ell$ is connected and disjoint from $E_c$, hence contained in $P_c$ or in $Q_c$.
    But $\tilde T_\ell$ meets $E_a \subset P_c$ and also meets $E_{a'} \subset Q_c$, a contradiction.

    We claim $I_\ell'$ contains $J$.
    Suppose $I_\ell' \ne \emptyset$ and $J \notin I_\ell'$, and put $b := \max I_\ell' < J$.
    Since $E_J$ is connected and disjoint from $\tilde T_\ell$, it lies in $A$ or in $B$; as $E_J$ meets $\tilde\gamma_i \subset A$, we get $E_J \subset A$.
    The points of $\tilde\gamma$ lying in $A$ are exactly those of the initial ray, so the point $\tilde\gamma \cap E_J$ precedes $x$, and therefore $x \in Q_J$.
    On the other hand $\tilde T_\ell$ is connected, disjoint from $E_J$, and meets $E_b \subset P_J$, so $\tilde T_\ell \subset P_J$ and hence $x \in P_J$.
    This is a contradiction.

    For $r \in I_\ell'$ the lift $E_r$ has non-empty intersection with $\tilde T_\ell$ by definition, and with $\tilde T_i$ because $\tilde\gamma_i \subset \tilde T_i$ penetrates $E_r$.
    Thus $I_\ell'$ injects into the set of lifts of $T_j$ meeting both $\tilde T_i$ and $\tilde T_\ell$, and Proposition \ref{prop:two_cylinders} bounds the latter by $D(\Gamma)$.
\end{proof}

\begin{corollary} \label{cor:end_interval}
    In the situation of Lemma \ref{lem:end_block}, the set $I_\ell'$ is contained in one of the two end intervals
    \[
        \{1, \ldots, D(\Gamma)\}
        \qquad\text{or}\qquad
        \{J - D(\Gamma) + 1, \ldots, J\} .
    \]
\end{corollary}

\begin{proof}
    By Lemma \ref{lem:end_block} the set $I_\ell'$ is an interval of length at most $D(\Gamma)$ containing $1$ or $J$.
\end{proof}

\subsection{Fellow travelers}

For a positive integer $N$, we say a homotopy class $\gamma$ of an essential simple closed curve \emph{$N$--fellow-travels with $T_i$ (with respect to $\Gamma$)} if 
\begin{enumerate}
    \item $\gamma$ penetrates $T_i$ and
    \item there are lifts $\tilde\gamma$ of $\gamma$ and $\tilde\gamma_i$ of $\gamma_i$ such that for each $T_j \in \Gamma \setminus \operatorname{st}(T_i)$, at least $N$ lifts of $T_j$ are penetrated by both $\tilde\gamma$ and $\tilde\gamma_i$.
\end{enumerate}
For each $i$, let $\mathcal{FT}\!_{N}(T_i, \Gamma)$ denote the collection of $N$--fellow travelers of $T_i$ with respect to $\Gamma$.
By definition, the sequence \[\mathcal{FT}\!_1(T_i, \Gamma) \supseteq \mathcal{FT}\!_2(T_i,\Gamma) \supseteq \mathcal{FT}\!_3(T_i, \Gamma) \supseteq \dots\] is satisfied.
% Define
% \[
%     \Delta(\Gamma) := \max_{i,j,\ell = 1,\dots,k} \operatorname{lcm}(i(\gamma_{i}, T_{j}),\, i(\gamma_{\ell}, T_{j})),
% \]
% where \(\operatorname{lcm}\) denotes the least common multiple.
Then the following holds.


\begin{lemma} \label{lem:fellow-travel}
    Let $\gamma$ be a homotopy class of an essential simple closed curve that $(2\,D(\Gamma)+3)$--fellow-travels with $T_i$.
    Then there exist lifts $\tilde\gamma$ of $\gamma$ and $\tilde\gamma_i$ of $\gamma_i$ such that for every $T_j \in \Gamma \setminus \operatorname{st}(T_i)$ there exist consecutive lifts $\tilde T_j, \tilde T_j', \tilde T_j''$ of $T_j$ with the property that
    \begin{enumerate}
        \item each of $\tilde T_j, \tilde T_j', \tilde T_j''$ is penetrated by both $\tilde\gamma$ and $\tilde\gamma_i$, and
        \item whenever $\tilde T_\ell$ is a lift of some $T_\ell \in \operatorname{lk}(T_i)$ penetrated by $\tilde\gamma$, the cylinder $\tilde T_\ell$ is disjoint from each of $\tilde T_j, \tilde T_j', \tilde T_j''$.
    \end{enumerate}
\end{lemma}

\begin{proof}
    Fix representatives in simultaneous minimal position and choose lifts $\tilde\gamma$ and $\tilde\gamma_i$ as in the definition of a fellow traveler.
    Let $T_j \in \Gamma \setminus \operatorname{st}(T_i)$.
    By Corollary \ref{cor:common_block} the lifts of $T_j$ penetrated by both $\tilde\gamma$ and $\tilde\gamma_i$ are consecutive; write them as $E_1, \ldots, E_J$ with $J \ge 2\,D(\Gamma) + 3$, indexed in the order in which $\tilde\gamma$ penetrates them.

    By Corollary \ref{cor:end_interval}, for every $T_\ell \in \operatorname{lk}(T_i)$ and every lift $\tilde T_\ell$ of $T_\ell$ penetrated by $\tilde\gamma$, the indices $r$ with $\tilde T_\ell \cap E_r \ne \emptyset$ all lie in $\{1, \ldots, D(\Gamma)\}$ or all lie in $\{J - D(\Gamma) + 1, \ldots, J\}$.
    Hence every such $\tilde T_\ell$ is disjoint from each of the middle lifts
    \[
        E_{D(\Gamma)+1}, \ldots, E_{J - D(\Gamma)} .
    \]
    Since $J - 2\,D(\Gamma) \ge 3$, there are at least three of these, and they are consecutive.
    Taking $\tilde T_j, \tilde T_j', \tilde T_j''$ to be three consecutive middle lifts proves the lemma.
\end{proof}

\subsection{Dehn twists and fellow travelers}

For each \(i\), we write \(\delta_i\) for the Dehn twist \(\delta_{\vec{T}_i}\) along \(T_i\), choosing some orientation of \(T_i\).
Define
\[
    M(\Gamma) := \max_{i, j = 1, \dots, k} i(\gamma_i, T_j),
    \qquad
    I(\Gamma) := \max_{a \ne b} i(T_a, T_b),
\]
so that $D(\Gamma) = 7\, I(\Gamma)^2$ in the notation of the previous subsection.
The next proposition shows that, after applying a suitable power of \(\delta_j\) to certain homotopy classes of essential simple closed curves, the resulting curve fellow-travels with \(T_j\) for a long time.
As a consequence, each set \(\mathcal{FT}\!_{2\,D(\Gamma)+3}(T_j,\Gamma)\) contains infinitely many homotopy classes of essential simple closed curves.

\begin{lemma} \label{lem:core_lifts_inequivalent}
    Let $T \in V$ and let $\tilde T \ne \tilde T'$ be distinct lifts of $T$, with cores $\tilde\gamma$ and $\tilde\gamma'$.
    Then $\tilde\gamma$ and $\tilde\gamma'$ are not equivalent.
\end{lemma}

\begin{proof}
    Suppose they are equivalent, so they have the same pair of endpoints in $\partial\widetilde\MM$.
    Write $\operatorname{Stab}(\tilde T) = \langle g \rangle$ and $\operatorname{Stab}(\tilde T') = \langle g' \rangle$; these are the stabilizers of the two lifts and $g, g'$ are conjugates of the core of $T$, hence indivisible.
    Both $g$ and $g'$ fix the common endpoint pair, so they lie in a common maximal cyclic subgroup $\langle u \rangle$ of $\FF$, and indivisibility gives $\langle g \rangle = \langle u \rangle = \langle g' \rangle$.
    Thus $\tilde T$ and $\tilde T'$ are distinct lifts with the same stabilizer, so $u$ preserves both.
    Since $\FF$ acts on the dual tree $X_T$ with $\langle g \rangle$ the stabilizer of the edge $\tilde T$, and edge stabilizers of distinct edges of a one-edge $\ZZ$--splitting are distinct subgroups unless the edges coincide, we obtain $\tilde T = \tilde T'$, a contradiction.
\end{proof}

\begin{lemma} \label{lem:orbit_separation}
    Let $T_j, T_\ell \in V$ be co-penetrating, let $\tilde T_j'$ be a lift of $T_j$ with core $\tilde\gamma_j'$, and let $g$ generate $\operatorname{Stab}(\tilde T_j')$.
    Let $\tilde\gamma''$ be a lift of the core of some lift of $T_j$ distinct from $\tilde T_j'$, and let $\tilde T_\ell$ be a lift of $T_\ell$ penetrated by $\tilde\gamma_j'$.
    Then there is $a_0 \in \ZZ$ such that for every $b, c > 0$ the lift $g^{a_0} \tilde T_\ell$ separates $g^{-b} \tilde\gamma''$ from $g^{c} \tilde\gamma''$, provided no translate $g^a \tilde T_\ell$ is penetrated by $\tilde\gamma''$; and if some $g^{a_1} \tilde T_\ell$ is penetrated by $\tilde\gamma''$, then $g^{a_1} \tilde T_\ell$ separates $g^{-b-1}\tilde\gamma''$ from $g^{c+1}\tilde\gamma''$ for all $b, c > 0$.
\end{lemma}

\begin{proof}
    Let $X_\ell$ be the dual tree of $T_\ell$ and put $\tau := i(\gamma_j, T_\ell)$, which is positive because $T_j$ and $T_\ell$ co-penetrate.
    Let $A \subset X_\ell$ be the geodesic determined by $\tilde\gamma_j'$ as in Lemma \ref{lem:geodesic_in_dual_tree}\eqref{item:geodesic:geodesic}.
    Since $g$ preserves $\tilde\gamma_j'$, the geodesic $A$ is $g$--invariant, and $g$ acts on $A$ by a translation of length $\lVert g \rVert_{X_\ell} = \tau$.
    Thus $A$ is the axis of $g$.
    Let $e \subset A$ be the edge corresponding to $\tilde T_\ell$, and write $e_a := g^a e$ for $a \in \ZZ$, so that $e_a$ corresponds to $g^a \tilde T_\ell$.
    Let $B \subset X_\ell$ be the geodesic determined by $\tilde\gamma''$; for each $m \in \ZZ$, the geodesic determined by $g^m \tilde\gamma''$ is $g^m B$.

    By Lemma \ref{lem:core_lifts_inequivalent}, the lifts $\tilde\gamma_j'$ and $\tilde\gamma''$ are not equivalent, so Proposition \ref{prop:parallel_torus} applied to them and to $T_\ell$ gives
    \begin{equation} \label{eq:orbit_sep_bound}
        \# E(A \cap B)
        \;\le\; 2 \operatorname{lcm}\bigl( i(\gamma_j, T_\ell), \, i(\gamma_j, T_\ell) \bigr)
        \;=\; 2\tau .
    \end{equation}
    In particular $A \cap B$ is a finite geodesic segment.
    Orient $A$ in the direction in which $g$ translates, and speak accordingly of the positive and negative side of an edge of $A$.
    We treat the two cases of the statement separately.

    \medskip
    \noindent
    \textbf{Case 1.} No edge $e_a$ is contained in $B$.

    \medskip
    Let $S \subset A$ be the nearest-point projection of $B$ to $A$.
    Since $A \cap B$ is finite, $S$ is a finite segment: it equals $A \cap B$ when this is nonempty, and otherwise it is the endpoint on $A$ of the unique bridge joining $A$ to $B$.
    Because $B$ contains no edge $e_a$, there is an integer $a_0$ such that $S$ lies between the consecutive orbit edges $e_{a_0 - 1}$ and $e_{a_0}$.

    We claim that $g^m B$ is disjoint from $e_{a_0}$ for every $m \ne 0$.
    Indeed, $e_{a_0} \subset g^m B$ holds if and only if $e_{a_0 - m} \subset B$, which is excluded by hypothesis.
    Since $g^m B$ is connected, it therefore lies entirely on one side of $e_{a_0}$, and since the projection of $g^m B$ to $A$ is $g^m S$, that side is the negative one for $m < 0$ and the positive one for $m > 0$.
    Consequently $e_{a_0}$ separates $g^{-b} B$ from $g^{c} B$ for all $b, c > 0$.
    Translating back to $\widetilde\MM$, the lift $g^{a_0} \tilde T_\ell$ separates $g^{-b} \tilde\gamma''$ from $g^{c} \tilde\gamma''$ for all $b, c > 0$.

    \medskip
    \noindent
    \textbf{Case 2.} Some edge $e_{a_1}$ is contained in $B$.

    \medskip
    We first claim that
    \begin{equation} \label{eq:orbit_sep_gap}
        e_{a_1 + k} \not\subset B
        \qquad \text{whenever } \lvert k \rvert \ge 2 .
    \end{equation}
    Suppose $e_{a_1 + k} \subset B$ with $\lvert k \rvert \ge 2$.
    Then $e_{a_1}$ and $e_{a_1 + k}$ both lie in $A \cap B$, and the intersection of two geodesics in a tree is convex, so $A \cap B$ contains the whole segment of $A$ joining them.
    Since $g$ translates $A$ by $\tau$ edges, that segment contains at least $2\tau + 1$ edges, contradicting \eqref{eq:orbit_sep_bound}.
    This proves \eqref{eq:orbit_sep_gap}.

    Fix $m \ge 2$.
    The geodesic $g^m B$ contains $g^m e_{a_1} = e_{a_1 + m}$, which lies on the positive side of $e_{a_1}$.
    Moreover $e_{a_1} \subset g^m B$ would give $e_{a_1 - m} \subset B$, which \eqref{eq:orbit_sep_gap} forbids; hence $g^m B$ is disjoint from $e_{a_1}$, and being connected it lies entirely on the positive side.
    Symmetrically, for $m \le -2$ the geodesic $g^m B$ contains $e_{a_1 + m}$, which lies on the negative side of $e_{a_1}$, and is disjoint from $e_{a_1}$ by \eqref{eq:orbit_sep_gap}, so it lies entirely on the negative side.
    Therefore $e_{a_1}$ separates $g^{m} B$ with $m \le -2$ from $g^{m'} B$ with $m' \ge 2$.
    Taking $m = -b-1$ and $m' = c+1$ with $b, c > 0$ and returning to $\widetilde\MM$, the lift $g^{a_1} \tilde T_\ell$ separates $g^{-b-1} \tilde\gamma''$ from $g^{c+1} \tilde\gamma''$ for all $b, c > 0$.
\end{proof}

\begin{lemma} \label{lem:core_shares_lifts}
    Let $T_i, T_j \in V$ be co-penetrating.
    Then there is a lift $\tilde\gamma_i$ of $\gamma_i$ penetrating infinitely many lifts of $T_j$.
\end{lemma}

\begin{proof}
    Since $T_i$ penetrates $T_j$, we have $i(\gamma_i, T_j) > 0$, so some lift $\tilde\gamma_i$ penetrates some lift $\tilde T_j$ of $T_j$.
    Let $g$ generate $\operatorname{Stab}(\tilde\gamma_i)$.
    Then $g^m \tilde T_j$ is penetrated by $\tilde\gamma_i$ for every $m$, and these lifts are pairwise distinct: if $g^m \tilde T_j = g^{m'} \tilde T_j$ with $m \ne m'$, then $g^{m - m'}$ stabilizes $\tilde T_j$, so the cores of $T_i$ and $T_j$ are commensurable, contradicting Lemma \ref{lem:compatible_noncomm}.
\end{proof}

\begin{proposition} \label{prop:co-penetrating}
    If an integer $n$ satisfies $\lvert n \rvert \geq M(\Gamma) + 2\,D(\Gamma) + 8$, then we have
    \[
        \delta_j^n\bigl(\mathcal{FT}\!_3(T_i, \Gamma) \cup \{\gamma_i\}\bigr) \subseteq \mathcal{FT}\!_{2\,D(\Gamma)+3}(T_j, \Gamma)
    \]
    for all co-penetrating pairs $T_i$ and $T_j$.
\end{proposition}

\begin{proof}
    Let $\gamma \in \mathcal{FT}\!_3(T_i, \Gamma) \cup \{\gamma_i\}$ be given.
    If $\gamma \in \mathcal{FT}\!_3(T_i, \Gamma)$, then by definition there are lifts $\tilde\gamma$ of $\gamma$ and $\tilde\gamma_i$ of $\gamma_i$ penetrating at least three lifts of $T_j$ in common.
    If $\gamma = \gamma_i$, take $\tilde\gamma = \tilde\gamma_i$ as in Lemma \ref{lem:core_shares_lifts}, which penetrates infinitely many lifts of $T_j$.
    In either case, by Corollary~\ref{cor:common_block} the common lifts are consecutive; denote three of them by $\tilde T_j, \tilde T_j', \tilde T_j''$, in the order in which $\tilde\gamma$ penetrates them.
    Let $T_\ell \in \Gamma \setminus \operatorname{st}(T_j)$, so that $T_j$ and $T_\ell$ co-penetrate.
    Let $\tilde\gamma_j, \tilde\gamma_j', \tilde\gamma_j''$ be the cores of $\tilde{T}_j, \tilde{T}_j', \tilde{T}_j''$, respectively.
    Let $\tilde\delta_j$ be a lift of $\delta_j$ that preserves $\tilde{T}_j$ and $\tilde{T}_j'$.
    Then there exists a nontrivial element $g \in \FF$ preserving $\tilde\gamma_j'$ such that $\tilde\delta_j^n(\tilde\gamma_j'') = g^n \tilde\gamma_j''$.

    We claim that the number of lifts of $T_\ell$ that are penetrated by $\tilde{\gamma}_j'$ and that separate $\tilde{\gamma}_j''$ from $g^{\,n}\tilde{\gamma}_j''$ is at least
    \[
        (\lvert n \rvert - 3)\, i(\gamma_j, T_\ell).
    \]
    Without loss of generality, assume $n > 0$; the case $n < 0$ follows by the same argument applied to $g^{-1}$ in place of $g$.
    Let $\tilde{T}_\ell$ be any lift of $T_\ell$ penetrated by $\tilde{\gamma}_j'$.

    If $g^a\tilde{T}_\ell$ is never penetrated by $\tilde{\gamma}_j''$ for any $a \in \ZZ$, then there exists $a_0 \in \ZZ$ such that $g^{a_0}\tilde{T}_\ell$ separates $g^{-b}\tilde{\gamma}_j''$ from $g^{c}\tilde{\gamma}_j''$ for all $b, c > 0$.
    Consequently,
    \[
        g^{a_0}\tilde{T}_\ell,\ g^{a_0+1}\tilde{T}_\ell,\ \dots,\ g^{a_0+n-1}\tilde{T}_\ell
    \]
    all separate $\tilde{\gamma}_j''$ from $g^{\,n}\tilde{\gamma}_j''$.
    If instead $g^{a_1}\tilde{T}_\ell$ is penetrated by $\tilde{\gamma}_j''$ for some $a_1 \in \ZZ$, then by Lemma~\ref{lem:orbit_separation} the lift $g^{a_1}\tilde{T}_\ell$ separates $g^{-b-1}\tilde\gamma_j''$ from $g^{c+1}\tilde\gamma_j''$ for all $b, c > 0$.
    It follows that
    \[
        g^{a_1+2}\tilde{T}_\ell,\ g^{a_1+3}\tilde{T}_\ell,\ \dots,\ g^{a_1+n-2}\tilde{T}_\ell
    \]
    all separate $\tilde{\gamma}_j''$ from $g^{\,n}\tilde{\gamma}_j''$.

    Since the lifts of $T_\ell$ penetrated by $\tilde{\gamma}_j'$ consist of exactly $i(\gamma_j, T_\ell)$ distinct $\langle g \rangle$--orbits, the total number of such separating lifts is at least
    \[
        (n-3)\, i(\gamma_j, T_\ell) = (\lvert n \rvert - 3)\, i(\gamma_j, T_\ell),
    \]
    which proves the claim.

    Let $\mathcal{T}$ be the set of lifts of $T_\ell$ that are penetrated by $\tilde{\gamma}_j'$ and that separate $\tilde{\gamma}_j''$ from $g^{\,n}\tilde{\gamma}_j''$.
    By Proposition \ref{prop:parallel_torus}, at most $2\, i(\gamma_j, T_\ell)$ lifts in $\mathcal{T}$ are penetrated by $\tilde{\gamma}_j$.
    If a lift in $\mathcal{T}$ is not penetrated by $\tilde{\gamma}_j$, then it separates $\tilde{\gamma}_j$ from either $\tilde{\gamma}_j''$ or $g^{\,n}\tilde{\gamma}_j''$.
    By Proposition \ref{prop:between}, the number of such lifts separating $\tilde{\gamma}_j$ from $\tilde{\gamma}_j''$ is at most $i(\gamma_i, T_\ell)$.
    Therefore the number of lifts in $\mathcal{T}$ separating $\tilde{\gamma}_j$ from $g^{\,n}\tilde{\gamma}_j''$ is at least
    \[
        (\lvert n \rvert - 3)\, i(\gamma_j, T_\ell) - 2\, i(\gamma_j, T_\ell) - i(\gamma_i, T_\ell)
        \;\ge\; \lvert n \rvert - 5 - M(\Gamma)
        \;\ge\; 2\,D(\Gamma) + 3,
    \]
    where the middle inequality uses $i(\gamma_j, T_\ell) \ge 1$, which holds because $T_j$ and $T_\ell$ co-penetrate.

    Any lift of $T_\ell$ separating $\tilde{\gamma}_j$ from $g^{\,n}\tilde{\gamma}_j''$ is penetrated by $\tilde{\delta}_j^{\,n}(\tilde{\gamma})$.
    Hence, for each $T_\ell \in \Gamma \setminus \operatorname{st}(T_j)$, both $\tilde{\delta}_j^{\,n}(\tilde{\gamma})$ and $\tilde{\gamma}_j'$ penetrate at least $2\,D(\Gamma) + 3$ lifts of $T_\ell$.
    Therefore $\delta_j^{\,n}(\gamma) \in \mathcal{FT}\!_{2\,D(\Gamma)+3}(T_j, \Gamma)$.
\end{proof}

\begin{proposition} \label{prop:commuting}
    Let $T_i, T_{j_1}, \dots, T_{j_m}$ be pairwise distinct and pairwise disjoint members of $V$, that is, $i(T_a, T_b) = 0$ for all $a, b \in \{i, j_1, \dots, j_m\}$.
    Then
    \[
        \delta_{j_m}^{n_m} \dots \delta_{j_1}^{n_1}\bigl(\mathcal{FT}\!_{2\,D(\Gamma)+3}(T_i, \Gamma)\bigr) \subseteq \mathcal{FT}\!_3(T_i, \Gamma)
    \]
    for all $n_1, \dots, n_m \in \ZZ$.
\end{proposition}

\begin{proof}
    Choose $\gamma \in \mathcal{FT}\!_{2\,D(\Gamma)+3}(T_i, \Gamma)$ and let $\tilde\gamma$, $\tilde\gamma_i$ be lifts as in Lemma \ref{lem:fellow-travel}.
    Fix $T_j \in \Gamma \setminus \operatorname{st}(T_i)$ and let $\tilde T_j, \tilde T_j', \tilde T_j''$ be the three lifts the lemma provides.
    Since $T_{j_1}, \dots, T_{j_m} \in \operatorname{lk}(T_i)$, clause (2) of the lemma says that each of these three cylinders is disjoint from every lift of every $T_{j_a}$ penetrated by $\tilde\gamma$.
    We show that $\tilde T_j$ is penetrated by the image of $\tilde\gamma$; the arguments for $\tilde T_j'$ and $\tilde T_j''$ are identical.

    For each $a$, let $\tilde\delta_{j_a}$ be the lift of $\delta_{j_a}$ fixing $\tilde\gamma_i$ pointwise, which exists because $i(T_i, T_{j_a}) = 0$ makes $\tilde\gamma_i$ disjoint from every lift of $T_{j_a}$.
    Put $\tilde\delta := \tilde\delta_{j_m}^{n_m} \dots \tilde\delta_{j_1}^{n_1}$.
    Each $\tilde\delta_{j_a}$ is supported in a union of regular neighborhoods of the lifts of $T_{j_a}$ and preserves each such lift; hence $\tilde\delta$ preserves every lift of every $T_{j_a}$, and fixes pointwise the complement of the union of those neighborhoods.

    Let $\mathcal{O}$ be the set of lifts of $T_{j_1}, \dots, T_{j_m}$ penetrated by $\tilde\gamma$.
    By Lemma \ref{lem:geodesic_in_dual_tree} \eqref{item:geodesic:disjoint}, $\tilde\gamma$ is disjoint from every lift of every $T_{j_a}$ outside $\mathcal{O}$, and by \eqref{item:geodesic:once} it meets each member of $\mathcal{O}$ in exactly one point; order the members of $\mathcal{O}$ by the order in which $\tilde\gamma$ meets them.
    If $\mathcal{O} = \emptyset$, then $\tilde\gamma$ is disjoint from the support of $\tilde\delta$, so $\tilde\delta$ fixes $\tilde\gamma$ pointwise and the conclusion is immediate.
    Assume therefore $\mathcal{O} \ne \emptyset$.

    We claim that $\mathcal{O}$ is bi-infinite in this order.
    Fix $\tilde S \in \mathcal{O}$, say a lift of $T_{j_a}$, and let $h$ generate $\operatorname{Stab}(\tilde\gamma)$.
    Since $h$ preserves $\tilde\gamma$, it carries penetrated lifts to penetrated lifts, so $h^p \tilde S \in \mathcal{O}$ for every $p \in \ZZ$.
    These are pairwise distinct: if $h^p \tilde S = h^{p'} \tilde S$ with $p \ne p'$, then $h^{p - p'}$ is a nontrivial element of $\operatorname{Stab}(\tilde S)$, so $\gamma$ and $\gamma_{j_a}$ are commensurable and Lemma \ref{lem:penetration_noncomm} gives $i(\gamma, T_{j_a}) = 0$, contradicting $\tilde S \in \mathcal{O}$.
    As $h$ acts on $\tilde\gamma$ by translation, the points $h^p(\tilde\gamma \cap \tilde S)$ leave every compact subset of $\tilde\gamma$ in both directions, which proves the claim.

    Since $\tilde T_j$ is connected and, by clause (2) of Lemma \ref{lem:fellow-travel}, disjoint from every member of $\mathcal{O}$, it lies in a single component of $\widetilde\MM \setminus \bigsqcup_{\tilde S \in \mathcal{O}} \tilde S$; as $\tilde\gamma$ meets $\tilde T_j$, this is a component through which $\tilde\gamma$ passes, and by the previous paragraph it is bounded by two members $\tilde S, \tilde S'$ of $\mathcal{O}$, consecutive in the order along $\tilde\gamma$.
    Let $\alpha$ be the arc of $\tilde\gamma$ in this component, so that $\alpha$ runs from $\tilde S$ to $\tilde S'$ and meets $\tilde T_j$ in exactly one point.
    Writing $P$ and $Q$ for the two sides of $\tilde T_j$, each of $\tilde S$ and $\tilde S'$ is connected and disjoint from $\tilde T_j$, hence contained in $P$ or in $Q$; since the endpoints of $\alpha$ lie on opposite sides of $\tilde T_j$, we may label so that
    \[
        \tilde S \subset P, \qquad \tilde S' \subset Q .
    \]

    Now let $c$ be an arbitrary representative of $\tilde\delta(\tilde\gamma)$.
    Each $\tilde\delta_{j_b}$ preserves every lift of $T_{j_b}$, so $\tilde\delta$ preserves every member of $\mathcal{O}$; hence $\tilde\delta(\tilde\gamma)$ penetrates $\tilde S$ and $\tilde S'$, and by the definition of penetration $c$ meets both.
    Therefore $c$ has a point in $P$ and a point in $Q$, and being connected it meets $\tilde T_j$.
    As $c$ was arbitrary, $\tilde\delta(\tilde\gamma)$ penetrates $\tilde T_j$.
    The same argument applies to $\tilde T_j'$ and $\tilde T_j''$.
    Finally $\tilde\delta$ fixes $\tilde\gamma_i$ pointwise, so these three lifts remain penetrated by $\tilde\gamma_i$ as well.
    Therefore $\delta_{j_m}^{n_m} \dots \delta_{j_1}^{n_1}(\gamma) \in \mathcal{FT}\!_3(T_i, \Gamma)$.
\end{proof}

\begin{remark}
    The requirement that the cores of a compatible collection be pairwise distinct and indivisible is essential for Proposition \ref{prop:commuting}.
    By Lemma \ref{lem:compatible_noncomm} it guarantees that the cores are pairwise non-commensurable, which is what Proposition \ref{prop:two_cylinders}, and hence the bound $D(\Gamma)$, requires.
    Without it the composition $\delta_{j_m}^{n_m} \dots \delta_{j_1}^{n_1}$ could agree with a power of $\delta_i$ on the relevant lifts and send an element of $\mathcal{FT}\!_{2\,D(\Gamma)+3}(T_i, \Gamma)$ outside $\mathcal{FT}\!_3(T_i, \Gamma)$.
\end{remark}


\subsection{Right-angled Artin group generated by powers of Dehn twists}

In this section, we look into a specific collection of essential tori whose sufficiently large powers of Dehn twists generate a right-angled Artin group.
Note that there exist finitely many Dehn twists that generate $\Out(\FF)$, which is far more complicated than a right-angled Artin group.
This section is inspired by Koberda \cite{MR3000498} and Seo \cite{MR4299673}.

Let us recall the relevant definitions.
A collection $V = \{T_1, \dots, T_k\}$ of homotopy classes of essential tori is \emph{compatible} if their cores are pairwise distinct and indivisible, and whenever $i(T_i, T_j) > 0$ the classes $T_i$ and $T_j$ co-penetrate.
The coincidence graph $\Gamma$ of $V$ is the simple graph with vertices $T_1, \dots, T_k$ in which $T_i$ and $T_j$ are joined by an edge if and only if $i(T_i, T_j) = 0$.
Recall also that
\[
    M(\Gamma) = \max_{i, j} i(\gamma_i, T_j),
    \qquad
    D(\Gamma) = 7 \, I(\Gamma)^2 .
\]

\begin{lemma} \label{lem:core_not_fellow}
    Let $T_0, T \in V$ and let $\gamma_0$ be the core of $T_0$.
    Then $\gamma_0 \notin \mathcal{FT}\!_N(T, \Gamma)$ for every $N \ge 1$.
\end{lemma}

\begin{proof}
    Suppose $\gamma_0 \in \mathcal{FT}\!_N(T, \Gamma)$.
    Condition (1) requires $i(\gamma_0, T) > 0$; since $i(\gamma_0, T_0) = 0$, this forces $T \ne T_0$ and $i(T_0, T) > 0$, so $T_0 \in \Gamma \setminus \operatorname{st}(T)$.
    Condition (2), applied with $T_j = T_0$, then requires some lift of $\gamma_0$ to penetrate at least $N \ge 1$ lifts of $T_0$, contradicting $i(\gamma_0, T_0) = 0$.
\end{proof}

\begin{theorem} \label{thm:main}
    Let $V = \{ T_1, \dots, T_k \}$ be a compatible collection of essential tori with cores $\gamma_1, \dots, \gamma_k$, and let $\Gamma$ be the coincidence graph of $V$.
    For each $i$, let $\delta_i = \delta_{\vec{T}_i}$ be a Dehn twist along $T_i$.
    Then for all
    \[
        \lvert n_1 \rvert, \dots, \lvert n_k \rvert \;\ge\; M(\Gamma) + 2\, D(\Gamma) + 8,
    \]
    the homomorphism
    \[
        \Phi \colon A(\Gamma) \longrightarrow \Out(\FF), \qquad T_i \longmapsto \delta_i^{\,n_i}
    \]
    is injective.
\end{theorem}

\begin{proof}
    Call a simple graph \emph{centerless} if it admits no decomposition as a join $K * \Lambda$ with $K$ a nonempty complete subgraph.
    Equivalently, $\Gamma$ is centerless if and only if $A(\Gamma)$ has trivial center, and if and only if no vertex of $\Gamma$ is adjacent to all other vertices.
    We first prove the theorem when $\Gamma$ is centerless, and then deduce the general case.
    
    For a nontrivial $x \in A(\Gamma)$, the \emph{support} $\operatorname{supp}(x) \subseteq V(\Gamma)$ is the set of vertices occurring in some reduced word representing $x$; this is well defined, since any two reduced words representing $x$ have the same support.
    An element of $A(\Gamma)$ is \emph{central} if its support spans a complete subgraph of $\Gamma$; such an element generates, together with the vertices of its support, a free abelian subgroup.
    Every element of $A(\Gamma)$ can be written as a product of central elements, and such an expression is called a \emph{central form}.
    A central form $x = x_r \dots x_1$ is in \emph{left greedy normal form} if for each $s < r$ and each $v \in \operatorname{supp}(x_s)$ there is $u \in \operatorname{supp}(x_{s+1})$ with $[u, v] \ne 1$ in $A(\Gamma)$, that is, with $u$ and $v$ non-adjacent in $\Gamma$.
    Every nontrivial element of $A(\Gamma)$ admits a left greedy normal form; see Koberda \cite{MR3000498}.
    
    Assume $\Gamma$ is centerless and let $x \in A(\Gamma)$ be nontrivial, written in left greedy normal form $x = x_r \dots x_1$.

    Since $\Gamma$ is centerless, $\operatorname{st}(\operatorname{supp}(x_1)) \ne V(\Gamma)$, so we may choose a vertex
    \[
        T_0 \in \Gamma \setminus \operatorname{st}(\operatorname{supp}(x_1))
    \]
    and a vertex $T_{i_1} \in \operatorname{supp}(x_1)$ with $i(T_{i_1}, T_0) > 0$.
    Let $\gamma_0$ be the core of $T_0$.
    Since $V$ is compatible and $i(T_{i_1}, T_0) > 0$, the classes $T_{i_1}$ and $T_0$ co-penetrate.

    Because $x_1$ is central, its support spans a complete subgraph, so we may write
    \[
        x_1 = T_{i_1}^{\,m_1} y_1,
        \qquad m_1 \ne 0,
        \qquad \operatorname{supp}(y_1) \subseteq \operatorname{lk}(T_{i_1}) \cap \operatorname{supp}(x_1),
    \]
    and the vertices of $\operatorname{supp}(y_1)$ are pairwise adjacent and adjacent to $T_{i_1}$.
    Now $\Phi(T_{i_1}^{\,m_1}) = \delta_{i_1}^{\,m_1 n_{i_1}}$ and
    \[
        \lvert m_1 n_{i_1} \rvert \;\ge\; \lvert n_{i_1} \rvert \;\ge\; M(\Gamma) + 2\,D(\Gamma) + 8 ,
    \]
    so Proposition \ref{prop:co-penetrating}, applied to the co-penetrating pair $T_0$ and $T_{i_1}$ and to $\gamma_0$, gives
    \[
        \Phi(T_{i_1}^{\,m_1})(\gamma_0) \in \mathcal{FT}\!_{2\,D(\Gamma)+3}(T_{i_1}, \Gamma).
    \]
    The element $\Phi(y_1)$ is a product of powers of the twists along the vertices of $\operatorname{supp}(y_1)$, which are pairwise distinct and pairwise disjoint members of $V$, all lying in $\operatorname{lk}(T_{i_1})$.
    Hence Proposition \ref{prop:commuting} applies and yields
    \[
        \Phi(x_1)(\gamma_0) = \Phi(y_1)\bigl(\Phi(T_{i_1}^{\,m_1})(\gamma_0)\bigr) \in \mathcal{FT}\!_3(T_{i_1}, \Gamma).
    \]

    We now iterate.
    Suppose that for some $s$ with $2 \le s \le r$ we have chosen $T_{i_{s-1}} \in \operatorname{supp}(x_{s-1})$ with
    \[
        \Phi(x_{s-1} \dots x_1)(\gamma_0) \in \mathcal{FT}\!_3(T_{i_{s-1}}, \Gamma).
    \]
    By the left greedy condition there is $T_{i_s} \in \operatorname{supp}(x_s)$ with $T_{i_s}$ and $T_{i_{s-1}}$ non-adjacent, that is, $i(T_{i_s}, T_{i_{s-1}}) > 0$; by compatibility they co-penetrate.
    As above, write $x_s = T_{i_s}^{\,m_s} y_s$ with $m_s \ne 0$ and $\operatorname{supp}(y_s) \subseteq \operatorname{lk}(T_{i_s})$.
    Proposition \ref{prop:co-penetrating} applied to the co-penetrating pair $T_{i_{s-1}}$ and $T_{i_s}$ gives
    \[
        \Phi(T_{i_s}^{\,m_s})\bigl(\mathcal{FT}\!_3(T_{i_{s-1}}, \Gamma)\bigr) \subseteq \mathcal{FT}\!_{2\,D(\Gamma)+3}(T_{i_s}, \Gamma),
    \]
    and Proposition \ref{prop:commuting} applied to $\Phi(y_s)$ gives
    \[
        \Phi(x_s \dots x_1)(\gamma_0)
        = \Phi(y_s)\Bigl(\Phi(T_{i_s}^{\,m_s})\bigl(\Phi(x_{s-1}\dots x_1)(\gamma_0)\bigr)\Bigr)
        \in \mathcal{FT}\!_3(T_{i_s}, \Gamma).
    \]
    By induction,
    \[
        \Phi(x)(\gamma_0) \in \mathcal{FT}\!_3(T_{i_r}, \Gamma).
    \]

    On the other hand $\gamma_0 \notin \mathcal{FT}\!_3(T_{i_r}, \Gamma)$ by Lemma \ref{lem:core_not_fellow}.
    Therefore $\Phi(x)(\gamma_0) \ne \gamma_0$ and $\Phi(x) \ne 1$.
    Since $x$ was an arbitrary nontrivial element, $\Phi$ is injective when $\Gamma$ is centerless.
    
    Now let $\Gamma$ be arbitrary and write $\Gamma = K * \Gamma'$, where $K$ is the (possibly empty) complete subgraph spanned by the vertices adjacent to all other vertices, and $\Gamma'$ is the induced subgraph on the remaining vertices; then $\Gamma'$ is centerless.
    If $\Gamma' = \emptyset$, then $\Gamma$ is complete and $A(\Gamma) = \ZZ^{K}$, so the conclusion follows from the last paragraph of this proof alone; we therefore assume $\Gamma' \ne \emptyset$.
    Correspondingly
    \[
        A(\Gamma) \cong A(\Gamma') \times \ZZ^{K},
        \qquad
        \mathcal{Z}(A(\Gamma)) \cong \ZZ^{K} .
    \]
    If $K = \emptyset$ we are in the centerless case, so assume $K \ne \emptyset$.

    We first show $\ker(\Phi) \le \mathcal{Z}(A(\Gamma))$.
    Let $x \in \ker(\Phi)$ and write $x = zy$ with $z \in \ZZ^{K}$ and $y \in A(\Gamma')$.
    Then $\Phi(y) = \Phi(z)^{-1}$.
    Since $z$ is central in $A(\Gamma)$, the element $\Phi(z)$ commutes with every element of $\Phi(A(\Gamma'))$, hence so does $\Phi(y)$; that is, $\Phi(y)$ lies in the center of $\Phi(A(\Gamma'))$.
    The subcollection $V' \subseteq V$ corresponding to the vertices of $\Gamma'$ is again compatible, its coincidence graph is $\Gamma'$, and $M(\Gamma') \le M(\Gamma)$ and $D(\Gamma') \le D(\Gamma)$, so the given exponents satisfy the threshold for $\Gamma'$.
    By the centerless case applied to $V'$, the restriction of $\Phi$ to $A(\Gamma')$ is injective, so $\Phi(A(\Gamma')) \cong A(\Gamma')$, which has trivial center because $\Gamma'$ is centerless.
    Hence $\Phi(y) = 1$, so $y = 1$ and $x = z \in \mathcal{Z}(A(\Gamma))$.

    It remains to show that $\Phi$ is injective on $\mathcal{Z}(A(\Gamma)) \cong \ZZ^{K}$.
    The vertices of $K$ are pairwise adjacent in $\Gamma$, so the corresponding members of $V$ are pairwise disjoint, and they are pairwise distinct because the cores of $V$ are.
    By Corollary \ref{cor:disjoint_twists_free_abelian}, the twists $\delta_v$ with $v \in K$ generate a free abelian subgroup of $\Out(\FF)$ of rank $\lvert K \rvert$.
    Hence a relation $\prod_{v \in K} \delta_v^{\,n_v m_v} = 1$ forces $n_v m_v = 0$, and therefore $m_v = 0$, for every $v \in K$.
    Thus $\Phi$ is injective on the center, and combined with the previous paragraph, $\ker(\Phi) = 1$.
\end{proof}





\section{The space of primitive curves}





% \subsection{Iteration of a Dehn twist}


An essential simple closed curve is called \emph{primitive} if it represents a \(1\)-free factor of \(\mathbb{F}\).  
Let \(\mathcal{C}\) denote the set of all homotopy classes of primitive curves.  
Let \(\mathbb{R}_+\) denote the set of nonnegative real numbers, i.e., \([0,\infty)\).  
Let $\mathcal{S}$ denote the set of all isotopy classes of essential spheres.
Define a map  
\[
i_* : \mathcal{C} \to \mathbb{R}_+^{\mathcal{S}}, \qquad
i_*(\gamma)(S) = i(\gamma, S)
\]
for all \(\gamma \in \mathcal{C}\) and \(S \in \mathcal{S}\).  
Let \(P(\mathbb{R}_+^{\mathcal{S}})\) denote the projectivization of \(\mathbb{R}_+^{\mathcal{S}}\), and let  
\[
p : \mathbb{R}_+^{\mathcal{S}} \setminus \{0\} \to P(\mathbb{R}_+^{\mathcal{S}})
\]
be the natural projection.  
Since the image of \(i_*\) contains no zero vector, the composition \(p \circ i_*\) is well-defined.

% An essential simple closed curve is called \emph{primitive} if it represents a $1$-free factor of $\FF$.
% Let us write $\mathcal{C}$ for the set of all homotopy classes of primitive curves.
% Let $\mathbb{R}_+$ denote the space of nonnegative real numbers, that is, $[0, \infty)$.
% We define a map $i_*: \mathcal{C} \to \mathbb{R}_+^{\mathcal{S}}$ by \[i_*(\gamma)(S) = i(\gamma, S)\] for all $\gamma \in \mathcal{C}$ and $S \in \mathcal{S}$.
% Let $P(\mathbb{R}_+^\mathcal{S})$ denote the projective space of $\mathbb{R}_+^\mathcal{S}$ with the projection $p: \mathbb{R}_+^\mathcal{S} \setminus \{0\} \to P(\mathbb{R}_+^\mathcal{S})$.
% The image of $i_*$ does not contain a zero map, so $p \circ i_*$ is well-defined.

\begin{proposition} \label{prop:conti}
    The composition $p \circ i_* : \mathcal{C} \to P(\mathbb{R}_+^\mathcal{S})$ is injective.
\end{proposition}

\begin{proof}
    Assume that \(\gamma, \gamma' \in \mathcal{C}\) satisfy  
    \[
    (p \circ i_*)(\gamma) = (p \circ i_*)(\gamma').
    \]
    Then there exists \(\lambda > 0\) such that  
    \[
    i_*(\gamma) = \lambda \, i_*(\gamma').
    \]
    
    Since \(\gamma\) corresponds to a free factor, there exists a separating sphere \(s\) disjoint from a representative \(c \in \gamma\) such that the component \(C\) of \(\MM \setminus s\) containing \(c\) is homeomorphic to the once--punctured doubled solid torus $(\mathbb{S}^1 \times \mathbb{S}^2)\setminus \{x\}$ for some point \(x \in \mathbb{S}^1 \times \mathbb{S}^2\).
    This region \(C\) contains a unique isotopy class \(S'\) of essential spheres not isotopic to \(s\).  
    Choose a representative \(s' \in S'\) such that \(c\) intersects \(s'\) in exactly one point.  
    Thus,
    \[
    i(\gamma,[s]) = 0, \qquad i(\gamma, S') = \lvert c \cap s' \rvert = 1.
    \]
    
    By proportionality, it follows that  
    \[
    i(\gamma', [s]) = 0, \qquad i(\gamma', S') = \frac{1}{\lambda} > 0.
    \]
    So \(\gamma'\) has a representative \(c'\) disjoint from \(s\).  
    Because \(s\) is separating, the curve \(c'\) lies entirely inside \(C\).  
    Algebraically, the free factor determined by \(c'\) is contained in the free factor determined by \(C\).  
    Since \(c'\) is primitive, these free factors must be equal.  
    Consequently, \(c'\) is homotopic to \(c\), and hence \(\gamma = \gamma'\).
\end{proof}

By Proposition~\ref{prop:conti}, we may regard \(\mathcal{C}\) as a subspace of \(P(\mathbb{R}_+^{\mathcal{S}})\).

\begin{proposition} \label{prop:C_compact}
    The closure of $(p \circ i_*)(\mathcal{C})$ is compact.
\end{proposition}

Before proving Proposition~\ref{prop:C_compact}, we require the following technical lemma. This result is analogous to \cite[Lemma~4.2]{MR3053012}.

\begin{lemma} \label{lem:bounded}
    Let $\Sigma$ be a maximal sphere system of $\MM_g$.
    Then we have \[i(\gamma, S) \leq (2g-2) \, i(\gamma, \Sigma) \, i(S, \Sigma) \] for all $\gamma \in \mathcal{C}$ and $S \in \mathcal{S}$.
\end{lemma}

\begin{proof}
    Choose representatives \(c \in \gamma\), \(s \in S\), and \(\sigma \in \Sigma\) that intersect pairwise minimally.
    Cutting \(\MM_g\) along \(\sigma\) yields \(2g-2\) thrice-punctured spheres.
    For each such component \(N\), the intersection \(s \cap N\) is a disjoint union of annuli, pairs of pants, and disks, with at most \(i(S,\Sigma)\) components.
    Similarly, \(c \cap N\) is a disjoint union of at most \(i(\gamma,\Sigma)\) arcs.
    Each of these arcs intersects \(s\) in at most \(i(S,\Sigma)\) points.
    So the number of intersection points of \(c\) and \(s\) inside \(N\) is at most \(i(\gamma,\Sigma)\,i(S,\Sigma)\).
    Summing over all thrice-punctured components gives the desired inequality.
\end{proof}

The following proof is analogous to \cite[Proposition 4.1]{MR3053012}.

\begin{proof}[Proof of Proposition \ref{prop:C_compact}]
    Fix an isotopy class of a maximal sphere system \(\Sigma\).  
    By Lemma~\ref{lem:bounded}, the set $\{ i_*(\gamma)/i(\gamma, \Sigma) \mid \gamma \in \mathcal{C} \}$ is a subset of the bounded set
    \[
        \{ f \in \mathbb{R}_{+}^{\mathcal{S}} \mid f(S) \le (2g - 2)\, i(S, \Sigma) \text{ for all } S \in \mathcal{S} \}.
    \]
    By Tychonoff's theorem, the closure of 
    \(\{ i_*(\gamma)/i(\gamma, \Sigma) \mid \gamma \in \mathcal{C} \}\) 
    is compact.  
    Consequently, the closure of \((p \circ i_*)(\mathcal{C})\) is compact.
\end{proof}




The following is motivated by \cite[Proposition A.1]{MR3053012}.

\begin{lemma} \label{lem:disjoint_twist_growth}
    Let $T_1, \dots, T_r$ be pairwise disjoint homotopy classes of essential tori with cores $\gamma_1, \dots, \gamma_r$, and fix an orientation of each.
    Let $\alpha$ be a homotopy class of an essential simple closed curve and let $S$ be an isotopy class of an essential sphere.
    Then for all $N_1, \dots, N_r \in \ZZ$,
    \[
        \left\lvert\;
        i\bigl(\delta_{\vec T_1}^{N_1} \dots \delta_{\vec T_r}^{N_r}(\alpha),\, S\bigr)
        \;-\;
        \sum_{v=1}^{r} \lvert N_v \rvert \, i(\alpha, T_v) \, i(\gamma_v, S)
        \;\right\rvert
        \;\le\; i(\alpha, S).
    \]
\end{lemma}

\begin{proof}
    Choose representatives $a \in \alpha$, $t_v \in T_v$ and $s \in S$ in simultaneous minimal position, which is possible by Theorem \ref{thm:minimal-position}; since the $T_v$ are pairwise disjoint, the $t_v$ are pairwise disjoint.
    Choose pairwise disjoint regular neighborhoods $N(t_v)$, each disjoint from $a \cap s$, and let $a'$ be the representative of $\delta_{\vec T_1}^{N_1} \dots \delta_{\vec T_r}^{N_r}(\alpha)$ that agrees with $a$ outside $\bigcup_v N(t_v)$ and is obtained in each $N(t_v)$ by applying $\delta_{\vec t_v}^{N_v}$, arranged so that every component of $a' \cap N(t_v)$ meets $s$ minimally.

    Write $\mathcal{A}$ for the set of components of $a' \cap \bigcup_v N(t_v)$ and $\mathcal{B}$ for the set of components of $a' \setminus \bigcup_v N(t_v)$; the latter are subarcs of $a$.
    For each $v$ the set $a' \cap N(t_v)$ has exactly $i(\alpha, T_v)$ components, each meeting $s$ in $\lvert N_v \rvert \, i(\gamma_v, S)$ points.
    Since the neighborhoods are disjoint,
    \[
        \lvert a' \cap s \rvert
        = \lvert a \cap s \rvert + \sum_{v=1}^{r} \lvert N_v \rvert \, i(\alpha, T_v) \, i(\gamma_v, S),
    \]
    which gives the upper bound $i(\delta_{\vec T_1}^{N_1} \dots \delta_{\vec T_r}^{N_r}(\alpha), S) \le \lvert a' \cap s \rvert$.

    For the lower bound we bound the number of bigon reductions needed to put $a'$ and $s$ in minimal position.
    Suppose $a'$ and $s$ form a bigon, and choose one that is innermost, in the sense that the disk it bounds contains no other bigon between $a'$ and $s$.
    Let $a'' \subset a'$ and $b'' \subset s$ be its two sides.
    Since $a$ and $s$ are in minimal position, they form no bigon, so $a''$ is not contained in a single element of $\mathcal{B}$; hence $a''$ meets $\mathcal{A}$.
    On the other hand, no element of $\mathcal{A}$ is contained in $a''$: each such component meets $s$ minimally, so a subarc of it together with a subarc of $s$ bounding a disk would contradict innermostness.
    Therefore $a''$ meets at least one and at most two elements of $\mathcal{A}$, and in each it is a proper subarc; as the $N(t_v)$ are disjoint and $a''$ is connected, $a''$ decomposes as a concatenation
    \[
        a'' = a_1 \cdot c \cdot a_2 ,
    \]
    where $a_1$ and $a_2$ are (possibly empty) proper subarcs of elements of $\mathcal{A}$ and $c$ is a subarc of $a$, at least one of $a_1, a_2$ being nonempty.
    In particular an endpoint of $a''$ lies in the interior of an element of $\mathcal{A}$, so the intersection point of $a''$ with $s$ nearest that endpoint lies in $\mathcal{A} \cap s$, while the reduction across the bigon removes it together with one point of $c \cap s \subset a \cap s$.

    Thus each innermost bigon reduction destroys at least one point of $a \cap s$, and the points so destroyed are distinct for distinct reductions.
    Hence at most $\lvert a \cap s \rvert$ reductions occur, and since each removes two intersection points,
    \[
        i\bigl(\delta_{\vec T_1}^{N_1} \dots \delta_{\vec T_r}^{N_r}(\alpha), S\bigr)
        \;\ge\; \lvert a' \cap s \rvert - 2 \lvert a \cap s \rvert
        \;=\; \sum_{v=1}^{r} \lvert N_v \rvert \, i(\alpha, T_v) \, i(\gamma_v, S) - i(\alpha, S) .
    \]
    Combining the two bounds gives the stated inequality.
\end{proof}

% \begin{lemma}
%     Let $T$ be a homotopy class of an essential torus with a core $\gamma$.
%     If $\alpha$ is a homotopy class of a primitive curve and $S$ is an isotopy class of an essential sphere, then the following holds for all $n \in \mathbb{Z}$:
%     \[\lvert i(\delta_{\vec{T}}^n(\alpha), S) - \lvert n \rvert \, i(\alpha, T) \, i(\gamma, S) \rvert \leq i(\alpha,S).\]
% \end{lemma}

% \begin{proof}
%     Let $a \in \alpha$, $t \in T$ and $s \in S$ be representatives in minimal position.
%     Let $N(t)$ be a regular neighborhood of $t$ which is disjoint from $a \cap s$.
%     Let $a_n$ be a representative of $\delta_{\vec{T}}^n(\alpha)$ that follows $a$ on the outside of $N(t)$ such that each component of $a_n \cap N(t)$ minimally intersect with $s$.
    
%     Then the number of components of $a_n \cap N(t)$ is exactly $i(\alpha, T)$.
%     Furthermore, each component intersects with $s$ in $(\lvert n \rvert \, i(\gamma, S))$-times.
%     So we obtain the following.
%     \begin{align*}
%         i(\delta_{\vec{T}}^n(\alpha), S) &\leq \lvert a_n \cap s \rvert = \lvert ( a_n - N(t) ) \cap s \rvert + \lvert ( a_n \cap N(t) ) \cap s \rvert \\
%         &= \lvert a \cap s \rvert + \lvert i(\alpha, T) \rvert \, \lvert n \rvert \, i(\gamma, S) \\
%         & = i(\alpha, S) + \lvert n \rvert \, i(\alpha, T) \, i(\gamma, S)
%     \end{align*}

%     Suppose $a_n$ and $s$ form a bigon.
%     Then there exists an arc $a' \subset a_n$ and $b' \subset s$ such that $a' \cup b'$ bounds a disk.
%     Since $a$ and $s$ do not form a bigon, $a'$ is not contained in $a$, that is, $a'$ and $a_n \cap N(t)$ intersect.
%     Any component of $a_n \cap N(t)$ cannot belong to $a'$ because this minimally intersects with $s$.
%     So we can decompose $a'$ into two arcs $a''$ and $a'''$ such that $a'' \in a$ and $a''' \in a_n \cap N(t)$.

%     By the above, reducing a bigon by homotoping $a_n$ induces the decrease of the number of $(a \cap a_n) \cap s$.
%     So there are at most $\lvert a \cap s \rvert$ reductions of bigons to make $a_n$ and $s$ become in minimal position.
%     Since each bigon makes two intersection points, we get
%     \begin{align*}
%         i(\delta_{\vec{T}}^n(\alpha), S) &\geq \lvert a_n \cap s \rvert - 2 \, \lvert a \cap s \rvert \\
%         & = \lvert n \rvert \, i(\alpha, T) \, i(\gamma, S) - i(\alpha, S)
%     \end{align*}
%     Therefore, the inequality of the statement holds.
% \end{proof}

The above lemma immediately induces the following theorem.

\begin{theorem} \label{thm:multi-converge}
    Let $T_1, \dots, T_r$ be pairwise disjoint homotopy classes of essential tori
    with cores $\gamma_1, \dots, \gamma_r$, and fix an orientation of each, so that
    the Dehn twists $\delta_{\vec T_1}, \dots, \delta_{\vec T_r}$ pairwise commute
    by Theorem \ref{thm:commuting}.
    For $\vec N = (N_1, \dots, N_r) \in \ZZ^r$, write
    \[
        \Delta_{\vec N} := \delta_{\vec T_1}^{N_1} \dots \delta_{\vec T_r}^{N_r} \in \Out(\FF).
    \]
    Let $\alpha \in \mathcal{C}$, and let $(\vec N^{(n)})_{n \ge 1}$ be a sequence in
    $\ZZ^r$ such that
    \[
        w_n := \sum_{v=1}^{r} \lvert N_v^{(n)} \rvert \, i(\alpha, T_v)
        \longrightarrow \infty
        \qquad\text{and}\qquad
        \frac{\lvert N_v^{(n)} \rvert \, i(\alpha, T_v)}{w_n} \longrightarrow c_v
        \quad (1 \le v \le r),
    \]
    where necessarily $c_v \ge 0$ and $\sum_v c_v = 1$.
    Then
    \[
        (p \circ i_*)\bigl(\Delta_{\vec N^{(n)}}(\alpha)\bigr)
        \longrightarrow
        \left[\, \sum_{v=1}^{r} c_v \, i_*(\gamma_v) \,\right]
        \qquad \text{in } P(\mathbb{R}_+^{\mathcal{S}}).
    \]
\end{theorem}

\begin{proof}
    Outer automorphisms carry free factors to free factors, so
    $\Delta_{\vec N^{(n)}}(\alpha) \in \mathcal{C}$ and the statement makes sense.
    Put $\eta := \sum_v c_v \, i_*(\gamma_v) \in \mathbb{R}_+^{\mathcal{S}}$.
    By Lemma \ref{lem:disjoint_twist_growth}, for every $S \in \mathcal{S}$ and every $n$,
    \[
        \left\lvert\; i\bigl(\Delta_{\vec N^{(n)}}(\alpha), S\bigr)
        - \sum_{v=1}^{r} \lvert N_v^{(n)} \rvert \, i(\alpha, T_v)\, i(\gamma_v, S)
        \;\right\rvert \;\le\; i(\alpha, S).
    \]
    Dividing by $w_n$ and letting $n \to \infty$, the error term tends to $0$ and the
    middle sum tends to $\eta(S)$; hence
    $\tfrac{1}{w_n}\, i_*\bigl(\Delta_{\vec N^{(n)}}(\alpha)\bigr) \to \eta$
    in the product topology on $\mathbb{R}_+^{\mathcal{S}}$.
    The limit is nonzero: choose $v_0$ with $c_{v_0} > 0$, write $\gamma_{v_0}$
    cyclically reduced with respect to a fixed basis of $\FF$, and let $S_0$ be the
    nonseparating essential sphere dual to a basis element occurring in it, as in the
    proof of Corollary \ref{cor:disjoint_twists_free_abelian}; then
    $\eta(S_0) \ge c_{v_0}\, i(\gamma_{v_0}, S_0) > 0$.
    Since $p$ is continuous and invariant under scaling by positive reals,
    $(p \circ i_*)\bigl(\Delta_{\vec N^{(n)}}(\alpha)\bigr) \to p(\eta)$.
\end{proof}

\begin{corollary}\label{cor:multitwist-parabolic}
    Fix $\vec m \in \ZZ^r$ and put $\Delta := \Delta_{\vec m}$.
    For every $\alpha \in \mathcal{C}$, exactly one of the following holds.
    \begin{enumerate}
        \item \label{enum:multi-para:1} $i(\alpha, T_v) = 0$ for every $v$ with $m_v \ne 0$; in this case
              $\Delta^n(\alpha) = \alpha$ for all $n \in \ZZ$.
        \item \label{enum:multi-para:2} Otherwise,
        \[
            (p \circ i_*)\bigl(\Delta^{n}(\alpha)\bigr)
            \longrightarrow
            \left[\, \sum_{v=1}^{r} \lvert m_v \rvert \, i(\alpha, T_v)\, i_*(\gamma_v) \,\right]
            \qquad (\lvert n \rvert \to \infty).
        \]
    \end{enumerate}
    Moreover, every point $\bigl[\sum_v s_v\, i_*(\gamma_v)\bigr]$ with
    $s_1, \dots, s_r \ge 0$ not all zero is fixed by $\Delta$.
\end{corollary}

\begin{proof}
    In case \eqref{enum:multi-para:1}, Theorem \ref{thm:minimal-position} provides representatives
    $a \in \alpha$ and $t_v \in T_v$ in simultaneous minimal position; then $a$ is
    disjoint from $t_v$ for every $v$ with $m_v \ne 0$, so a representative of
    $\Delta$ supported in disjoint regular neighborhoods of these $t_v$ fixes $a$,
    and $\Delta^n(\alpha) = \alpha$ for all $n$.
    
    In case \eqref{enum:multi-para:2}, apply Theorem \ref{thm:multi-converge} with $\vec N^{(n)} = n \vec m$:
    here $w_n = \lvert n \rvert \sum_v \lvert m_v \rvert\, i(\alpha, T_v) \to \infty$ and
    $c_v = \lvert m_v \rvert\, i(\alpha, T_v) \big/ \sum_w \lvert m_w \rvert\, i(\alpha, T_w)$.
    For the last claim, note first that $\Delta(\gamma_v) = \gamma_v$ for each $v$.
    Indeed, for $w \ne v$ the representative $t_w$ is disjoint from $N(t_v)$, so $\delta_{\vec t_w}$ fixes a representative of $\gamma_v$ lying on $t_v$; and the twist homeomorphism $\delta_{\vec t_v}$ preserves the curve $h^{-1}(\{1/2\} \times \{0\} \times \mathbb{S}^1)$ setwise, since in the coordinates of $N(t_v)$ it acts on this circle by a rotation.
    Hence every representative of the core of $t_v$ is fixed up to homotopy by each factor of $\Delta$.
    Since $i(\phi(\gamma), S) = i(\gamma, \phi^{-1}(S))$ for every $\phi \in \Out(\FF)$, the function $i_*(\gamma_v)$ is fixed by the coordinate action of $\Delta$, and therefore so is $\bigl[\sum_v s_v\, i_*(\gamma_v)\bigr]$.
\end{proof}

\begin{corollary}\label{cor:core_simplex}
    With notation as above, let $\alpha \in \mathcal{C}$ and put
    $P_\alpha := \{\, v \mid i(\alpha, T_v) > 0 \,\}$.
    Then for every family $(s_v)_{v \in P_\alpha}$ of nonnegative reals, not all zero,
    \[
        \left[\, \sum_{v \in P_\alpha} s_v \, i_*(\gamma_v) \,\right]
        \;\in\; \overline{(p \circ i_*)(\mathcal{C})} .
    \]
\end{corollary}

\begin{proof}
    For $v \in P_\alpha$ set $N_v^{(n)} := \lfloor n s_v / i(\alpha, T_v) \rfloor$, and
    set $N_v^{(n)} := 0$ for $v \notin P_\alpha$.
    Then $\lvert N_v^{(n)} \rvert\, i(\alpha, T_v) = n s_v + O(1)$, so
    $w_n = n \sum_v s_v + O(1) \to \infty$ and $c_v = s_v \big/ \sum_w s_w$.
    Theorem \ref{thm:multi-converge} exhibits
    $\bigl[\sum_{v \in P_\alpha} s_v\, i_*(\gamma_v)\bigr]$ as a limit of points of
    $(p \circ i_*)(\mathcal{C})$.
\end{proof}

% \begin{theorem} \label{thm:converge}
%     Let $T$ be a homotopy class of an essential torus with the core $\gamma$.
%     Let $\alpha$ be a homotopy class of a primitive curve.
%     Then, we have $i(\alpha, T) > 0$ if and only if $(p \circ i_*)(\delta_{\vec{T}}^n(\alpha))$ converges to $(p \circ i_*)(\gamma)$ as $n \to \infty$.
% \end{theorem}

% \begin{proof}
%     If $i(\alpha, T) = 0$, then we have $\delta_{\vec{T}}^n(\alpha) = \alpha$ for all $n \in \mathbb{Z}$.
%     Otherwise, we have \[\left\lvert \frac{1}{\lvert n \rvert \, i(\alpha, T)} \, i(\delta_{\vec{T}}^n(\alpha), S) - i(\gamma, S) \right\rvert \leq \frac{1}{\lvert n \rvert \, i(\alpha, T)} \, i(\alpha, S) \]
%     for each $S \in \mathcal{S}$, so $(p \circ i_*) (\delta_{\vec{T}}^n(\alpha))$ converges to $(p \circ i_*)(\gamma)$ as $n \to \infty$.
% \end{proof}

% By Theorem \ref{thm:converge}, a core of a homotopy class of an essential torus is contained in the closure of $(p \circ i_*)(\mathcal{C})$.
% Note a core of a homotopy class of an essential torus may not be a primitive curve.
% That is, $\overline{(p \circ i_*)(\mathcal{C})}$ contains homotopy classes of essential simple closed curves that is not primitive.





% \subsection{Finite dimension}

% So, what is $\overline{(p \circ i_*)(\mathcal{C})}$?
% What property holds in $\overline{(p \circ i_*)(\mathcal{C})}$?
% In this section, we will show the following.

% \begin{theorem} \label{thm:finite_dimension}
%     There is a collection of $6g-6$ essential spheres $\{S_1, \dots, S_{9g-9}\}$ such that the map \[\overline{(p \circ i_*)(\mathcal{C})} \to P(\mathbb{R}_+^{9g-9})\] defined by \[[f] \mapsto [f(S_1), \dots, f(S_{9g-9})]\] is well-defined, injective and continuous.
% \end{theorem}

% Choose a primitive curve $\gamma$.
% Let $\Sigma$ be a maximal sphere system of $\MM$ in which each sphere does not bound a once-punctured doubled torus, for example, a maximal nonesparating sphere system.
% For each $S \in \Sigma$, there are pairwise disinct essential spheres $S$, $S'$ and $S''$ such that $\Sigma$, $(\Sigma \setminus \{S\}) \cup \{S'\}$ and $(\Sigma \setminus \{S\}) \cup \{S''\}$ are all maximal sphere systems of $\MM$.
% Let $N(\Sigma)$ denote the collection of such $S, S', S''$ for all $S \in \Sigma$.

% First, let $\rho: \widetilde{\MM} \to \MM$ be a universal cover.
% Let $\widetilde{N(\Sigma)}$ be the set of essential spheres in $\widetilde{\MM}$ whose projective images are contained in $N(\Sigma)$.
% Note that $\rho^{-1}(\gamma)$ is a collection of pairwise disjoint bi-infinite lines.
% For each $S \in N(\Sigma)$ and its lift $\tilde{S}$, we have $\lvert \rho^{-1}(\gamma) \cap \tilde{S} \rvert = i(\gamma, S)$.
% Let $\tilde{S}_0$ be an essential sphere in $\widetilde{\MM}$ which intersects exactly two spheres $\tilde{S}_6$ and $\tilde{S}_7$ for some $\tilde{S}_6, \tilde{S}_7 \in \widetilde{\Sigma}$.
% Let $\tilde{S}_1, \dots, \tilde{S}_5 \in \widetilde{\Sigma}$ be the five spheres that bounds a submanifold $N$ of $\MM$ that contains $\tilde{S}_1$ and $\tilde{S}_2$.
% Without loss of generality, suppose that $\tilde{S}_6$ separates $\tilde{S}_1, \tilde{S}_2, \tilde{S}_3$ from $\tilde{S}_4, \tilde{S}_5$ and that $\tilde{S}_7$ separates $\tilde{S}_1, \tilde{S}_2$ from $\tilde{S}_3, \tilde{S}_4, \tilde{S}_5$.
% Then we can get a statement for $\tilde{S}_6', \tilde{S}_6'', \tilde{S}_7', \tilde{S}_7''$.
% \begin{itemize}
%     \item $\tilde{S}_6'$ separates $\tilde{S}_1, \tilde{S}_2, \tilde{S}_4$ from $\tilde{S}_3, \tilde{S}_5$.
%     \item $\tilde{S}_6''$ separates $\tilde{S}_1, \tilde{S}_2, \tilde{S}_5$ from $\tilde{S}_3, \tilde{S}_4$.
%     \item $\tilde{S}_7'$ separates $\tilde{S}_1, \tilde{S}_3$ from $\tilde{S}_2, \tilde{S}_4, \tilde{S}_5$.
%     \item $\tilde{S}_7''$ separates $\tilde{S}_2, \tilde{S}_3$ from $\tilde{S}_1, \tilde{S}_4, \tilde{S}_5$.
% \end{itemize}
% Note $N$ should contain $\tilde{S}_0$.

% Meanwhile, $N \cap \rho^{-1}(\gamma)$ is a disjoint union of arcs connecting boundary components.
% Let $a_{ij}$ be the number of arcs joining $\tilde{S}_i$ and $\tilde{S}_j$ for all $i < j$.
% Then we have ten equations.
% \begin{align*}
%     i(\gamma, S_1) &= a_{12} + a_{13} + a_{14} + a_{15} \\
%     i(\gamma, S_2) &= a_{12} + a_{23} + a_{24} + a_{25} \\
%     i(\gamma, S_3) &= a_{13} + a_{23} + a_{34} + a_{35} \\
%     i(\gamma, S_4) &= a_{14} + a_{24} + a_{34} + a_{45} \\
%     i(\gamma, S_5) &= a_{15} + a_{25} + a_{35} + a_{45} \\
%     i(\gamma, S_6) &= a_{14} + a_{15} + a_{24} + a_{25} + a_{34} + a_{35} \\
%     i(\gamma, S_6') &= a_{13} + a_{15} + a_{23} + a_{25} + a_{34} + a_{45} \\
%     i(\gamma, S_6'') &= a_{13} + a_{14} + a_{23} + a_{24} + a_{35} + a_{45} \\
%     i(\gamma, S_7) &= a_{13} + a_{14} + a_{15} + a_{23} + a_{24} + a_{25} \\
%     i(\gamma, S_7') &= a_{12} + a_{14} + a_{15} + a_{23} + a_{34} + a_{35} \\
%     i(\gamma, S_7'') &= a_{12} + a_{13} + a_{24} + a_{25} + a_{34} + a_{35} \\
% \end{align*}
% This gives a matrix equation as follows.
% \[
%     \begin{pmatrix}
%         1 & 1 & 1 & 1 & 0 & 0 & 0 & 0 & 0 & 0 \\
%         1 & 0 & 0 & 0 & 1 & 1 & 1 & 0 & 0 & 0 \\
%         0 & 1 & 0 & 0 & 1 & 0 & 0 & 1 & 1 & 0 \\
%         0 & 0 & 1 & 0 & 0 & 1 & 0 & 1 & 0 & 1 \\
%         0 & 0 & 0 & 1 & 0 & 0 & 1 & 0 & 1 & 1 \\
%         0 & 0 & 1 & 1 & 0 & 1 & 1 & 1 & 1 & 0 \\
%         0 & 1 & 0 & 1 & 1 & 0 & 1 & 1 & 0 & 1 \\
%         0 & 1 & 1 & 0 & 1 & 1 & 0 & 0 & 1 & 1 \\
%         0 & 1 & 1 & 1 & 1 & 1 & 1 & 0 & 0 & 0 \\
%         1 & 0 & 1 & 1 & 1 & 0 & 0 & 1 & 1 & 0 \\
%         1 & 1 & 0 & 0 & 0 & 1 & 1 & 1 & 1 & 0
%     \end{pmatrix}
%     \begin{pmatrix}
%         a_{12} \\ a_{13} \\ a_{14} \\ a_{15} \\ a_{23} \\ a_{24} \\ a_{25} \\ a_{34} \\ a_{35} \\ a_{45}
%     \end{pmatrix}
%     =
%     \begin{pmatrix}
%         i(\gamma, S_1) \\ i(\gamma, S_2) \\ i(\gamma, S_3) \\ i(\gamma, S_4) \\ i(\gamma, S_5) \\ i(\gamma, S_6) \\ i(\gamma, S_6') \\ i(\gamma, S_6'') \\ i(\gamma, S_7) \\ i(\gamma, S_7') \\ i(\gamma, S_7'')
%     \end{pmatrix}
% \]

% \[
%     \begin{pmatrix}
%         1 & 1 & 1 & 1 & 0 & 0 & 0 & 0 & 0 & 0 \\
%         1 & 0 & 0 & 0 & 1 & 1 & 1 & 0 & 0 & 0 \\
%         0 & 1 & 0 & 0 & 1 & 0 & 0 & 1 & 1 & 0 \\
%         0 & 0 & 1 & 0 & 0 & 1 & 0 & 1 & 0 & 1 \\
%         0 & 0 & 0 & 1 & 0 & 0 & 1 & 0 & 1 & 1 \\
%         0 & 1 & 1 & 1 & 1 & 1 & 1 & 0 & 0 & 0 \\
%         1 & 1 & 0 & 0 & 0 & 1 & 1 & 1 & 1 & 0 \\
%         0 & 1 & 1 & 0 & 1 & 1 & 0 & 0 & 1 & 1 \\
%         0 & 0 & 1 & 1 & 0 & 1 & 1 & 1 & 1 & 0 \\
%         1 & 1 & 1 & 0 & 0 & 0 & 1 & 0 & 1 & 1
%     \end{pmatrix}
% \]

% Let us prove Theorem \ref{thm:finite_dimension}.
% Let $\Sigma = \{S_1, \dots, S_{3g-3}\}$ be a maximal sphere system that does not contain a separating sphere bounding a punctured doubled torus.
% For each $i = 1, \dots, 3g-3$, one can take three pairwise distinct essential spheres $S_i$, $S_i'$ and $S_i''$ such that \[i(S_i, S_j) = i(S_i', S_j) = i(S_i'', S_j) = 0\] for all $j \neq i$.
% Let $\mathcal{F}$ be the collection of $S_i, S_i', S_i''$'s for all $i$.

% \begin{lemma}
%     The map \[\mathcal{C} \to P(\mathbb{R}_+^{\mathcal{F}}), \quad \gamma \mapsto [i_*(\gamma)|_{\mathcal{F}}]\] is well-defined and injective.
% \end{lemma}

% \begin{proof}
%     First, since $S_1, \dots, S_{3g-3}$ form a maximal sphere system, for each $\gamma \in \mathcal{C}$, there exists $i = 1, \dots, 3g-3$ such that $i(\gamma, S_i)$ is nonzero.
%     If $r_\mathcal{F}: \mathbb{R}_+^{\mathcal{S}} \to \mathbb{R}_+^{\mathcal{F}}$ is the projection induced by the restriction to $\mathcal{F}$ and $p_\mathcal{F}: \mathbb{R}_+^{\mathcal{F}} \setminus \{0\} \to P(\mathbb{R}_+^{\mathcal{F}})$ is the natural projection, then the map $p_\mathcal{F} \circ r_\mathcal{F} \circ i_*: \mathcal{C} \to P(\mathbb{R}_+^\mathcal{F})$ defined by $\gamma \mapsto [i_*(\gamma) |_{\mathcal{F}}]$ is well-defined.

%     Let $S_1'''$ be an essential sphere such that $\{S_1''', S_2', S_3, \dots, S_{3g-3}\}$ is a maximal sphere system.
%     We claim that $i(\gamma, S_1''')$ is a linear combination of $i(\gamma, S_1), \dots, i(\gamma, S_{9g-9})$.
%     If $S_1''' \in \{S_1, S_1', S_1''\}$, then $i(\gamma, S_1''')$ is already determined.

%     Suppose $S_1''' \not\in \{S_1, S_1', S_1''\}$.
%     Consider the universal cover $\widetilde{\MM}$ and a lift $\tilde{S}_1'''$.
%     If $\tilde{S}_1$
% \end{proof}

% \begin{lemma}
%     For each $[f] \in \overline{(p \circ i_*)(\mathcal{C})}$, there exists $S \in \mathcal{F}$ such that $f(S) > 0$.
% \end{lemma}

% \begin{proof}
%     Suppose $f(S) = 0$ for all $S \in \mathcal{F}$.
%     Notice that there are sequences $\{ c_j \}$ and $\{ \gamma_j \}$ such that $c_j \, i_*(\gamma_j) \to f$ as $j \to \infty$.
%     Let $\widetilde{\MM}$ be the universal cover of $\MM$.
% \end{proof}

% \begin{proof}[Proof of Theorem \ref{thm:finite_dimension}]
%     If $\pi_\mathcal{F}: \mathbb{R}_+^{\mathcal{S}} \to \mathbb{R}_+^{\mathcal{F}}$ is the coordinate projection, then $\pi_\mathcal{F} (i_*(\mathcal{C}))$ does not contain the zero vector.
%     So, if $p_\mathcal{F}: \mathbb{R}_+^{\mathcal{F}} \setminus \{0\} \to P(\mathbb{R}_+^{\mathcal{F}})$ is the projection, then the map $p_\mathcal{F} \circ \pi_\mathcal{F}|_{i_*(\mathcal{C})}$ is well-defined.
%     This implies the map $p_\mathcal{F} \circ \pi_\mathcal{F} \circ p^{-1} |_{(p \circ i_*)(\mathcal{C})}: (p \circ i_*)(\mathcal{C}) \to P(\mathbb{R}_+^\mathcal{F})$ is well-defined.
% \end{proof}


\subsection{Fully irreducible elements}

A \emph{geodesic current} on $\FF$ is an $\FF$-invariant, flip-invariant, positive Radon measure on $\partial^2\FF$.
We write $\operatorname{Curr}(\FF)$ for the space of geodesic currents on $\FF$, and $\operatorname{\mathbb{P}Curr}(\FF)$ for its projectivization.
Let $cv(\FF)$ denote unprojectivized Outer space, and let $\overline{cv(\FF)}$ be its closure.
By Kapovich and Lustig \cite{MR2496058}, there is a unique $\Out(\FF)$-equivariant continuous \emph{intersection form}
\[
    \langle\,\cdot\,,\,\cdot\,\rangle \colon \overline{cv(\FF)} \times \operatorname{Curr}(\FF) \longrightarrow \mathbb{R}_+
\]
satisfying $\langle T, \eta_g \rangle = \| g \|_T$ for all $T \in \overline{cv(\FF)}$ and $g \in \FF$, where $\| g \|_T$ denotes the translation length of $g$ on $T$ and $\eta_g$ the counting current of $g$.

Each homotopy class $S$ of an essential sphere determines a one-edge free splitting whose dual tree $T_S$ lies in $\overline{cv(\FF)}$; in this way $\mathcal{S}$ is realized as a discrete subset of $\overline{cv(\FF)}$.
Moreover, every primitive simple closed curve is a counting current, which yields an embedding $\mathcal{C} \hookrightarrow \operatorname{Curr}(\FF)$.
Under these identifications, the geometric intersection number of a sphere and a primitive curve agrees with the value of the intersection form:
\[
    i(\gamma, S) = \langle T_S, \eta_\gamma \rangle.
\]

By Martin \cite{MR2693216}, the space $\operatorname{\mathbb{P}Curr}(\FF)$ contains a unique minimal nonempty closed $\Out(\FF)$-invariant subset $\mathcal{M}$, and the projectivized counting currents of primitive elements are dense in $\mathcal{M}$.
Consequently the coordinate map
\[
    j \colon \mathcal{M} \longrightarrow \mathbb{P}\big(\mathbb{R}_+^{\mathcal{S}}\big), \qquad [\mu] \longmapsto \big[\,(\langle T_S, \mu\rangle)_{S \in \mathcal{S}}\,\big]
\]
restricts on the dense set of primitive classes to $p \circ i_*|_{\mathcal{C}}$, and hence its image is the closure $\overline{(p \circ i_*)(\mathcal{C})}$.
We record the basic structural properties of $j$ before turning to dynamics.

\begin{proposition}\label{prop:surjection}
    The map $j \colon \mathcal{M} \to \overline{(p \circ i_*)(\mathcal{C})}$ is a continuous $\Out(\FF)$-equivariant surjection of compact Hausdorff spaces.
\end{proposition}

\begin{proof}
    Continuity follows from continuity of each pairing $\langle T_S, \cdot\rangle$ in the current variable.
    The space $\mathcal{M}$ is a closed subset of the compact space $\operatorname{\mathbb{P}Curr}(\FF)$, hence compact, and $\overline{(p \circ i_*)(\mathcal{C})} \subseteq \mathbb{P}(\mathbb{R}_+^{\mathcal{S}})$ is compact (Proposition~\ref{prop:C_compact}) and Hausdorff.
    Surjectivity holds because $j(\mathcal{M})$ is compact, hence closed, and contains the dense set $(p \circ i_*)(\mathcal{C})$.
    Finally, the family $\{T_S : S \in \mathcal{S}\}$ is $\Out(\FF)$-invariant, and equivariance of the intersection form gives $\langle T_S, \phi\mu\rangle = \langle \phi^{-1}T_S, \mu\rangle = \langle T_{\phi^{-1}S}, \mu\rangle$ for every $\phi \in \Out(\FF)$, which is the coordinate-permuting action of $\phi$ on $\mathbb{P}(\mathbb{R}_+^{\mathcal{S}})$; thus $j$ is equivariant.
\end{proof}

We now describe the dynamics of fully irreducible elements on the compact space $\overline{(p \circ i_*)(\mathcal{C})}$.
Throughout, $\phi \in \Out(\FF)$ is a fully irreducible outer automorphism, with expansion factors $\lambda_+ := \lambda(\phi) > 1$ and $\lambda_- := \lambda(\phi^{-1})^{-1} \in (0,1)$.
We invoke the following dynamical input.

\begin{theorem}[Martin \cite{MR2693216}, Uyanik \cite{MR3273532, MR3370027}]\label{thm:NS-currents}
    Let $\phi \in \Out(\FF)$ be a fully irreducible outer automorphism.
    Then $\phi$ acts on the minimal set $\mathcal{M}$ with north–south dynamics: there are distinct fixed points $[\mu_+], [\mu_-] \in \mathcal{M}$, with representatives normalized so that $\phi \cdot \mu_\pm = \lambda_\pm\, \mu_\pm$, such that for every compact $K \subseteq \mathcal{M} \setminus \{[\mu_-]\}$ and every neighborhood $U$ of $[\mu_+]$ in $\mathcal{M}$, one has $\phi^n(K) \subseteq U$ for all sufficiently large $n$, and symmetrically for $\phi^{-1}$ and $[\mu_-]$.
\end{theorem}

When $\phi$ is atoroidal, the north–south dynamics above holds on all of $\operatorname{\mathbb{P}Curr}(\FF)$ by Martin \cite{MR2693216} and Uyanik \cite{MR3273532}.
For a general fully irreducible $\phi$ the full projective space carries additional fixed currents, but Uyanik \cite{MR3370027} shows that north–south dynamics persists on the minimal set $\mathcal{M}$, which is all that the descent argument below requires.
In particular, for any primitive simple closed curve $\gamma$ with $[\eta_\gamma] \neq [\mu_-]$, the iterates $\phi^n[\eta_\gamma]$ converge to $[\mu_+]$ within $\mathcal{M}$.

The next lemma is a general principle: north–south dynamics descend through any equivariant continuous surjection of compact spaces, provided the two fixed points are not identified.
It is the only place where separation information about $j$ is needed.

\begin{lemma}\label{lem:descent}
    Let $q \colon X \to Y$ be a continuous, $\langle\phi\rangle$-equivariant surjection of compact Hausdorff spaces.
    Suppose $\phi$ acts on $X$ with north–south dynamics, with attracting and repelling fixed points $x_+$ and $x_-$.
    If $q(x_+) \neq q(x_-)$, then $\phi$ acts on $Y$ with north–south dynamics, with fixed points $y_\pm := q(x_\pm)$.
\end{lemma}

\begin{proof}
    Equivariance gives $\phi(y_\pm) = q(\phi(x_\pm)) = q(x_\pm) = y_\pm$, and $y_+ \neq y_-$ by hypothesis.
    Fix a compact set $K \subseteq Y \setminus \{y_-\}$ and an open neighborhood $V$ of $y_+$.
    Then $q^{-1}(K)$ is compact and avoids $x_-$ (otherwise $y_- = q(x_-) \in K$), so $q^{-1}(K) \subseteq X \setminus \{x_-\}$.
    The set $U := q^{-1}(V)$ is an open neighborhood of $x_+$, so by north–south dynamics on $X$ there is $N$ with $\phi^n\big(q^{-1}(K)\big) \subseteq U$ for all $n \geq N$.
    Using surjectivity and equivariance,
    \[
        \phi^n(K) = \phi^n\big(q(q^{-1}(K))\big) = q\big(\phi^n(q^{-1}(K))\big) \subseteq q(U) = q\big(q^{-1}(V)\big) \subseteq V
    \]
    for all $n \geq N$.
    The symmetric statement for $\phi^{-1}$ and $y_-$ follows by applying the argument to $\phi^{-1}$.
\end{proof}

It remains to separate the two fixed points under $j$.
We first show that degeneracy never occurs on the minimal set; this both confirms that $j$ is well-defined on all of $\mathcal{M}$ and supplies the non-degeneracy needed for the separation argument.
Call a current $\mu$ \emph{non-degenerate} if $\langle T_S, \mu\rangle > 0$ for some $S \in \mathcal{S}$.

\begin{lemma}\label{lem:nondegeneracy}
    Every class $[\mu] \in \mathcal{M}$ is non-degenerate.
    In particular, the attracting current $\mu_+$ satisfies $\langle T_S, \mu_+\rangle > 0$ for some $S \in \mathcal{S}$.
\end{lemma}

\begin{proof}
    Let $\mathcal{Z} := \{[\mu] \in \operatorname{\mathbb{P}Curr}(\FF) : \langle T_S, \mu\rangle = 0 \text{ for all } S \in \mathcal{S}\}$.
    Each pairing $\langle T_S, \cdot\rangle$ is continuous and positively homogeneous in the current variable, so the condition $\langle T_S, \mu\rangle = 0$ is well-defined on projective classes and closed; hence $\mathcal{Z}$, an intersection of closed sets, is closed.
    As in the proof of Proposition~\ref{prop:surjection}, the family $\{T_S : S \in \mathcal{S}\}$ is $\Out(\FF)$-invariant and the intersection form is equivariant, so $\mathcal{Z}$ is $\Out(\FF)$-invariant.
    Now let $\gamma$ be the homotopy class of a primitive simple closed curve, and let $S$ be the homotopy class of the non-separating essential sphere dual to a representative, so that the curve crosses the sphere exactly once.
    Then $\langle T_S, \eta_\gamma\rangle = i(\gamma, S) = 1$.
    It follows that $[\eta_\gamma] \notin \mathcal{Z}$.
    Since $[\eta_\gamma] \in \mathcal{M}$, the closed $\Out(\FF)$-invariant set $\mathcal{Z} \cap \mathcal{M}$ is a proper subset of $\mathcal{M}$; by minimality of $\mathcal{M}$ it is therefore empty.
    Hence no class in $\mathcal{M}$ is degenerate, and in particular $\mu_+ \in \mathcal{M}$ is non-degenerate.
\end{proof}

\begin{lemma}\label{lem:reciprocity}
    The fixed points are separated by $j$, that is, $j([\mu_+]) \neq j([\mu_-])$.
\end{lemma}

\begin{proof}
    Write $f_\pm(S) := \langle T_S, \mu_\pm\rangle$.
    By equivariance of the intersection form and the relation $\phi\cdot\mu_\pm = \lambda_\pm\mu_\pm$,
    \[
        f_\pm(\phi^{-1}S) = \langle T_{\phi^{-1}S}, \mu_\pm\rangle = \langle T_S, \phi\,\mu_\pm\rangle = \lambda_\pm\, f_\pm(S)
    \]
    for all $S \in \mathcal{S}$.
    Suppose, for contradiction, that $j([\mu_+]) = j([\mu_-])$.
    Then there is $\kappa > 0$ with $f_+(S) = \kappa\, f_-(S)$ for all $S \in \mathcal{S}$.
    Evaluating this identity at $\phi^{-1}S$ and using the displayed scaling,
    \[
        \lambda_+\, f_+(S) = f_+(\phi^{-1}S) = \kappa\, f_-(\phi^{-1}S) = \kappa\,\lambda_-\, f_-(S) = \lambda_-\, f_+(S)
    \]
    for all $S \in \mathcal{S}$.
    Since $\lambda_+ > 1 > \lambda_-$, this forces $f_+(S) = 0$ for every $S$, contradicting the non-degeneracy of $\mu_+$ due to Lemma~\ref{lem:nondegeneracy}.
    Hence $j([\mu_+]) \neq j([\mu_-])$.
\end{proof}

Combining Theorem~\ref{thm:NS-currents} with Lemmas~\ref{lem:descent}, \ref{lem:nondegeneracy}, and~\ref{lem:reciprocity} yields the main dynamical statement of this subsection.

\begin{theorem}\label{thm:loxodromic}
    Let $\phi \in \Out(\FF)$ be a fully irreducible outer automorphism.
    Then $\phi$ acts on the compact space $\overline{(p \circ i_*)(\mathcal{C})}$ with north–south dynamics; that is, $\phi$ acts loxodromically, with attracting and repelling fixed points $j([\mu_+])$ and $j([\mu_-])$.
\end{theorem}

\begin{proof}
    By Theorem~\ref{thm:NS-currents}, $\phi$ acts on the compact Hausdorff space $\mathcal{M}$ with north–south dynamics, with fixed points $[\mu_\pm] \in \mathcal{M}$.
    By Proposition~\ref{prop:surjection}, the map $j \colon \mathcal{M} \to \overline{(p \circ i_*)(\mathcal{C})}$ is a continuous $\langle\phi\rangle$-equivariant surjection between compact spaces, and $j([\mu_+]) \neq j([\mu_-])$ by Lemma~\ref{lem:reciprocity}.
    The conclusion now follows from Lemma~\ref{lem:descent}.
\end{proof}

% \subsection*{Geodesic currents and iwip elements}

% A \emph{geodesic current} on $\FF$ is an $\FF$-invariant and flip-invariant positive Radon measure on $\partial^2\FF$.
% Let $\operatorname{Curr}(\FF)$ denote the space of geodesic currents on $\FF$.
% If $cv(\FF)$ is the unprojectived outer space, then there exists a unique $\Out(\FF)$-equivariant continuous \emph{intersection form}
% \[
%     \langle \cdot, \cdot \rangle : \overline{cv(\FF)} \times \operatorname{Curr}(\FF) \longrightarrow \mathbb{R}_+
% \]
% satisfying that $\langle T, \eta_g \rangle = \| g \|_T$ for all $T \in \overline{cv(\FF)}$ and $g \in \FF$.
% See Kapovich and Lustig \cite{MR2496058}.

% Notice that $\mathcal{S}$ is considered as a discrete subset of $cv(\FF)$, in addition, every primitive simple closed curve can be seen as a counting current so that there exists an embedding $\mathcal{C} \to \operatorname{Curr}(\FF)$.
% So the intersection number between a sphere and a primitive simple closed curve is the value of the intersection form between them, that is, 
% \[
%     i(\gamma, S) = \min\{\langle T_s, \eta_c \rangle \mid s \in S, \, c \in \gamma\}.
% \]

% By Martin \cite{MR2693216}, there exists a unique $\Out(\FF)$-invariant minimal set $\mathcal{M}$ of $\operatorname{\mathbb{P}Curr}(\FF)$.
% The set of projective counting measures on primitive elements is dense in $\mathcal{M}$.
% That is, there exists a surjection $\mathcal{M} \to \overline{(p \circ i_*)(\mathcal{C})}$.

% \begin{proposition}
%     The surjection $\mathcal{M} \to \overline{(p \circ i_*)(\mathcal{C})}$ is a homeomorphism.
% \end{proposition}












\bibliographystyle{amsalpha}
\bibliography{references}

\end{document}
